\documentclass[11pt,a4paper]{article}
\usepackage[utf8]{inputenc}
\usepackage[T1]{fontenc}
\usepackage[french]{babel}
\usepackage{amsmath,amssymb}
\usepackage{geometry}
\usepackage{hyperref}
\usepackage{xcolor}
\usepackage{booktabs} 
\usepackage[most]{tcolorbox}
\usepackage{colortbl}

\geometry{margin=2.5cm}

\title{\Large \textbf{Rapport d'Avancement de Stage de Recherche}\\ \huge \textbf{Auto-Amélioration des Agents LLM en Environnement Multi-tour}\\ \large Criteo AI Lab}
\author{\textbf{Vadim Lagresle} \\ Encadrants : Alberto Lumbreras, Patrick Gallinari, Alain Rakotomamonjy, Sylvain Lamprier}
\date{\today}

\begin{document}

\maketitle

\vspace{2cm}



\begin{abstract}
\noindent Ce rapport présente une étude expérimentale, en partie exploratoire, du post-entraînement d'agents fondés sur des modèles de langage par apprentissage par renforcement, en environnement interactif multi-tour. Le passage à l'échelle des environnements d'entraînement, par leur diversité, leur difficulté et leurs vérificateurs, est devenu le principal moteur des capacités agentiques. Notre question est celle qui se pose une fois l'environnement donné : comment y stabiliser l'apprentissage, et comment le rendre efficient. Nous l'étudions sur TextCraft, un environnement de l'écosystème AgentGym-RL à récompense vérifiable, avec pour horizon l'adaptation de ces méthodes aux agents d'achat de Criteo.

Nos deux travaux principaux répondent à ces deux questions. L'efficience vient de LoRA, une adaptation de rang faible qui n'entraîne qu'une petite fraction des poids du modèle et rend l'expérimentation possible sur un seul GPU ; nous montrons à quelles conditions elle reste stable en multi-tour. La stabilité vient de curriculums, c'est-à-dire de régimes d'entraînement qui augmentent progressivement la difficulté des épisodes présentés à l'agent, par le nombre de tours, le budget de génération ou la complexité des tâches. À recette identique et à compute contrôlé, l'entraînement sans curriculum s'effondre là où les trois curriculums tiennent. Le meilleur porte un modèle de trois milliards de paramètres à 82\,\% de réussite, au-dessus du résultat publié par le système de référence sur une infrastructure multi-GPU, et l'ensemble de l'entraînement est documenté pour en permettre la reproduction.

Ce parcours fournit ainsi un retour d'expérience détaillé sur les obstacles rencontrés, du masquage des observations à la conception de la récompense, qui pourra faire gagner du temps au développement des agents d'achat de Criteo. Enfin, nous analysons ces résultats au-delà de la stabilisation : ce que le curriculum change à l'ordre d'acquisition des difficultés et à l'élimination des erreurs, et ce qu'il dit de l'hypothèse selon laquelle le renforcement ne ferait que concentrer les capacités du modèle de base. Nos mesures vont dans le sens contraire, et ouvrent sur des questions plus larges de méta-apprentissage et d'apprentissage continu.
\end{abstract}



\newpage




\section*{Remerciements}
\addcontentsline{toc}{section}{Remerciements}

Je tiens d'abord à remercier Patrick Gallinari, Alberto Lumbreras et Alain Rakotomamonjy, qui m'ont suivi chaque semaine depuis le début de ce stage et dont les conseils ont beaucoup compté. Patrick m'a appris à lire les évaluations par les pairs des travaux sur lesquels je m'appuyais ; Alberto m'a fait profiter de sa connaissance des harnais agentiques et de ses lectures sur LoRA ; Alain m'a guidé dans les miennes. Merci également à Maxime, mon manager, et aux équipes Agentic Commerce et Agentic Core Team de Criteo, pour m'avoir fait découvrir leurs missions, leurs savoir-faire et le contexte agentique de l'entreprise. Ma reconnaissance va enfin à Sylvain Lamprier pour ses conseils de lecture sur les agents autotéliques et l'organisation hiérarchique des compétences, à Flavian et Imad pour la découverte de DeepShopper, un projet interne à Criteo, et à Otmane Sakhi pour nos discussions autour de l'échantillonnage préférentiel auto-normalisé.



\text




\vspace{0.5cm}

\newpage
\setcounter{tocdepth}{2} % Limite le sommaire aux sections et sous-sections (exclut les subsubsections)
\tableofcontents

\newpage
\section*{Notations et Abréviations}
\addcontentsline{toc}{section}{Notations et Abréviations}

\subsection*{Notations}
\begin{tabular}{ll}
\toprule
Symbole & Signification \\
\midrule
$\pi_\theta$ & politique entraînée, paramétrée par les poids $\theta$ \\
$\pi_{\text{ref}}$ & politique de référence de la pénalité KL (ancre) \\
$\tau$ & trajectoire (épisode complet) \\
$s_t,\ a_t,\ o_t$ & état, action et observation au pas $t$ \\
$T$ & horizon de l'épisode (nombre maximal de tours) \\
$R(\tau)$, $R_i$ & récompense terminale d'une trajectoire ; celle du rollout $i$ \\
$\hat{A}$, $\hat{A}_i$ & estimateur d'avantage ; avantage relatif du rollout $i$ dans son groupe \\
$V_\phi$ & fonction de valeur estimée par le réseau critique (PPO) \\
$r_t(\theta)$ & ratio d'importance $\pi_\theta(a_t \mid s_t)/\pi_{\theta_{\text{old}}}(a_t \mid s_t)$ \\
$G$ & taille de groupe GRPO (trajectoires par tâche dans un même groupe) \\
$k$ & nombre de tirages indépendants par tâche (budget du pass@$k$) \\
$\beta$ & coefficient de la pénalité $D_{\text{KL}}(\pi_\theta \parallel \pi_{\text{ref}})$ \\
$\epsilon$ & rayon de \emph{clipping} de l'objectif PPO/GRPO \\
$r,\ \alpha$ & rang et facteur d'échelle d'un adaptateur LoRA \\
$d$ & profondeur (\emph{depth}) d'une tâche TextCraft, $d \in \{1,\dots,4\}$ \\


\bottomrule
\end{tabular}

\medskip
\noindent\textbf{pass@$k$.} Fraction des tâches de l'ensemble d'évaluation résolues par
au moins un tirage parmi $k$ tirages indépendants de la politique. Le pass@$1$
mesure la fiabilité d'un tirage unique ; le pass@$k$ pour $k$ grand approche le
support de la politique. Toute valeur de pass@$k$ est rapportée avec son budget
$k$ et sa température d'échantillonnage.

\subsection*{Abréviations}
\begin{tabular}{ll}
\toprule
Abréviation & Signification \\
\midrule
LLM & grand modèle de langage (\emph{Large Language Model}) \\
SFT & \emph{Supervised Fine-Tuning} \\
RL & apprentissage par renforcement \\
RLHF & RL à partir de préférences humaines \\
RLVR & RL à récompense vérifiable \\
NLL & log-vraisemblance négative \\
TRPO & \emph{Trust Region Policy Optimization} \\
PPO & \emph{Proximal Policy Optimization} \\
GRPO & \emph{Group Relative Policy Optimization} \\
DPO & \emph{Direct Preference Optimization} \\
KL & divergence de Kullback-Leibler \\
(PO)MDP & processus de décision markovien (partiellement observable) \\
LoRA & \emph{Low-Rank Adaptation} \\
CoT & \emph{Chain-of-Thought} \\
PRM & \emph{Process Reward Model} \\
MCP & \emph{Model Context Protocol} \\
SNIS & échantillonnage préférentiel auto-normalisé \\
\bottomrule
\end{tabular}



\newpage




\section{Introduction}

\subsection{Évolution du signal d'apprentissage}

Les agents auto-améliorants (\emph{self-improving agents}) comptent aujourd'hui parmi les sujets les plus actifs de l'apprentissage automatique. On le voit aux cours universitaires qui leur sont consacrés, comme le récent CS329A de Stanford~\cite{cs329a}, au nombre de publications dans les grandes conférences du domaine, et aux investissements des constructeurs de modèles. Pour comprendre pourquoi ce sujet attire autant, il est utile de retracer l'évolution du signal d'apprentissage des LLM au cours de la dernière décennie. Au-delà des opportunités qu'ouvrent les applications agentiques, cette évolution prolonge les mutations précédentes du domaine : à chaque fois, c'est la nature de la donnée, et du signal qu'on en extrait, qui a changé.

La rupture initiale introduite par l'architecture Transformer \cite{vaswani2017attention} a permis de passer à l'échelle sur la tâche d'autocomplétion autorégressive, en prédisant le token suivant sur d'immenses corpus textuels issus du web. Ce pré-entraînement massif confère au modèle une vaste représentation linguistique du monde. Toutefois, ce signal d'apprentissage reste implicite et passif : le modèle apprend à imiter la distribution du texte web sans garantie d'alignement avec les exigences de précision, de logique ou de sécurité attendues par un utilisateur humain.

Pour surmonter ces contraintes, la phase de post-training s'est imposée comme l'étape centrale d'alignement comportemental. Dans un premier temps, le \emph{Supervised Fine-Tuning} a structuré la forme des réponses en instructions~\cite{wei2022flan}. Par la suite, l'apprentissage par renforcement à partir de préférences humaines (RLHF) \cite{ouyang2022training} a introduit un signal d'évaluation global, permettant de pénaliser les réponses toxiques ou hors sujet. Plus récemment, les efforts se sont concentrés sur l'émergence de capacités de raisonnement complexe, où le modèle est encouragé à produire une trace de réflexion (\emph{Chain-of-Thought}) avant sa réponse \cite{wei2022cot}. Dans ce cadre, la nature du signal a de nouveau évolué vers des récompenses vérifiables de manière déterministe ($0$ ou $1$), notamment en mathématiques ou en programmation, qui permettent d'évaluer la justesse d'une séquence de pensée sans dépendre d'annotations humaines subjectives. Cet apprentissage par renforcement à récompense vérifiable (RLVR) a permis de faire passer à l'échelle les LLM sur des tâches de raisonnement complexes, et il est au cœur des performances des modèles actuels \cite{openai_o1_2024, deepseekr1_2025}. Son intérêt pratique est double : il s'affranchit de l'entraînement et de l'inférence d'un modèle de récompense, et il renoue avec les questions classiques de l'apprentissage par renforcement, formulées en objectifs, en tâches et en environnements, que la vague des LLM avait temporairement éclipsées. De nombreuses applications restent hors de portée du RLVR, comme l'évaluation de résumés ou la recommandation ; nous ne les traitons pas ici, et visons à maîtriser les techniques qui passent à l'échelle avant d'aborder leurs applications spécifiques.

Dans le cadre du RLVR, les trajectoires d'entraînement se heurtent aujourd'hui à une limite matérielle et économique critique : la rareté et le coût élevé des données humaines de haute qualité dans les domaines exigeants. La donnée annotée constitue désormais le goulot d'étranglement majeur de la discipline, ce qui impose un changement de paradigme vers la génération de données synthétiques et l'apprentissage autonome. Des approches d'auto-amélioration telles que STaR \cite{zelikman2022star} ou ReST \cite{singh2024beyond} ont démontré qu'un modèle pouvait générer ses propres traces de raisonnement, filtrer les trajectoires menant à la bonne réponse, puis se réentraîner sur ses propres succès. Ce processus itératif accroît l'autonomie du modèle tout en réduisant considérablement la dépendance envers la supervision humaine directe.

Une thèse théorique influente interprète cette dynamique avec scepticisme : l'auto-amélioration sur données synthétiques ne créerait pas d'information fondamentalement nouvelle, et s'analyserait comme un mécanisme de \emph{sharpening} \cite{huang2024self}, où le réentraînement concentre la masse de probabilité de la politique sur les trajectoires à haute récompense, amortissant le coût du calcul à l'inférence sans étendre les capacités. Ce rapport traite cette thèse comme une hypothèse de travail à mettre à l'épreuve, non comme un acquis : une part importante de notre étude consiste précisément à mesurer si une structuration adéquate de l'apprentissage permet de dépasser les frontières que cette lecture prédit. L'émergence de l'IA agentique s'inscrit dans la continuité directe de cette évolution : il ne s'agit plus seulement de générer une trace de raisonnement en un seul jet, mais de déployer ces capacités au sein d'environnements numériques interactifs, où l'agent apprend de manière autonome à partir des retours de ses actions.

Le cadre agentique reformule enfin le goulot d'étranglement de la donnée sous une forme nouvelle : ce qui doit désormais passer à l'échelle n'est plus un corpus statique de textes ou d'annotations, mais l'environnement d'interaction lui-même, c'est-à-dire la définition conjointe d'un état initial, d'une surface d'outils, d'une tâche et d'un vérificateur. Les rapports techniques récents des constructeurs de modèles convergent sur ce point. L'équipe Qwen identifie explicitement le passage à l'échelle des environnements (\emph{environment scaling}) comme le verrou central du RL agentique \cite{qwen3coder2025}, et attribue l'essentiel des gains de ses dernières générations à l'extension massive de la difficulté et de la diversité de ses tâches d'entraînement \cite{qwen35_2026}. Kimi K2 fait reposer ses capacités agentiques sur un pipeline de synthèse à grande échelle d'outils, de tâches et de trajectoires vérifiées, couplé à un entraînement RL conjoint sur environnements réels et simulés \cite{kimik2_2025}. Symétriquement, les benchmarks de référence du domaine ($\tau^2$-bench \cite{tau2bench2025}, Toolathlon \cite{toolathlon2025}, MCP-Atlas \cite{mcpatlas2026}) sont moins des jeux de données que des travaux de définition du problème : y construire une tâche, c'est spécifier un monde, ses règles et son critère de succès. La définition de l'environnement devient ainsi un objet scientifique à part entière, au même titre que la curation des corpus l'a été pour le pré-entraînement. Ce constat structure l'ensemble de ce rapport.

\subsection{Rappel et similitudes des algorithmes d'entraînement}

Nous partons du principe que les techniques des Transformers et du traitement du langage sont connues du lecteur, et renvoyons aux cours de Stanford sur la modélisation du langage (\href{https://cs336.stanford.edu/}{CS336}) et sur les grands modèles de langage (\href{https://stanford-cs324.github.io/winter2022/}{CS324}) pour une présentation complète. De même, cette sous-section ne présente que les éléments d'apprentissage par renforcement nécessaires à la lecture du rapport, en particulier GRPO, notre algorithme de travail. L'annexe~A développe le formalisme complet, du cadre classique du RL (POMDP) aux méthodes acteur-critique ; pour un traitement approfondi du RL profond, nous renvoyons aux cours de Stanford (\href{https://cs224r.stanford.edu/}{CS224R}) et de Berkeley (\href{https://rail.eecs.berkeley.edu/deeprlcourse/}{CS285}).

L'évolution de la nature des données et du signal d'apprentissage s'est accompagnée d'une diversité des algorithmes d'optimisation utilisés à chaque phase. Le pré-entraînement repose sur l'architecture Transformer \cite{vaswani2017attention}, qui a supplanté les architectures récurrentes et séquentielles. Ces dernières souffraient d'un goulot d'étranglement computationnel empêchant la parallélisation et d'un phénomène de disparition du gradient sur les longs contextes. Le mécanisme d'auto-attention permet de calculer en parallèle les dépendances entre tous les tokens d'une séquence, indépendamment de leur distance relative. Cette capacité de parallélisation massive sur GPU a rendu possible le traitement de corpus web bruts à très grande échelle.

Dans la même logique, la phase de \emph{Supervised Fine-Tuning} conserve la même architecture autorégressive et le même principe de prédiction de token que le pré-entraînement, mais s'en distingue par trois ajustements techniques, au-delà du volume restreint et de la haute qualité des données. Tout d'abord, l'injection de tokens de contrôle structurants (\emph{chat templates}) transforme le texte continu en une machine à états alternant les rôles d'utilisateur et d'assistant, ce qui conditionne la génération vers un comportement d'agent conversationnel. Ensuite, le régime d'optimisation est adapté : un taux d'apprentissage nettement plus faible sur un nombre réduit d'époques concentre la masse de probabilité sur les structures de réponse attendues, sans détruire les représentations acquises. Enfin, alors que le pré-entraînement calcule le gradient sur l'intégralité du texte, le SFT applique un masquage de la perte (\emph{loss masking}) sur les tokens de l'instruction $x = (x_1, \dots, x_N)$, afin de ne rétropropager le signal que sur les tokens de la réponse cible $y = (y_1, \dots, y_M)$. Le modèle minimise la log-vraisemblance négative moyenne calculée exclusivement sur les $M$ tokens de la réponse :
\begin{equation}
\mathcal{L}_{\text{SFT}}(\theta) = -\frac{1}{M} \sum_{m=1}^{M} \log \pi_\theta(y_m \mid x, y_{<m})
\end{equation}

Contrairement au SFT, dont la mécanique de mise à jour des poids reste proche de celle du pré-entraînement, l'alignement par renforcement introduit un signal et un cadre d'optimisation fondamentalement différents. Il convient d'abord d'établir la différence fondamentale entre les deux. En fine-tuning, le modèle apprend à reproduire, token par token, une séquence cible : le signal d'apprentissage est le token suivant sachant la séquence qui précède. En RL, le signal est la récompense (\emph{reward}), un scalaire quantifiant la qualité globale d'une réponse produite par l'agent. L'objectif devient de favoriser les trajectoires menant aux récompenses les plus élevées : dans le cas d'une récompense \emph{sparse} binaire, typique des benchmarks mathématiques par exemple, on renforce les trajectoires ayant atteint la bonne réponse et l'on dévalorise les autres. On parle d'apprentissage \emph{par renforcement} parce que l'algorithme renforce la distribution de sortie du modèle vers une distribution cible qu'il contenait déjà, en germe, à l'issue du pré-entraînement.

Pour exploiter un signal de récompense scalaire $R(\tau)$ défini sur une trajectoire entière $\tau$, il faut exprimer un gradient, mais la récompense est rarement dérivable contrairement aux fonctions de pertes utilisées en fine-tuning supervisé. Les méthodes de gradient de politique s'appuient sur l'astuce de REINFORCE (dérivée du logarithme), qui réécrit le gradient de l'espérance de la récompense sous une forme calculable par échantillonnage Monte-Carlo :
\begin{equation}
\nabla_\theta \mathbb{E}_{\tau \sim \pi_\theta}[R(\tau)] = \mathbb{E}_{\tau \sim \pi_\theta} \left[ \sum_{t=1}^T \nabla_\theta \log \pi_\theta(a_t \mid s_t) R(\tau) \right]
\end{equation}
Sans cette réécriture, le gradient ne pourrait pas être estimé par échantillonnage.

\textbf{Remarque.} Le SFT et le RL ne sont pas fondamentalement éloignés dans leur mécanique de mise à jour des poids, bien que les théories sous-jacentes diffèrent. Sous une récompense \emph{sparse} et une stratégie strictement \emph{on-policy}, c'est-à-dire lorsque l'entraînement s'effectue sur des générations produites par la politique courante (annexe~A), l'algorithme de gradient de politique se réduit exactement à un SFT appliqué aux seules trajectoires ayant atteint la bonne réponse. Cette équivalence illustre à la fois la richesse et l'instabilité du RL : il identifie sans annotation humaine les trajectoires à privilégier, mais le seul signal disponible est la réussite finale, alors que des stratégies erronées peuvent occasionnellement mener à la bonne réponse. Le faible coût d'annotation se paie par une difficulté accrue de stabilisation, ces algorithmes pouvant s'effondrer (\emph{collapse}) en quelques itérations ; les sections~2.3 et~4.2 y reviennent en détail.

Utilisée seule, la récompense $R(\tau)$ constitue un signal brut : elle ne dit pas si une trajectoire est bonne en soi, seulement quel gain elle a rapporté, sans point de comparaison. Le cadre acteur-critique introduit un réseau critique $V_\phi(s)$, entraîné à estimer la fonction de valeur, c'est-à-dire l'espérance du retour obtenu depuis l'état $s_t$ en suivant la politique courante. Cette quantité ne dépend que de l'état, et non de l'action jouée : elle fournit une référence permettant de recentrer le signal sous la forme d'un avantage, dont l'estimateur le plus simple s'écrit $\hat{A}_t = R_t - V_\phi(s_t)$ (l'annexe~A détaille les estimateurs employés en pratique). Substituer $\hat{A}_t$ à $R(\tau)$ revient à ne renforcer une action que dans la mesure où elle a fait mieux que la moyenne attendue, ce qui réduit considérablement la variance. À cette structure s'ajoute, pour les LLM, une pénalité de divergence de Kullback-Leibler ancrant la politique entraînée $\pi_\theta$ à une politique de référence figée $\pi_{\text{ref}}$, généralement le modèle issu du SFT, afin d'éviter la dérive catastrophique et l'oubli des capacités de base :
\begin{equation}
\mathcal{L}_{\text{RL}}(\theta) = -\mathbb{E}_{(s_t, a_t)} \left[ \log \pi_\theta(a_t \mid s_t) \hat{A}_t \right] + \beta \, D_{\text{KL}}(\pi_\theta \parallel \pi_{\text{ref}})
\end{equation}

C'est dans ce cadre acteur-critique que s'inscrit \emph{Proximal Policy Optimization} (PPO) \cite{schulman2017proximal}, conçu pour le contrôle et la robotique bien avant les LLM modernes, puis repris à partir de 2022 comme l'algorithme central du RLHF \cite{ouyang2022training}. Réutiliser un même lot de trajectoires pour plusieurs pas de gradient améliore considérablement l'efficacité en données, mais impose de corriger l'écart entre la politique courante et celle qui a généré les trajectoires. Un ratio d'importance $r_t(\theta) = \pi_\theta(a_t \mid s_t) / \pi_{\theta_{\text{old}}}(a_t \mid s_t)$ mesure cet écart, et un mécanisme de \emph{clipping} le borne, retirant à la politique tout intérêt à s'éloigner fortement de $\pi_{\theta_{\text{old}}}$ :
\begin{equation}
\mathcal{L}_{\text{PPO}}(\theta) = -\mathbb{E}_{(s_t, a_t)} \left[ \min\left( r_t(\theta)\, \hat{A}_t, \ \text{clip}\big(r_t(\theta),\, 1-\epsilon,\, 1+\epsilon\big)\, \hat{A}_t \right) \right]
\end{equation}

Notre algorithme de travail, \emph{Group Relative Policy Optimization} (GRPO) \cite{shao2024deepseekmath}, en est l'évolution la plus utilisée en RLVR. Il supprime le réseau critique, très coûteux en mémoire GPU, en échantillonnant un groupe de $G$ trajectoires $\{\tau_1, \dots, \tau_G\}$ pour une même tâche, et en remplaçant l'estimateur d'avantage par un avantage relatif normalisé au sein du groupe, l'objectif clippé de PPO restant conservé :
\begin{equation}
\hat{A}_i = \frac{R_i - \text{mean}(\mathbf{R})}{\text{std}(\mathbf{R}) + \delta}
\end{equation}
où $R_i$ est la récompense de la trajectoire $i$ et $\mathbf{R}$ le vecteur des récompenses du groupe. Cette construction a une conséquence directe qui traverse tout le rapport : si toutes les trajectoires d'un groupe échouent, ou si toutes réussissent, l'écart-type est nul, l'avantage aussi, et la tâche ne produit aucun gradient. La quantité de signal d'apprentissage dépend donc de l'adéquation entre la difficulté des tâches et le niveau courant de la politique.

C'est dans ce cadre que s'inscrit le présent travail. Parmi les leviers d'amélioration des agents, ingénierie du harnais, calcul à l'inférence et optimisation des poids, nous avons choisi le dernier : l'amélioration de la politique elle-même par renforcement, en environnement multi-tour, là où l'espace des trajectoires croît exponentiellement avec l'horizon $T$ et où la récompense reste terminale. Le passage à l'échelle des environnements relève des constructeurs de modèles ; notre question est celle qui se pose une fois l'environnement donné. Nous l'étudions sur un environnement contrôlé à récompense vérifiable, TextCraft, issu de la plateforme AgentGym-RL~\cite{xi_agentgym-rl_2025}, qui permet d'isoler les variables de l'apprentissage, avec pour horizon les agents d'achat de Criteo.

Nos deux travaux principaux portent sur l'efficience et la stabilité de cet apprentissage. Le premier part de la réplication du système de référence sur un seul GPU, puis remplace son \emph{full finetuning} par une adaptation de rang faible, LoRA, bien moins coûteuse. La recette publiée pour LoRA s'effondre dans notre régime ; nous en analysons la cause et proposons un correctif, une ancre KL mobile. Ce parcours est rapporté comme un retour d'expérience, organisé par constats de stabilisation. Le second travail porte sur les curriculums. Un curriculum, au sens de Bengio et al.~\cite{bengio2009curriculum} puis de la littérature du renforcement~\cite{narvekar2020curriculum}, est un régime d'entraînement qui augmente progressivement la difficulté de ce qui est présenté à l'apprenant, soit en ordonnant les tâches, soit en dosant les moyens qu'on lui accorde pour les résoudre. Nous en comparons trois, sur l'horizon d'interaction, le budget de génération et la profondeur des tâches, à recette identique et à compute contrôlé, face à des références sans curriculum. Un troisième volet analyse ces résultats : ce que le curriculum change à l'ordre d'acquisition des difficultés et à l'élimination des erreurs, et ce que nos mesures disent de l'hypothèse du plafond pass@$k$. Les pistes explorées sans aboutir, dont la recomposition de trajectoires par SNIS et le transfert entre planification et contrôle interactif, sont résumées dans le corps du rapport et développées en annexe.

La suite de ce rapport s'organise comme suit. La section~2 dresse l'état de l'art en six temps : les systèmes agentiques et leurs harnais ; les environnements et bancs d'essai, comme outils de définition des objectifs de l'IA agentique ; la stabilité et la mesure du RL en environnement interactif ; le curriculum, l'autocurriculum et les agents autotéliques ; le transfert entre planification mono-tour et contrôle multi-tour ; et les autres pistes considérées. La section~3 présente le cadre technique, la plateforme AgentGym-RL, l'environnement TextCraft et notre positionnement. La section~4 expose le protocole et les résultats : stabilisation, curriculums, comparaison à compute contrôlé, puis expériences complémentaires. La section~5 délimite le périmètre et ouvre les perspectives, avant la conclusion. Cinq annexes complètent le document : les fondements du RL profond (A), la dérivation du SNIS (B), une analyse critique du parcours (C), le transfert planification-contrôle (D) et l'autocurriculum MAGELLAN (E).




%-----------------------------------------------------------------------------

\newpage
\section{État de l'art}

\subsection{Définition des Systèmes Agentiques dans la Continuité des LLMs}

L'émergence de l'IA agentique ne se réduit pas à l'ajout d'interfaces d'appel d'outils aux modèles existants. Elle découle en grande partie du goulot d'étranglement des données décrit en introduction : les corpus statiques s'épuisent et les tâches d'évaluation mono-tour ne suffisent plus à discriminer les modèles. Le déploiement des modèles au sein d'environnements numériques interactifs constitue alors la suite logique des travaux sur l'alignement et l'auto-amélioration. Les rapports techniques des constructeurs documentent ce basculement : l'équipe Qwen attribue les gains agentiques de ses dernières générations au passage à l'échelle de la génération d'environnements et de l'apprentissage par renforcement sur ces environnements \cite{qwen3coder2025, qwen35_2026}, et Kimi K2 comme Magistral décrivent des pipelines analogues de synthèse de tâches et d'entraînement RL \cite{kimik2_2025, magistral2025}. La section 2.2 détaille ce constat.

Nous appelons agent un modèle de langage placé dans une boucle d'interaction avec un environnement. À chaque pas, il reçoit une observation, sélectionne une action au sein d'un espace d'outils exposé par cet environnement, puis observe le résultat de cette action avant de décider du pas suivant, jusqu'à terminaison de l'épisode. Le cadre \emph{ReAct} \cite{yao2023react} en constitue l'instanciation canonique : il alterne explicitement pensées (\texttt{Thought}), actions (\texttt{Action}) et observations (\texttt{Observation}) au sein d'une même séquence autorégressive. Ce cadre distingue formellement l'agentique du raisonnement mono-tour, qui évalue le modèle sur un problème statique à information complète, résolu en un seul rollout autorégressif. L'agentique relève au contraire du processus de décision markovien partiellement observable, le cadre usuel de l'apprentissage par renforcement : le modèle n'accède jamais à l'état complet du système et ne dispose que de la trace des observations accumulées dans son contexte. La difficulté réside alors dans la gestion d'un horizon décisionnel long, où chaque action modifie l'état du système et conditionne les observations ultérieures. Qu'il s'agisse de jeux textuels et d'environnements incarnés \cite{shridhar2020alfworld}, d'interfaces web \cite{zhou2023webarena} ou de bases de code réelles \cite{jimenez2023swebench}, ces environnements se caractérisent par leur nature interactive, évolutive et médiée par des outils.

Ce passage au multi-tour ne modifie pas seulement la nature de la tâche, il modifie la difficulté même de l'optimisation. La trajectoire $\tau$ n'est plus intégralement générée par la politique : elle contient des observations produites par l'environnement, sur lesquelles aucun gradient ne doit être rétropropagé. Il faut donc un masquage analogue à celui du SFT, appliqué cette fois à l'échelle de la trajectoire. Surtout, l'horizon $T$ croît de plusieurs ordres de grandeur tandis que la récompense demeure généralement terminale. L'attribution de crédit, c'est-à-dire l'identification des actions responsables du succès ou de l'échec d'un épisode, devient considérablement plus ardue que dans le cas d'une réponse unique. Les recettes de stabilisation validées en RL mono-tour ne se transfèrent donc pas telles quelles au régime interactif. Elles fournissent un point de départ raisonnable, mais doivent y être revalidées ; la section 2.3 précise ces instabilités et la section 4.2 en documente plusieurs exemples.

L'interaction apporte en retour un avantage propre, et la littérature offre des réponses complémentaires à la tension décrite ci-dessus. D'une part, la boucle de rétroaction de l'environnement, par ses messages d'erreur et ses observations d'état, compense partiellement les déficits de planification \emph{a priori} des LLMs \cite{yao2023react, shinn2023reflexion, kambhampati2024llmmodulo} ; la section 2.5 revient sur cette asymétrie entre boucle ouverte et boucle fermée. D'autre part, une partie de la littérature densifie le signal en estimant une récompense à l'échelle de l'étape (\emph{step-level credit}), au moyen de modèles de récompense de processus (\emph{Process Reward Models}, PRM) \cite{lightman2023verify}, au prix d'un coût d'annotation ou d'estimation supplémentaire. C'est précisément cette tension entre horizon, sparsité du signal et attribution de crédit qui motive l'étude des algorithmes de la section 1.2 dans le contexte agentique.

Au sein de cet écosystème, un modèle de langage n'agit jamais de manière isolée : il reste indissociable de son harnais, la superstructure logicielle qui l'entoure, de la gestion du contexte aux parseurs d'actions, en passant par les protocoles de communication comme MCP \cite{anthropic2024mcp} et les mécanismes de vérification. Il convient de distinguer les workflows rigides, où les chemins d'exécution sont codés en dur, des agents autonomes où le modèle dirige lui-même sa propre boucle de décision \cite{anthropic2024agents}. Les notes méthodologiques des benchmarks récents montrent que les performances mesurées dépendent autant de cet environnement d'exécution et du budget de tokens alloué que des poids du modèle \cite{openai2024swebenchverified}. Cette interdépendance conduit les laboratoires à livrer des paires modèle-harnais développées conjointement, à l'image du co-design entre matériel et logiciel, le modèle étant explicitement post-entraîné sur ses propres outils et prompts système \cite{anthropic2024agents, kimik2_2025}. Le harnais est enfin devenu un objet d'optimisation à part entière : les systèmes industriels s'appuient sur l'orchestration dynamique de sous-agents pour paralléliser la décomposition de tâches complexes \cite{kimik2_2025}, et la recherche explore la découverte automatique d'architectures agentiques, où le code du harnais, gestion de la mémoire, filtrage du contexte ou génération de sous-agents, est lui-même appris ou évolué \cite{hu_adas_2025}.

\subsection{Construction des Environnements et Benchmarks Agentiques}

Si la section précédente a défini l'agent comme une politique en boucle d'interaction, la définition du problème qu'il doit résoudre est un enjeu au moins aussi déterminant. C'est elle qui conditionne à la fois ce que l'on mesure et ce que le modèle peut apprendre. Un environnement agentique se spécifie formellement comme un quadruplet : un état initial (le contenu du monde au début de l'épisode), un ensemble d'outils (les actions exposées à l'agent, aujourd'hui de plus en plus médiées par le protocole MCP), une tâche exprimée en langage naturel, et un vérificateur exécutable qui décide du succès. Chacune de ces composantes est un choix de conception lourd de conséquences. Un vérificateur fondé sur l'exécution, par comparaison d'états finaux ou scripts de test, mesure autre chose qu'une comparaison textuelle à une référence. Un état initial pauvre rend la tâche triviale quelle que soit la sophistication des outils.

Les benchmarks académiques illustrent chacun un axe de difficulté distinct, et leur lecture conjointe dessine une décomposition opérationnelle de la compétence agentique. $\tau$-bench et son extension $\tau^2$-bench \cite{taubench2024, tau2bench2025} introduisent la dimension conversationnelle et normative : l'agent doit résoudre la requête d'un utilisateur simulé, lui-même doté d'outils dans le régime de contrôle dual, tout en respectant une politique de domaine (compagnie aérienne, télécoms). Les performances s'effondrent dès que l'utilisateur doit être activement guidé, avec un pass@1 de l'ordre de 34\,\% en télécoms pour des modèles frontière. Toolathlon \cite{toolathlon2025} cible la diversité et le réalisme des états initiaux : 32 applications réelles et 604 outils MCP, 108 tâches d'environ 20 tours nécessitant de coordonner plusieurs applications sur des états de départ authentiques, tels que des cours Canvas peuplés de dizaines d'étudiants ou des feuilles de calcul financières réelles. Le meilleur modèle évalué n'y dépasse pas 38,6\,\% de réussite. MCP-Atlas \cite{mcpatlas2026} et MCP-Universe \cite{mcpuniverse2025} isolent quant à eux la découverte d'outils : les énoncés ne mentionnent ni serveur ni outil, ce qui contraint l'agent à identifier les bons instruments parmi des distracteurs sémantiquement plausibles, puis à composer des flux multi-serveurs. Ces scores, systématiquement inférieurs à 55\,\% pour les modèles frontière, indiquent que le facteur limitant n'est plus seulement la capacité brute de raisonnement. C'est aussi la conjonction entre horizon long, observabilité partielle, ambiguïté de la surface d'outils et vérification stricte, c'est-à-dire précisément ce que la définition de l'environnement encode.

Du côté de l'entraînement, les rapports techniques des constructeurs révèlent que ces environnements ne sont plus seulement des instruments d'évaluation, mais la matière première du post-training. L'équipe Qwen décrit une trajectoire éloquente : une vingtaine de tâches RL généralistes pour Qwen3 \cite{qwen3_2025}, puis 20\,000 environnements d'exécution parallèles pour le RL agentique de Qwen3-Coder, dont le rapport désigne l'\emph{environment scaling} comme le défi central \cite{qwen3coder2025}, jusqu'à Qwen3.5, dont les gains sont attribués au passage à l'échelle de « pratiquement toutes les tâches et tous les environnements de RL concevables », avec une distribution de tâches à complexité progressive \cite{qwen35_2026}. Kimi K2 systématise cette logique par un pipeline de synthèse en trois étapes, de la constitution d'un répertoire de spécifications d'outils réels et synthétiques à la production de trajectoires multi-tours vérifiées, alimentant un étage de RL conjoint sur environnements simulés et réels \cite{kimik2_2025}. Mistral rapporte de même le rôle d'environnements RL dédiés dans l'entraînement de ses modèles de code agentiques \cite{magistral2025}.

Cette insistance sur la diversité n'est pas rhétorique. Elle répond à un mode d'échec bien identifié du RL agentique : une politique optimisée contre un environnement fixe apprend à exploiter les particularités de sa surface d'outils, de son vérificateur ou de ses gabarits de tâches, de sorte que les gains mesurés ne se transfèrent pas hors de la distribution d'entraînement. La multiplicité des environnements ne joue donc pas le rôle d'un simple élargissement de couverture, mais celui d'un régularisateur : c'est la variété des vérificateurs et des dynamiques qui empêche la politique de sur-ajuster un simulateur particulier et force l'apprentissage de compétences décisionnelles transférables. Ce constat ouvre un axe de recherche en plein essor, qui consiste à traiter l'environnement lui-même comme un artefact à générer, qu'un agent propose ses propres tâches et leurs vérificateurs \cite{zhou2025sca} ou qu'un modèle de monde synthétise intégralement les expériences d'entraînement \cite{dreamgym2025}. Il prolonge directement les approches de génération synthétique de données évoquées en introduction. La génération synthétique d'environnements est du reste une pratique établie en apprentissage robotique : l'agent y évolue dans un jumeau numérique, et l'algorithme ajuste ou régénère l'environnement en fonction des performances courantes de l'agent, jusqu'à la co-évolution ouverte d'environnements et de politiques \cite{wang2019poet}. Cette adaptation de l'environnement aux capacités de l'apprenant relie la génération d'environnements au méta-apprentissage et au curriculum, concepts détaillés en section 2.4.

Pour la présente étude, ce panorama fixe deux exigences et délimite un positionnement. Première exigence : comprendre finement l'anatomie de l'environnement de travail, soit pour TextCraft son espace d'actions compositionnel, son vérificateur exact et sa gradation de difficulté par profondeur de recettes. C'est un préalable à toute interprétation des dynamiques d'apprentissage, car un gain mesuré n'a de sens que rapporté à ce que le vérificateur mesure réellement. Seconde exigence : la valeur d'un résultat obtenu dans un environnement contrôlé se juge à ce qui se transfère hors de sa distribution, faute de quoi l'on ne mesure que l'ajustement à un simulateur. Quant au positionnement, notre sujet n'est pas de générer des environnements, mais de comprendre comment optimiser l'apprentissage d'un agent en leur sein, et d'y étudier les idées d'auto-amélioration et de méta-apprentissage. La génération d'environnements à grande échelle relève avant tout d'une ingénierie logicielle que seuls les constructeurs peuvent soutenir ; l'étude des dynamiques d'apprentissage dans ces environnements constitue la question de recherche, et c'est elle qui correspond au sujet du stage. Ce double constat motive le choix d'un cadre expérimental délibérément restreint mais finement instrumenté (section 3), et l'attention portée aux propriétés de stabilisation et de transfert plutôt qu'aux seuls scores absolus (section 4).

\subsection{Apprentissage par Renforcement en Environnement Interactif : Stabilité et Mesure}

Les algorithmes présentés en section 1.2 ont été conçus et validés dans des régimes mono-tour : contrôle continu pour PPO, alignement RLHF puis raisonnement RLVR pour ses variantes appliquées aux LLMs. Leur transposition aux environnements interactifs est récente, et deux questions transversales structurent la littérature comme notre travail : comment maintenir la stabilité de cet apprentissage, et comment mesurer ce qu'il apprend réellement. Cette sous-section pose ces deux points, que les sections 4.2 et 4.5 développeront expérimentalement.

Le RL profond à récompense sparse est notoirement instable, et le régime multi-tour aggrave chacune de ses pathologies connues. Le détournement de récompense (\emph{reward hacking}) désigne le cas où la politique optimise la récompense mesurée en divergeant de l'intention du concepteur \cite{skalse2022defining} ; les récompenses intermédiaires ajoutées pour densifier le signal en sont un vecteur classique. La dérive de politique désigne l'éloignement progressif de la politique entraînée par rapport à ses capacités d'origine ; la pénalité de divergence KL vers une politique de référence, héritée du RLHF \cite{ouyang2022training}, en est le garde-fou standard, mais le choix de la référence et du coefficient devient critique sur les longs horizons. Enfin, le collapse d'entropie désigne l'effondrement précoce de la diversité des sorties de la politique : Cui et al. \cite{cui2025entropy} montrent que l'entropie chute dès le début de l'entraînement RL, que la performance se paie en entropie, et que la saturation de l'une accompagne l'épuisement de l'autre. Sans intervention, la politique devient quasi déterministe et l'exploration s'éteint, ce qui prive les algorithmes de groupe comme GRPO de tout signal. Ces trois pathologies ne sont pas des curiosités théoriques : nous les avons toutes rencontrées, et la section 4.2 présente les pratiques de stabilisation que cette étude a permis d'identifier, du masquage des observations à l'ancrage KL mobile.

Mesurer ce que le RL apprend exige ensuite un instrument plus fin que le taux de réussite en un tirage. Le pass@$k$, fraction des tâches résolues par au moins un tirage parmi $k$, approche pour $k$ grand le support de la politique, c'est-à-dire l'ensemble des problèmes que le modèle sait résoudre au moins occasionnellement. Cet instrument fonde une thèse forte sur la nature du RLVR : Yue et al. \cite{yue2025rlvr} observent que les modèles entraînés dominent leur modèle de base à petit $k$ mais lui deviennent inférieurs à grand $k$, et concluent que le RL concentre la masse de probabilité sur les chemins déjà récompensés sans étendre la frontière de compétence, en écho au mécanisme de \emph{sharpening} \cite{huang2024self} présenté en introduction. Cette thèse est discutée : ProRL~\cite{liu2025prorl} observe, avec un entraînement bien plus long, que les modèles entraînés dépassent leur modèle de base à tous les budgets de tirages, et attribue les conclusions précédentes à la brièveté des entraînements étudiés.
 La question de l'attribution s'y ajoute : ce qu'un post-entraînement apporte dépend de ce que le modèle de base contient déjà, le déclenchement de capacités dites émergentes étant conditionné par la présence préalable des structures correspondantes \cite{liu2025drgrpo}, et les générations récentes de modèles ouverts étant massivement exposées aux traces agentiques dès le mid-training \cite{qwen3coder2025, kimik2_2025}. Conformément au positionnement de ce rapport, nous traitons la thèse du plafond pass@$k$ comme une hypothèse à mettre à l'épreuve et non comme un acquis ; la section 4.5 la confronte à nos mesures, ainsi qu'à nos ablations inter-générations (Qwen2.5, Qwen3, Qwen3.5).

\subsection{Curriculum Learning, Autocurriculum et Agents Autotéliques}


Le \emph{curriculum learning} désigne une stratégie d'entraînement dans laquelle les exemples sont présentés au modèle par difficulté croissante, plutôt que dans un ordre arbitraire \cite{bengio2009curriculum}. Dans le cadre du réentraînement par renforcement d'agents LLM, cette idée prend un sens très concret. La récompense y est terminale et binaire : sur une tâche trop difficile pour la politique courante, tous les rollouts d'un groupe échouent, l'avantage relatif est nul et GRPO ne produit aucun gradient ; sur une tâche déjà maîtrisée, tous réussissent et le gradient est nul également. L'ordre dans lequel les tâches sont présentées détermine donc directement la quantité de signal d'apprentissage disponible à chaque étape. C'est ce qui fait du curriculum un levier candidat pour stabiliser et accélérer l'apprentissage face à des horizons longs et des récompenses sparses, et la littérature récente des agents auto-améliorants le répertorie comme l'un de ses paradigmes constitutifs \cite{gao2025selfevolving}.

Une première famille de curriculums est conçue par l'utilisateur (\emph{user-designed}) : l'ordre de présentation des tâches ou les capacités du modèle suivent une heuristique de difficulté fixée avant l'entraînement. La méthode ScalingInter d'AgentGym-RL \cite{xi_agentgym-rl_2025} en est un exemple sur l'axe de l'horizon : le nombre maximal de tours d'interaction est d'abord restreint, ce qui force la politique à exploiter des trajectoires courtes, puis étendu par paliers. Le séquencement par complexité du domaine en est un second, sur l'axe de la tâche elle-même : dans TextCraft, entraîner d'abord sur les recettes de faible profondeur, puis élargir progressivement aux profondeurs supérieures. Ces schémas ont l'avantage de la simplicité et de la lisibilité. Ils ont aussi deux limites. La difficulté qu'ils encodent est celle perçue par le concepteur humain, qui ne coïncide pas nécessairement avec la difficulté effective pour le modèle. Et leur calendrier est fixé à l'avance, aveugle aux dynamiques d'apprentissage réelles : un palier franchi trop tôt prive la politique de signal, un palier maintenu trop longtemps la fait stagner sur des tâches acquises.

Ces limites motivent la seconde famille, l'\emph{autocurriculum}, où la sélection des tâches est ajustée en continu à partir des performances de la politique elle-même. L'idée fondatrice est de sélectionner les tâches selon le progrès d'apprentissage (\emph{learning progress}) qu'elles procurent, c'est-à-dire la vitesse à laquelle la compétence du modèle y évolue \cite{graves2017automated}. L'intuition rejoint la notion de zone proximale de développement : les tâches les plus informatives sont celles de difficulté intermédiaire, ni acquises ni hors de portée, précisément celles où les groupes de rollouts sont mixtes et où GRPO reçoit du signal. Cette approche a récemment reçu des garanties théoriques dans le contexte des LLMs : un autocurriculum réduit de manière exponentielle le besoin de démonstrations supervisées en SFT, et découple le coût du RL de la qualité du modèle de référence \cite{rajaraman2026curriculum}.

L'autocurriculum suppose encore un ensemble de tâches donné, dans lequel l'agent choisit. Le cadre des agents autotéliques franchit un pas de plus : l'agent y définit, poursuit et organise ses propres objectifs au sein d'environnements ouverts \cite{colas2022autotelic, forestier2017imgep}. La sélection de buts devient alors le cœur du problème, et elle se heurte à une difficulté d'échelle : lorsque l'espace des buts est très vaste, estimer la compétence de l'agent sur chaque but par comptage de ses succès devient impraticable, chaque but n'étant visité que rarement. Le framework MAGELLAN \cite{gaven2025magellan} répond à cette difficulté par une approche métacognitive : un estimateur appris, branché sur les représentations internes du LLM, prédit la probabilité de succès de l'agent sur chaque but, et le progrès d'apprentissage est mesuré comme l'évolution de cette prédiction dans le temps. Parce que des buts sémantiquement proches ont des représentations proches, l'estimateur généralise entre buts et fournit une carte de compétence dense à partir d'observations rares. C'est ce mécanisme que nous portons et évaluons en section 4. La même logique de sélection à la frontière des capacités s'étend aux architectures hiérarchiques : dans HERAKLES \cite{carta2025herakles}, une politique haut niveau choisit des sous-buts atteignables que la politique bas niveau apprend à compiler en compétences réutilisables, ce qui constitue un curriculum appliqué aux buts que l'agent se donne lui-même.

Notre étude se situe dans la première famille et emprunte à la troisième son instrument de mesure. Nous comparons, à recette d'entraînement identique, trois curriculums conçus par l'humain sur trois leviers de TextCraft : le nombre de tours d'interaction, le budget de génération par tour et la profondeur des recettes (section~4.3). Nous avons aussi porté MAGELLAN et couru un autocurriculum sur la même recette ; son résultat, négatif et instructif, est présenté en section~4.5 et en annexe~E. Il instruit une question que les schémas manuels laissent ouverte : la difficulté perçue par l'humain, ici la profondeur, est-elle la vraie structure de compétence du modèle ? La perspective plus lointaine, un générateur d'environnements dont la difficulté suivrait les capacités courantes du modèle, rejoint l'axe de génération d'environnements décrit en section~2.2 ; elle est discutée en ouverture (section~5).


\subsection{Transfert de Compétences : Planification Mono-tour et Contrôle Interactif Multi-tour}


Le transfert entre la planification mono-tour en boucle ouverte (\emph{open-loop planning}) et l'exécution interactive multi-tour en boucle fermée (\emph{closed-loop control}) représente un enjeu à la fois économique et algorithmique. D'un point de vue computationnel, l'entraînement par renforcement en mono-tour est nettement plus frugal : il ne nécessite qu'un seul rollout autorégressif par problème, ce qui supprime les dizaines d'allers-retours avec le serveur d'environnement qui dominent le temps de chaque étape d'optimisation en multi-tour. Si des compétences de planification abstraite acquises en mono-tour se transféraient au contrôle interactif, on disposerait d'un régime de pré-alignement à bas coût, susceptible d'accélérer la convergence du RL multi-tour et d'en réduire la facture de calcul.

La littérature sur les capacités de planification des LLMs invite toutefois à la prudence sur le premier terme de ce transfert. Les travaux de Valmeekam, Kambhampati et leurs collaborateurs \cite{valmeekam2023planning, kambhampati2024llmmodulo} établissent que la génération autonome de plans complets et valides reste une faiblesse structurelle des LLMs : le raisonnement transitif sur de longues chaînes de dépendances se dégrade rapidement avec la profondeur du problème. Symétriquement, la littérature du contrôle interactif, de ReAct \cite{yao2023react} aux approches à rétroaction verbale \cite{shinn2023reflexion}, montre que la boucle de retour de l'environnement compense partiellement ces déficits de planification \emph{a priori}. Les messages d'erreur et l'observation de l'état agissent comme un support dynamique que le plan abstrait ne peut fournir seul.

La question du transfert est donc doublement ouverte. Dans le sens mono-tour vers multi-tour : un pré-entraînement en planification accélère-t-il le contrôle interactif ? Dans le sens inverse : le RL interactif améliore-t-il la planification abstraite ? On peut également concevoir des régimes hybrides alternant phases de planification en boucle ouverte et phases d'exécution interactive. L'annexe~D rapporte les mesures préliminaires que nous avons pu faire sur TextCraft, et la section~5 reprend la question.


Ce couplage entre apprentissage par observation et apprentissage par action dépasse enfin notre cas d'étude et rejoint une question posée par les sciences cognitives. Dupoux, LeCun et Malik \cite{dupoux2026why} partent du même constat de mur des données que notre introduction, et proposent une architecture d'apprentissage autonome inspirée du développement humain : un système qui apprend par observation, un système qui apprend par action en environnement, et un méta-contrôleur qui orchestre les deux. La planification mono-tour en boucle ouverte et le contrôle interactif en boucle fermée en constituent une instanciation naturelle dans le cadre des agents LLM. Un régime qui alternerait ces deux phases et dont l'orchestration serait elle-même apprise relèverait du méta-apprentissage, au même titre que l'autocurriculum de la section 2.4. Nous n'avons pas eu le temps de développer cette direction au-delà des mesures préliminaires de l'annexe~D ; elle constitue à nos yeux l'un des prolongements les plus naturels de ce travail.

\subsection{Autres Pistes Considérées}

Deux directions complémentaires ont été étudiées au cours du stage, puis dépriorisées au profit du curriculum ; elles ne relèvent pas directement du méta-apprentissage et leur développement est renvoyé en annexe. La première est le \emph{scaling} du calcul à l'inférence : plutôt que d'entraîner davantage, allouer plus de calcul au moment de la génération. Ce levier va de l'apprentissage en contexte, qui adapte le comportement du modèle par simple injection d'exemples sans modification des poids \cite{brown2020gpt3}, aux stratégies d'échantillonnage multiple et de recherche, dont les analyses de Snell et al. \cite{snell2024scaling} fournissent le cadre d'arbitrage : à coût de calcul égal, un modèle plus petit doté d'une stratégie d'inférence adaptée à la difficulté peut surpasser un modèle plus grand interrogé naïvement. Nos expériences few-shot en section 4 relèvent de ce levier, et la frontière entre calcul d'inférence et génération de données d'entraînement s'estompe dans les systèmes générateurs de tâches évoqués en section 2.2 \cite{zhou2025sca, dreamgym2025}.

La seconde piste est la recomposition de trajectoires par échantillonnage préférentiel auto-normalisé (\emph{Self-Normalized Importance Sampling}, SNIS) \cite[chap.~9]{owen2013montecarlo}, formalisée à la suite d'échanges avec Otmane Sakhi (Criteo AI Lab, équipe Performance Science). L'idée consiste à réutiliser hors-politique les sous-composantes des trajectoires d'un même groupe GRPO, en recomposant chaînes de pensée et réponses finales repondérées par des poids auto-normalisés ; le mécanisme exploite le fait qu'évaluer la vraisemblance d'une recomposition coûte une seule passe avant, là où générer une trajectoire nouvelle est coûteux. La transposition au multi-tour se heurte à un verrou structurel que nos essais ont confirmé : l'état intègre l'historique strict des observations du simulateur, et recoller des sous-trajectoires issues de contextes discordants viole la cohérence causale. La formalisation des poids, la dérivation de l'estimateur et le bilan de nos observations sont regroupés en annexe B, avec la variante guidée par un modèle réviseur comme condition d'une reprise éventuelle.











%------------------------------------------------------------------------------------------------------------








\newpage
\section{Cadre Technique : AgentGym-RL et TextCraft}


\subsection{L'Écosystème AgentGym : une Plateforme Unifiée d'Environnements}

Les exigences dégagées en section 2, diversité des environnements, vérificateurs exécutables et mesure du transfert, plaident pour un cadre expérimental mutualisé plutôt que pour le développement d'un simulateur ad hoc. Ce travail s'appuie sur l'écosystème AgentGym, construit en deux temps par la même équipe (Fudan NLP). Le premier papier, AgentGym \cite{xi_agentgym_2025}, publié à ACL 2025, répond à l'absence de plateforme unifiée pour \emph{évaluer et entraîner} des agents LLM. Il fédère 14 environnements couvrant 89 tâches réparties en sept grands scénarios, de la navigation web (WebShop, WebArena) aux jeux textuels et de crafting (TextCraft, MAZE, Wordle), en passant par la planification incarnée (ALFWorld, BabyAI), le raisonnement scientifique (SciWorld) et l'utilisation d'outils applicatifs (Weather, Movie, Academia, Sheet, TODOList, BIRD), derrière une API HTTP uniforme qui découple strictement l'agent des serveurs d'environnements. Au-delà de l'infrastructure, la contribution comprend une base d'instructions élargie, une suite d'évaluation (AgentEval), un corpus de trajectoires expertes (AgentTraj) et une méthode d'auto-amélioration, AgentEvol, qui fait évoluer l'agent au-delà des démonstrations par exploration et imitation filtrée de ses propres succès, héritière directe des approches de type ReST présentées en introduction. L'ambition affichée est celle d'agents \emph{généralistes} : un même modèle entraîné et évalué sur l'ensemble des environnements de la suite.

Le second papier, AgentGym-RL \cite{xi_agentgym-rl_2025}, accepté comme présentation orale à ICLR 2026 \cite{xi_agentgymrl_openreview}, prolonge cette infrastructure vers l'apprentissage par renforcement multi-tour de bout en bout. Le premier AgentGym ne proposait que des méthodes d'entraînement supervisées, et son mécanisme d'auto-amélioration, AgentEvol, relevait de l'imitation filtrée des propres succès de l'agent. AgentGym-RL justifie le passage au renforcement en ligne en relevant que les tentatives antérieures de RL multi-tour « peinent à atteindre la stabilité et l'efficacité d'optimisation » ; la stabilité de cet apprentissage est donc, dès l'origine de la plateforme, la question centrale.
Le constat de départ rejoint exactement notre cadrage : la communauté manque d'un cadre RL interactif unifié capable d'entraîner des agents \emph{from scratch}, sans étape initiale de SFT sur démonstrations expertes, coûteuses et peu généralisables. L'architecture modulaire et découplée permet des rollouts distribués, concurrents et asynchrones, condition du passage à l'échelle du Deep RL en environnement interactif, et supporte les principaux algorithmes d'optimisation de politique (PPO, GRPO, variantes). L'entraînement RL porte sur un sous-ensemble de cinq scénarios et 27 tâches : navigation web, recherche approfondie, tâches incarnées et sciences, jeux dont TextCraft (Minecraft Textuel). Pour contenir l'instabilité du RLVR sur longs horizons, les auteurs introduisent ScalingInter-RL, un curriculum sur la durée d'interaction : l'entraînement débute avec un nombre maximal de tours restreint, ce qui force l'exploitation sur trajectoires courtes, puis l'horizon est progressivement étendu à mesure que la politique se stabilise. Les résultats sont marquants : un Qwen2.5-7B entraîné par ScalingInter-RL atteint 91\,\% de réussite moyenne sur la suite, égalant ou dépassant des modèles propriétaires de premier plan (GPT-4o, Gemini-2.5-Pro, o3) évalués sur les mêmes tâches \cite{xi_agentgym-rl_2025}.

\subsection{Positionnement Critique et Justification du Choix}

Ces résultats appellent cependant une lecture nuancée, qui délimite précisément l'espace de notre contribution. D'abord, la comparaison oppose de petits modèles optimisés directement sur la distribution exacte des environnements d'entraînement à des modèles généralistes évalués à froid, sans adaptation au format d'action ni au comportement du simulateur. Cette asymétrie d'évaluation relativise l'équivalence affichée. La démonstration de faisabilité reste néanmoins économiquement significative pour le déploiement de modèles spécialisés frugaux : concentrer la politique d'un modèle léger sur de longues séquences décisionnelles, avec un transfert partiel vers des tâches non vues. Ensuite, le processus d'évaluation par les pairs, consultable publiquement \cite{xi_agentgymrl_openreview}, situe la contribution davantage du côté de l'infrastructure que de l'algorithmique : ScalingInter-RL est une heuristique de curriculum fixée a priori, dont le gain sur un entraînement à horizon constant varie selon les tâches, et dont le mécanisme n'est ni analysé ni comparé à des alternatives adaptatives. Il faut enfin mesurer l'asymétrie de moyens : ce travail collectif mobilise plusieurs dizaines d'auteurs et une infrastructure multi-GPU, quand la présente étude est une réplication et une extension menées par une seule personne sur un GPU unique. C'est cette contrainte qui donne sa valeur pratique à l'étude de stabilisation de la section 4.2. Il faut y ajouter un point de méthode. Le papier ne rapporte ni le nombre de mises à jour ni le nombre d'époques de ses runs, ne donne pas de configuration pour le modèle 3B, ne publie ni courbe d'entraînement ni trace, et son texte et son code divergent sur l'horizon d'interaction (section~4.2.2). Une des motivations de notre étude est d'apporter cette transparence : rendre publics les journaux, les recettes et les échecs, pour que les résultats puissent être questionnés, les constats transmis, et de nouvelles questions posées. C'est aussi pourquoi nous avons commencé par une réplication à l'identique, avant d'ouvrir vers nos propres questions : sans elle, le défaut de masquage des observations décrit en section~4.2.1 serait probablement passé inaperçu.


Cette lecture critique n'épuise pourtant pas l'intérêt de l'idée. Faire porter le curriculum sur un paramètre du harnais, ici le budget d'interaction, plutôt que sur la sélection des données elles-mêmes, reste rare dans la littérature du curriculum, qui se concentre généralement sur l'ordre des tâches. L'idée nous semble généralisable au-delà de AgentGym-rl : à coût de calcul égal, des épisodes courts seront toujours préférables, et l'horizon utile sera toujours un bon proxy de difficulté des tâches, de sorte que la règle est plausiblement applicable à tout environnement interactif. C'est d'ailleurs le nombre de tours qu'un modèle propriétaire de référence met à résoudre une tâche, proxy que Toolathlon utilise pour graduer ses tâches. Nos expériences la prolongent d'ailleurs sur un second paramètre du même type, le budget de génération par tour (section 4).

ScalingInter-RL n'est qu'un point dans l'espace des régimes de structuration de l'apprentissage décrit en section 2.4 : un curriculum \emph{user-designed}, statique, portant sur une seule variable, l'horizon. L'environnement TextCraft en expose naturellement une seconde, la profondeur des recettes, et la littérature autotélique fournit le cadre d'une troisième famille, adaptative, pilotée par le progrès d'apprentissage, que nous avons explorée sans la mener au bout (annexe~E). Comparer systématiquement les régimes de curriculum à recette d'entraînement identique et à compute contrôlé, ce que le papier ne fait pas, constitue le cœur de notre second travail. Par ailleurs, la plateforme satisfait nos exigences méthodologiques : récompense vérifiable exacte, simulateur peu coûteux autorisant des ablations répétées, et multiplicité d'environnements offrant, à terme, le terrain de mesure du transfert hors distribution exigé en section 2.2.




\subsection{TextCraft : Anatomie du Banc d'Essai, Dataset et Questions Ouvertes}

TextCraft \cite{prasad2023adapt} est un jeu de crafting textuel inspiré de Minecraft, introduit initialement pour évaluer la décomposition de tâches par les LLMs. L'agent dispose de trois actions élémentaires : \texttt{get} pour récupérer un ingrédient de base, \texttt{craft} pour combiner des ingrédients selon une recette, et \texttt{inventory} pour consulter son inventaire. Il doit fabriquer un objet cible dont la recette dépend récursivement d'autres objets. Chaque tâche se caractérise par la \emph{profondeur} de son arbre de dépendances : à profondeur 1, l'objet se fabrique directement à partir d'ingrédients de base ; à profondeur $d$, il requiert la fabrication préalable d'objets de profondeur $d-1$. La figure~\ref{fig:textcraft_episode} présente un épisode annoté, avec le prompt système, l'observation initiale listant les commandes de craft disponibles, et l'alternance de pensées et d'actions de l'agent.

\begin{figure}[htbp]
\centering
\begin{tcolorbox}[
    colback=gray!4, colframe=black!65,
    title=\textbf{Exemple de tâche TextCraft --- prompt système et objectif},
    fonttitle=\small, width=0.96\textwidth,
    boxrule=0.6pt, arc=2mm, left=3mm, right=3mm, top=2mm, bottom=2mm]
{\small\ttfamily
You are given few useful crafting recipes to craft items in Minecraft. Crafting commands are of the format \mbox{"craft} [target object] using [input \mbox{ingredients]"}. Every round I will give you an observation, you have to respond an action based on the state and instruction. You can "get" an object (ingredients) from the inventory or the environment, look-up the game inventory by "inventory", or "craft" (target) using any of the crafting commands. You can use ONLY these crafting commands provided, do not use your own crafting commands. However, if the crafting command uses a generic ingredient like "planks", you can use special types of the same ingredient e.g. "dark oak planks" in the command instead.

\medskip
\textbf{Goal: craft flint and steel.}
}
\end{tcolorbox}
\caption{Prompt système de l'environnement TextCraft et objectif type, tels que fournis à l'agent en début d'épisode (exemple tiré de \cite{xi_agentgym-rl_2025}). L'agent répond ensuite tour par tour par une action \texttt{get}, \texttt{inventory} ou \texttt{craft}, l'environnement renvoyant une observation après chaque action, jusqu'à vérification de l'objet cible dans l'inventaire.}
\label{fig:textcraft_episode}
\end{figure}


L'instance utilisée compte 544 recettes, que nous partitionnons en trois ensembles : 374 recettes d'entraînement, 100 recettes de test et un réservoir de 70 recettes réservées à la construction des exemples few-shot, ce qui garantit l'absence de fuite entre démonstrations et évaluation. La distribution par profondeur est fortement déséquilibrée : l'entraînement compte 109 recettes de profondeur 1, 221 de profondeur 2, 43 de profondeur 3 et une seule de profondeur 4 ; le test en compte respectivement 31, 41, 25 et 3. Ce déséquilibre est une limite du dataset que le papier ne discute pas : avec une seule recette de profondeur 4 à l'entraînement, l'échec éventuel d'une méthode sur les tâches profondes est confondu avec la quasi-absence de support d'entraînement. Nous n'avons pas re-stratifié cette distribution, ce qui aurait constitué une étude à part entière, mais nous la gardons présente dans toutes nos interprétations, et la re-stratification par difficulté figure parmi les travaux futurs les plus directement actionnables (section 5).

Ce banc d'essai concentre, sous forme minimale, les propriétés identifiées en section 2.2 : un espace d'actions compositionnel dont la difficulté est \emph{paramétrable par construction}, la profondeur étant une variable de curriculum naturelle ; un vérificateur exact, l'objet cible figurant ou non dans l'inventaire final ; une observabilité partielle réelle, l'agent devant consulter l'inventaire et découvrir les recettes ; et un horizon ajustable. Sa légèreté computationnelle, sans rendu et à dynamique déterministe, autorise les ablations répétées qu'exigerait difficilement WebArena. Ces propriétés en font l'environnement adapté pour disséquer la dynamique des gradients GRPO avant extension aux autres environnements de la suite, extension envisagée si les régimes de curriculum se montrent concluants, à la manière de la validation multi-environnements de ScalingInter chez les auteurs.

Le tableau~\ref{tab:agentgym_textcraft} reproduit les résultats de référence du papier sur les 100 tâches de test, qui constituent nos lignes de base externes. Sa lecture par profondeur est la plus instructive. La profondeur 4 est un mur quasi universel : GPT-4o, o3 et la totalité des modèles ouverts de moins de 100 milliards de paramètres antérieurs à Qwen3 y échouent intégralement, et seuls Gemini-2.5-Pro, Qwen3-8B, Qwen3-32B et ScalingInter-7B y résolvent une tâche sur trois. La profondeur 3 sépare nettement les générations et les tailles de modèles. Enfin, la comparaison des trois modèles RL du papier montre que l'essentiel du gain de ScalingInter-7B sur AgentGym-RL-7B se joue en profondeur 3 et 4. Deux précisions de lecture. Le score de 33{,}33\,\% en profondeur 4 correspond à une tâche réussie sur trois. Et notre propre mesure de Qwen2.5-3B-Instruct non entraîné, dans le même harnais, donne 10{,}3\,\% en moyenne de vingt passes (section~4.5), compatible avec les 14\,\% rapportés ici ; le papier ne précise pas son décodage à l'évaluation.


Trois questions ouvertes, posées en section 2, prennent alors leur forme opérationnelle. Le curriculum améliore-t-il la vitesse de convergence, le plafond de performance, ou les deux ? Peut-on atteindre les profondeurs 3 et 4 par une gradation appropriée, malgré le support d'entraînement quasi nul en profondeur 4 ? Et les pass@$k$ du modèle nu constituent-ils un plafond infranchissable pour le RL, comme le soutient la thèse du sharpening, ou une bonne structuration de l'apprentissage permet-elle de les dépasser ?



\begin{table}[htbp]
\centering
\caption{Résultats de référence sur TextCraft (100 tâches de test, succès en \%, par profondeur $d$ et global), d'après la Table~3 de \cite{xi_agentgym-rl_2025}. Meilleur score de chaque colonne en gras.}
\label{tab:agentgym_textcraft}
\small
\setlength{\tabcolsep}{7pt}
\begin{tabular}{lccccc}
\toprule
\textbf{Modèle} & $\mathbf{d{=}1}$ & $\mathbf{d{=}2}$ & $\mathbf{d{=}3}$ & $\mathbf{d{=}4}$ & \textbf{Global} \\
\midrule
\rowcolor{gray!20}\multicolumn{6}{c}{\underline{Modèles propriétaires}} \\
GPT-4o & \textbf{100} & 87{,}80 & 64{,}00 & 0 & 83{,}00 \\
Qwen-Max & 93{,}55 & 75{,}61 & 36{,}00 & 0 & 69{,}00 \\
Gemini-2.5-Flash & \textbf{100} & 95{,}12 & 40{,}00 & 0 & 80{,}00 \\
OpenAI o4-mini & \textbf{100} & \textbf{100} & \textbf{84{,}00} & 0 & 93{,}00 \\
OpenAI o3 & \textbf{100} & \textbf{100} & \textbf{84{,}00} & 0 & 93{,}00 \\
Gemini-2.5-Pro & \textbf{100} & \textbf{100} & \textbf{84{,}00} & \textbf{33{,}33} & \textbf{94{,}00} \\
\midrule
\rowcolor{gray!20}\multicolumn{6}{c}{\underline{Modèles ouverts $\geq$ 100B}} \\
Qwen3-235B-A22B & \textbf{100} & \textbf{100} & \textbf{84{,}00} & 0 & 93{,}00 \\
DeepSeek-V3-0324 & 80{,}65 & 53{,}66 & 40{,}00 & 0 & 57{,}00 \\
DeepSeek-R1-0528 & \textbf{100} & \textbf{100} & \textbf{84{,}00} & 0 & 93{,}00 \\
\midrule
\rowcolor{gray!20}\multicolumn{6}{c}{\underline{Modèles ouverts $<$ 100B}} \\
Qwen2.5-3B-Instruct & 35{,}48 & 7{,}32 & 0 & 0 & 14{,}00 \\
Qwen2.5-7B-Instruct & 80{,}65 & 41{,}46 & 0 & 0 & 42{,}00 \\
Qwen2.5-72B-Instruct & 96{,}77 & 85{,}37 & 48{,}00 & 0 & 77{,}00 \\
Qwen3-4B & 87{,}10 & 36{,}59 & 12{,}00 & 0 & 45{,}00 \\
Qwen3-8B & \textbf{100} & 78{,}05 & 40{,}00 & \textbf{33{,}33} & 74{,}00 \\
Qwen3-32B & 90{,}32 & 92{,}68 & 72{,}00 & \textbf{33{,}33} & 85{,}00 \\
Llama-3.1-8B-Instruct & 74{,}19 & 56{,}10 & 4{,}00 & 0 & 47{,}00 \\
Llama-3.1-70B-Instruct & \textbf{100} & \textbf{100} & \textbf{84{,}00} & 0 & 93{,}00 \\
\midrule
\rowcolor{gray!20}\multicolumn{6}{c}{\underline{Modèles entraînés par RL dans le papier}} \\
AgentGym-RL-3B & \textbf{100} & 90{,}24 & 28{,}00 & 0 & 75{,}00 \\
AgentGym-RL-7B & \textbf{100} & 97{,}56 & 72{,}00 & 0 & 89{,}00 \\
ScalingInter-7B & \textbf{100} & 97{,}56 & 76{,}00 & \textbf{33{,}33} & 91{,}00 \\
\bottomrule
\end{tabular}
\end{table}





%----------------------------------------------------------------------------------------
\newpage

\section{Expériences et Résultats}

Cette section présente le travail expérimental du stage et ses résultats. L'ordre d'exposition est logique plutôt que chronologique : chaque sous-section s'appuie sur les acquis de la précédente.

\begin{itemize}
    \item La sous-section~4.1 fixe le cadre et le protocole communs à toutes les expériences : environnement, modèles, algorithme, évaluation et convention de comparaison.
    \item La sous-section~4.2 est une étude de stabilisation du GRPO multi-tour. Elle part de la réplication du papier en \emph{full finetuning}, puis documente le passage à LoRA, ses instabilités et le correctif que nous proposons, l'ancre KL mobile. C'est le premier de nos deux travaux principaux. LoRA divise par plus de cent la mémoire de l'optimiseur et la taille des points de contrôle. Cette mémoire libérée a permis d'élargir le cache du moteur de génération, d'augmenter la taille des groupes GRPO, et de sauvegarder puis reprendre les runs à faible coût. Le prix en est une instabilité propre, que cette sous-section analyse et corrige.
    \item Les sous-sections~4.3 et~4.4 portent le second travail principal : la comparaison de trois régimes de curriculum, sur l'horizon d'interaction, la profondeur des tâches et le budget de génération, à recette d'entraînement identique puis à compute contrôlé, face à deux références sans curriculum. Le curriculum y est étudié à la fois comme mécanisme de stabilisation et comme question de recherche sur l'exploration.
    \item La sous-section~4.5 regroupe les expériences complémentaires : mesures sans entraînement, expériences non refaites dans la configuration finale, et pistes exploratoires. On y confronte en particulier nos résultats à l'hypothèse du plafond pass@$k$ posée en section~2.3.
\end{itemize}





\vspace{2cm}

\subsection{Protocole Commun}\label{subsec:protocole}

\paragraph{Environnement et données.} Toutes les expériences portent sur TextCraft, dont l'anatomie, les partitions et les limites sont détaillées en section~3.3. Nous en rappelons ici les trois ensembles : 374 recettes d'entraînement, 100 recettes de test fixes, et un réservoir de 70 recettes réservé à la construction des exemples \emph{few-shot}. Aucune recette n'appartient à deux ensembles, ce qui exclut toute fuite entre démonstrations et évaluation. Le déséquilibre par profondeur, en particulier l'unique recette de profondeur 4 à l'entraînement, est une limitation assumée en section~3.3 et non étudiée ici ; nous le rappelons à chaque analyse stratifiée. Une extension sans coût aurait consisté à rendre le réservoir \emph{few-shot} à l'entraînement dans les runs qui n'utilisent pas d'exemples. Nous ne l'avons pas fait, afin de conserver les partitions du papier de référence et de rester comparables à ses scores. Le serveur d'environnement est exécuté localement et l'agent y accède par l'API HTTP d'AgentGym.

\paragraph{Modèles.} Le modèle principal est Qwen2.5-3B-Instruct, la cible du papier de référence, sur lequel portent tous les entraînements. Ce choix tient d'abord au budget de calcul d'un GPU unique. Il rejoint aussi une contrainte applicative : pour un agent de recommandation, la latence d'inférence d'un modèle léger est un critère de déploiement. Deux autres modèles sont évalués sans aucun entraînement, avec le même harnais. Qwen2.5-7B-Instruct mesure l'effet de la taille au sein d'une même génération. Qwen3.5-4B mesure l'effet de la génération à taille comparable. Pour la génération intermédiaire Qwen3, nous n'avons pas de mesure interne et nous appuyons sur le chiffre du papier (tableau~\ref{tab:agentgym_textcraft}). 



\paragraph{Optimisation.} L'algorithme est GRPO dans toutes les expériences, tel que présenté en section~1.2 et en annexe~A : un groupe de $G$ trajectoires par tâche, un avantage relatif calculé au sein du groupe, une récompense terminale binaire, et une pénalité KL vers une politique de référence. En multi-tour s'ajoute le masquage des tokens d'observation, détaillé en section~4.2. Deux recettes structurent la suite.

\begin{itemize}
    \item \textbf{La recette du papier}, répliquée en \emph{full finetuning} à hyperparamètres identiques, d'après son code publié : $G=8$, taux d'apprentissage $10^{-6}$, coefficient KL $10^{-3}$ vers une référence fixe (le modèle de base), température d'échantillonnage $1{,}0$, collectes de 256 trajectoires consommées en quatre pas d'optimisation clippés.
    \item \textbf{La recette stabilisée}, construite pas à pas en section~4.2 : adaptation LoRA de petit rang, taux d'apprentissage constant plus élevé, coefficient KL renforcé, et référence mobile, re-basée périodiquement sur la politique courante.
\end{itemize}

Toutes les mises à jour consomment 64 trajectoires. Dans la phase curriculum, la taille de groupe passe de $G=8$ à $G=16$, soit quatre tâches par mise à jour au lieu de huit. Ce changement est rendu possible par la mémoire libérée par LoRA, réaffectée au cache du moteur de génération. Un banc de débit l'a motivé : à $G=16$ avec le cache élargi, une mise à jour prend 46{,}9~s contre 57{,}1~s dans la configuration historique, soit 18\,\% de moins, les 16 trajectoires d'une même tâche partageant le cache de son prompt. Le banc mélange deux changements, la taille de groupe et le cache, et ne mesure qu'une vitesse. L'effet de $G$ sur l'apprentissage lui-même est contrôlé par un run dédié en section~4.3.



\paragraph{Infrastructure.} Toutes les expériences tournent sur un unique GPU (NVIDIA B200, 192~Go). L'entraînement repose sur la bibliothèque TRL~\cite{vonwerra2020trl}, qui implémente GRPO, et sur le moteur vLLM~\cite{kwon2023vllm}, qui génère les trajectoires et partage la mémoire du GPU avec l'optimisation. Les poids sont synchronisés en mémoire entre les deux à chaque mise à jour. La boucle d'interaction avec l'environnement, tour par tour, est notre code ; le masquage des observations qu'elle produit est décrit en section~4.2. Les points de contrôle périodiques sont conservés sur un stockage volatil, et seuls les meilleurs modèles, avec l'état de leur optimiseur, sont sauvegardés durablement ; l'annexe~C détaille les conséquences de cette contrainte sur le déroulement des runs. 



\textcolor{red}{Ajouter explication sur le tableau d'expérience ic : vous verrez qu'une expe peut avoir plusieurs numeros (20 20.1, 20.2, etc), c'est parce qeu notre infra coder est souvent remise à jour ce qui coupe l'entrainement. On a donc réalisé une save du last et optimizers qui permet derepartir après chaque coupure, et les .1 et .2 etc sont les reprises des expé qui ont été coupées}





\paragraph{Évaluation.} La métrique principale est le pass@$1$ sur les 100 tâches de test, mesuré en cours d'entraînement toutes les 47 mises à jour. Ce pas correspond à une époque à $G=8$ (374 tâches $\times$ 8 trajectoires / 64), donc à une demi-époque à $G=16$. Chaque évaluation rejoue les 100 tâches en épisodes complets, dans des conditions fixes : 30 tours au plus, 512 tokens par tour, température $1{,}0$ sans troncature de la distribution (top-$p$ = 1). Ces conditions sont celles de l'entraînement : l'évaluation mesure la politique telle qu'elle est échantillonnée, et non une version rendue déterministe.

Une évaluation est une proportion sur 100 tâches. Si l'on suppose une même probabilité de succès $p$ pour toutes, son écart-type vaut $\sqrt{p(1-p)/100}$, au plus 5 points. Cette hypothèse est fausse en toute rigueur, chaque tâche ayant sa propre probabilité ; la valeur est alors une borne haute, et elle ne couvre que l'aléa du décodage, pas la variabilité d'un entraînement à l'autre. Nous nous en tenons à cet ordre de grandeur, qui mériterait une étude propre, et n'y revenons pas dans la suite. Il a en revanche deux conséquences sur la façon de lire nos courbes. Le meilleur score d'un run est un maximum pris sur des dizaines d'évaluations bruitées ; il surestime le niveau réel, ce que nous avons vérifié : une reprise d'entraînement initialisée sur un pic isolé retombe au niveau des évaluations voisines. Nous rapportons donc chaque meilleur score avec la moyenne des évaluations qui l'entourent. De même, notre référence pour le modèle non entraîné est la moyenne de vingt passes (10{,}3\,\%), et non la passe unique historique (18\,\%). Le papier de référence ne rapporte ni incertitude ni nombre de passes ; nous préférons expliciter ces choix plutôt que de présenter des maxima bruts.


\paragraph{Métriques stratifiées et diagnostiques.} Au-delà du pass@$1$ global, chaque évaluation périodique enregistre un même ensemble de mesures pour toutes les expériences. Ces mesures servent à lire le comportement du modèle, et non seulement son score : quelles profondeurs il acquiert, quelles erreurs il cesse de commettre, et comment sa dynamique d'entraînement évolue selon le régime de curriculum. Côté évaluation, nous enregistrons le pass@$1$ par profondeur ($d=1$ à $4$), le nombre moyen de tours par épisode, et une taxonomie des erreurs renvoyées par l'environnement, classées d'après son message. Six classes sont distinguées : action mal formée, recette invalide, quantité d'ingrédients insuffisante, objet introuvable, nom d'objet mal écrit, et plusieurs actions dans un même tour. Les deux premières et la dernière relèvent du format ; les autres relèvent de la planification. Côté entraînement, chaque mise à jour enregistre la moyenne et l'écart-type de la récompense par groupe, la fraction de groupes dont toutes les trajectoires ont la même récompense, donc sans gradient, l'entropie de la politique, la divergence KL vers la référence, les longueurs de génération et la norme du gradient. Ces séries alimentent trois analyses : l'ordre d'acquisition des profondeurs, l'ordre d'élimination des classes d'erreurs (section~4.3), et les signaux qui précèdent un collapse (section~4.2).

Une dernière mesure complète le pass@$1$. Le pass@$k$, défini en tête de rapport, est la fraction des tâches résolues au moins une fois en $k$ tirages indépendants. Pour $k$ grand, il approche le \emph{support} de la politique : l'ensemble des tâches qu'elle sait résoudre au moins occasionnellement. Nous le mesurons par un oracle qui tire vingt trajectoires par tâche, à température $1{,}0$, et le rapportons par profondeur (section~4.5). Enfin, chaque configuration correspond à un unique run d'entraînement. Le budget d'un GPU unique n'a pas permis de répéter les runs avec plusieurs graines aléatoires, ce qui aurait mesuré la variabilité d'un entraînement à l'autre. Cette répétition est un travail futur sans difficulté technique. En son absence, la stabilité d'un résultat se lit sur la durée : un niveau tenu sur des dizaines d'évaluations successives, plutôt qu'un pic.

\paragraph{Comparaison à compute contrôlé.} L'axe naturel de comparaison est l'époque, un passage complet sur les tâches d'entraînement. Deux précautions s'imposent. La première tient à $G$ : à 64 trajectoires par mise à jour, une époque compte 46 mises à jour pour $G=8$ et 93 pour $G=16$, soit deux fois plus de trajectoires. Entre tailles de groupe, nous comparons donc en mises à jour ou en trajectoires cumulées, jamais en époques. La seconde tient aux curriculums eux-mêmes. Une époque n'égalise pas le calcul dépensé : un horizon réduit produit des épisodes plus courts, donc moins de tokens générés par trajectoire. C'est un avantage réel des curriculums de budget, et une comparaison en époques le masque. La section~4.4 complète donc chaque comparaison par trois indicateurs relevés dans les journaux d'entraînement.

\begin{itemize}
    \item \textbf{Les tokens générés cumulés.} La génération des trajectoires domine le temps de chaque mise à jour ; c'est l'indicateur le plus proche du calcul réellement consommé, indépendant du matériel.
    \item \textbf{Les heures GPU.} C'est le coût facturé. Il dépend du matériel et de l'infrastructure, mais c'est celui qu'un praticien doit budgéter.
    \item \textbf{Les tours d'environnement.} Chaque tour est un appel au simulateur. Sur TextCraft il est presque gratuit ; dans un environnement réel, appels d'API ou navigation web, il devient un coût à part entière, parfois le premier.
\end{itemize}

\paragraph{Nomenclature.} Les expériences sont désignées dans la suite par des noms courts. Le tableau~\ref{tab:nomenclature} donne la correspondance avec la numérotation interne du registre expérimental, conservée pour la traçabilité. Les expériences secondaires (évaluations \emph{few-shot}, planification mono-tour, pistes SNIS et MAGELLAN) sont nommées directement dans les sections concernées.




\begin{table}[htbp]
\centering
\caption{Nomenclature des expériences principales et correspondance avec la numérotation interne.}
\label{tab:nomenclature}
\small
\begin{tabular}{l p{7.2cm} l}
\toprule
\textbf{Désignation} & \textbf{Contenu} & \textbf{Numéros internes} \\
\midrule
\rowcolor{gray!20}\multicolumn{3}{c}{\underline{Stabilisation (section~4.2)}} \\
\textsc{Réplication-FT} & recette du papier en \emph{full finetuning}, référence KL fixe & exp23 à 23.4 \\
\textsc{Grille-LoRA} & LoRA sur la recette du papier : rang $\times$ taux d'apprentissage $\times$ $\beta$ & exp24 \\
\textsc{Ancre-Mobile} & LoRA $r=8$, référence KL déplacée toutes les 4 époques & exp25, 25.2 \\
\textsc{Carré-Ancre} & référence fixe ou mobile $\times$ $\beta$, autour d'\textsc{Ancre-Mobile} & exp30, 31, 31.1 \\
\midrule
\rowcolor{gray!20}\multicolumn{3}{c}{\underline{Références sans curriculum (sections~4.3 et~4.4)}} \\
\textsc{Ancre-Mobile} & recette stabilisée, $G=8$ & exp25, 25.2 \\
\textsc{Contrôle-G16} & recette stabilisée, $G=16$ & exp36, 36.1 \\
\midrule
\rowcolor{gray!20}\multicolumn{3}{c}{\underline{Curriculums (sections~4.3 et~4.4)}} \\
\textsc{Horizon} & nombre de tours d'interaction croissant par paliers & exp32 \\
\textsc{Profondeur} & profondeur des recettes croissante par paliers & exp33, 33.1, 33.2 \\
\textsc{Budget} & nombre de tokens par tour croissant par paliers & exp35, 35.1 \\
\bottomrule
\end{tabular}
\end{table}




\vspace{2cm}

\subsection{Stabilisation et optimisation de l'apprentissage par renforcement en environnement multi-tour}
\label{subsec:stabilisation}

Cette sous-section établit les conditions de stabilité de l'entraînement, sur lesquelles reposent toutes les comparaisons ultérieures. Elle suit l'ordre logique du problème plutôt que la chronologie du stage.

\begin{itemize}
    \item Les conditions de possibilité du gradient lui-même : le masquage des observations et la conception de la récompense (4.2.1).
    \item La réplication de la recette du papier en \emph{full finetuning}, qui sert de référence de stabilité (4.2.2).
    \item Le passage à LoRA, ses motivations, et les instabilités qu'il introduit dans notre régime (4.2.3).
    \item Le correctif que nous proposons, l'ancre KL mobile, et le croisement d'expériences qui l'isole (4.2.4).
    \item Les signaux qui précèdent un collapse, relevés sur l'ensemble de nos runs (4.2.5).
\end{itemize}

Chacune des trois pathologies annoncées en section~2.3, détournement de récompense, dérive de politique et collapse d'entropie, y trouve une illustration empirique et une contre-mesure. C'est la partie de ce travail que nous jugeons la plus directement transférable aux futurs agents de commerce de Criteo : les difficultés rencontrées tiennent au RL multi-tour lui-même, non à TextCraft.

\subsubsection{Conditions de possibilité : le gradient et la récompense}

\paragraph{Masquage des observations.} La première condition est simple à énoncer, mais son défaut est silencieux si l'on n'y pense pas. Elle tient au calcul du gradient en multi-tour. En mono-tour, la séquence entraînée est entièrement produite par la politique. En multi-tour, elle alterne des tokens produits par la politique, les actions, et des tokens produits par l'environnement, les observations. Deux exigences en découlent. Le gradient ne doit porter que sur les actions : entraîner le modèle à prédire les observations reviendrait à lui faire apprendre le simulateur. Mais les observations doivent rester dans la séquence : chaque action a été générée en les ayant sous les yeux, et les log-probabilités recalculées lors de la rétropropagation doivent l'être dans ce même contexte, faute de quoi les ratios d'importance et la divergence KL sont calculés sur des quantités qui ne correspondent à rien.

Notre implémentation initiale ne satisfaisait que la première exigence : elle transmettait à la bibliothèque d'entraînement les actions seules, sans les observations intercalées. Le run s'exécutait, la perte décroissait, et le modèle n'apprenait pas : après quatre époques, son score restait au niveau du modèle non entraîné. Le correctif consiste à conserver la séquence complète, observations et marqueurs de dialogue compris, et à la couvrir d'un masque binaire nul sur tout ce que la politique n'a pas produit (figure~\ref{fig:env_mask}). La rétropropagation voit alors le contexte exact de la génération, et le gradient ne porte que sur les tokens de l'assistant. C'est l'analogue, à l'échelle de la trajectoire, du masquage de l'instruction en SFT (section~1.2). Après correction, le même pipeline montrait une récompense positive dès les premières mises à jour. Ce défaut ne produit aucune erreur ni aucune métrique anormale ; seule la comparaison systématique au score du modèle non entraîné l'a révélé.

\begin{figure}[htbp]
\centering
\begin{tcolorbox}[colback=gray!4, colframe=black!65, width=0.96\textwidth,
    title=\textbf{Séquence d'entraînement d'un épisode et masque des observations},
    fonttitle=\small, boxrule=0.6pt, arc=2mm, left=3mm, right=3mm, top=2mm, bottom=2mm]
{\small\ttfamily
séquence\ \ : [prompt | action$_1$ | obs$_1$ | action$_2$ | obs$_2$ | ... | action$_T$]\\
masque\ \ \ \ : [\ \ 0\ \ \ \ |\ \ 1 1 1\ \  |0 0 0 |\ \ 1 1 1\ \ |0 0 0 | ... |\ \ 1 1 1\ \ \ ]\\
gradient\  \ : [\ \ non\ \ |\ \ oui\ \ \ \  |\ non\ \ |\ \ oui\ \  \ \ |\ non\ \ | ... |\ \ oui\ \ \ \ \ ]
}
\end{tcolorbox}
\caption{Une trajectoire multi-tour telle qu'elle est transmise à l'optimiseur. Les observations de l'environnement et les marqueurs de dialogue restent dans la séquence, pour que chaque action soit évaluée dans le contexte où elle a été générée, mais leur masque est nul : ni la perte, ni la divergence KL, ni les ratios d'importance ne sont calculés sur ces tokens. Seuls les tokens produits par la politique, actions et marqueur de fin de tour, portent le gradient.}
\label{fig:env_mask}
\end{figure}





\paragraph{Conception de la récompense.} La seconde condition concerne la récompense elle-même. Le signal terminal de TextCraft est binaire et rare : sur un modèle non entraîné, neuf trajectoires sur dix rapportent zéro. Notre première campagne a donc tenté de le densifier par des termes de forme, dans l'idée d'encourager des actions bien formées avant d'exiger la réussite. Deux schémas ont été essayés. Le premier ajoutait à la récompense terminale un bonus de $+0{,}02$ pour une sortie contenant exactement une action, une pénalité de $-0{,}05$ sinon, et $-0{,}01$ par action invalide. Le second appliquait le même bonus tour par tour, lorsque le nombre moyen d'actions par tour restait proche de un. Aucun des deux n'a amélioré le modèle ; le second l'a nettement dégradé sous son niveau initial.

Le mécanisme est lisible dans les métriques d'entraînement. Un épisode qui abandonne après un tour propre encaisse le bonus. Un épisode qui tente une chaîne de plusieurs \emph{craft} s'expose à des actions invalides et à des pénalités, et échoue le plus souvent. Au sein d'un groupe GRPO où aucune trajectoire ne réussit, le terme de forme est le seul qui varie : les échecs courts et bien formés reçoivent un avantage positif, les tentatives longues un avantage négatif. La politique apprend donc à abandonner vite. Les métriques le montrent : en une vingtaine de mises à jour, la longueur moyenne des sorties tombe sous cent tokens, la récompense moyenne dépasse un alors qu'aucune tâche n'est résolue, et les groupes deviennent uniformes. C'est un détournement de récompense au sens de \cite{skalse2022defining} : la politique optimise exactement la fonction écrite, qui n'encodait pas l'intention.


Cet épisode nous a conduits à une position que la littérature récente conforte. Le rapport DeepSeek-R1 n'emploie que des récompenses par règles, exactitude et format, et écarte les modèles de récompense appris parce qu'ils « peuvent souffrir de reward hacking » à grande échelle~\cite{deepseekr1_2025}. L'analyse de Dr.~GRPO~\cite{liu2025drgrpo} porte une leçon voisine sur l'objectif lui-même : la normalisation par la longueur de la réponse, anodine en apparence, pénalise moins les réponses fausses longues que les courtes, et allonge les sorties sans les améliorer. Dans les deux cas, un terme d'apparence raisonnable induit un comportement néfaste dès que l'on suit le calcul de l'avantage. À l'inverse, des travaux appliqués empilent des termes de récompense à coefficients posés sans justification : Shop-R1~\cite{shopr1_2026} combine cinq termes pondérés de 0{,}1 à 0{,}5 et un facteur d'échelle de 1000, dont l'ablation ne teste que la présence, jamais la valeur. Notre observation ne dit pas que toute densification est nuisible. Les modèles de récompense de processus~\cite{lightman2023verify} cherchent à la faire proprement, et une récompense à seuil, nulle tant que tous les tests ne passent pas puis binaire, est courante en génie logiciel~\cite{song2026swemaster}. Elle dit qu'un terme de forme ajouté par bon sens, sans suivre sa propagation dans l'avantage relatif, est plus dangereux que la rareté qu'il prétend corriger.

Nous avons donc conservé la récompense binaire exacte pour toute la suite. C'est le cas standard des environnements agentiques et des bancs de raisonnement, et celui que les constructeurs recommandent. Les applications de recommandation, où le signal n'est ni binaire ni exact, demanderont un travail propre sur la récompense que ce stage n'a pas mené.






\subsubsection{Full finetuning : la réplication comme référence de stabilité}

La lignée \textsc{Réplication-FT} porte la recette du papier sur notre pile logicielle. Nous avons repris ses hyperparamètres dans son code publié, et non seulement dans son texte : taux d'apprentissage $10^{-6}$, coefficient KL $10^{-3}$ vers le modèle de base, $G=8$, température $1{,}0$, collectes de 32 tâches soit 256 trajectoires, consommées en quatre pas d'optimisation de 64 trajectoires avec clipping (tableau~\ref{tab:hyperparams}). Le texte et le code divergent cependant sur un point. L'annexe du papier fixe « le nombre maximal d'interactions à 20 tours » ; les scripts publiés, d'entraînement comme d'évaluation, en fixent 30. Nous avons suivi le code, à l'entraînement et à l'évaluation. Plus largement, les conditions des expériences sur le modèle 3B ne sont pas documentées : le script d'entraînement publié est configuré pour le 7B, aucune configuration 3B n'est fournie, et le papier ne rapporte ni nombre de mises à jour ni nombre d'époques pour aucun modèle. Son script en fixe 30, sans que l'on sache si les résultats publiés en proviennent.

Cette lignée établit trois faits. Le premier concerne le budget. Après 30 époques, notre modèle atteint 40\,\%, contre 75\,\% rapportés par le papier pour le même modèle. En prolongeant la recette, la performance croît de façon quasi monotone : 58\,\% après une centaine d'époques, 64\,\% puis 69\,\% après environ 247 époques cumulées. Ces reprises successives depuis le meilleur point de contrôle étaient imposées par les purges hebdomadaires de notre environnement de calcul (annexe~C). Aucun run \emph{full finetuning} n'a dépassé ce 69\,\%. L'écart avec le 75\,\% publié n'a pas d'explication établie, la durée et les conditions exactes de leurs runs n'étant pas documentées. Une hypothèse est compatible avec l'écart texte-code relevé ci-dessus. Si leurs résultats proviennent bien d'un horizon de 20 tours, l'espace d'exploration est plus restreint que le nôtre : les trajectoires longues et improductives sont coupées, et l'apprentissage se concentre sur les tâches résolubles en peu de tours. Or ces tâches, de profondeur 1 et 2, forment 72 des 100 tâches de test. Un horizon court favoriserait ainsi le score global, au prix des tâches profondes. Nos propres résultats sur le curriculum d'horizon (section~4.3) vont dans ce sens, sans permettre de trancher. Nous avons néanmoins conservé 30 tours pour toutes nos recettes : c'est le dernier palier du curriculum d'horizon, et une référence entraînée au même horizon final est la seule qui permette de mesurer ce que le curriculum apporte, et non ce qu'un horizon plus court apporte à lui seul.



Le deuxième fait est la stabilité (figure~\ref{fig:fullft_stabilite}). À aucun moment la lignée n'a présenté de collapse. La divergence KL vers la référence fixe croît vite pendant les trente premières époques, de $10^{-3}$ à 0{,}024, puis de plus en plus lentement, pour atteindre 0{,}08 après 247 époques. L'entropie de la politique monte d'abord à 1{,}2 vers la quinzième époque, puis décroît régulièrement jusqu'à 0{,}25, sans à-coup : la politique s'affûte sans s'effondrer. La récompense d'entraînement et le score de test suivent la même montée, avec la dispersion la plus faible de toutes nos lignées. Sur les treize dernières évaluations, le score de test vaut 61\,\% en moyenne, pour un meilleur à 69\,\%. Cette stabilité est obtenue avec $\beta = 10^{-3}$ et sans aucun des mécanismes développés plus loin : c'est la référence à laquelle la suite se compare.



Le troisième fait est le coût. La progression est lente : le dernier segment de 72 époques, de 63 à 69\,\%, a demandé une soixantaine d'heures de GPU. L'état de l'optimiseur pèse 8~Go par point de contrôle, ce qui limite ce que l'on peut conserver et reprendre sur un disque de 35~Go. C'est ce coût, en temps et en stockage, qui motive le passage à LoRA de la sous-section suivante.

\begin{table}[htbp]
\centering
\caption{Les trois recettes d'entraînement de la section 4. Toutes utilisent GRPO avec récompense binaire, masquage des observations et 30 tours d'interaction au plus. « Collecte » : nombre de trajectoires générées avant optimisation, consommées en pas de 64. La référence KL fixe est le modèle de base ; la référence mobile est déplacée toutes les 4 époques. Tous les adaptateurs LoRA ont $\alpha = 32$.}
\label{tab:hyperparams}
\footnotesize
\setlength{\tabcolsep}{4pt}
\begin{tabular}{l l c c c c c}
\toprule
\textbf{Recette} & \textbf{Poids entraînés} & \textbf{Taux d'appr.} & $\boldsymbol{\beta}$ & \textbf{Réf. KL} & $\boldsymbol{G}$ & \textbf{Collecte} \\
\midrule
\textsc{Réplication-FT} & tous & $10^{-6}$ & $10^{-3}$ & fixe & 8 & 256 \\
\textsc{Grille-LoRA} & LoRA $r$ = 16, 32, 64 & $10^{-6}$ à $10^{-5}$ & $10^{-3}$, $10^{-2}$ & fixe & 8 & 256 \\
\textsc{Ancre-Mobile} & LoRA $r$ = 8 & $3 \times 10^{-6}$ & $10^{-2}$ & mobile & 8 puis 16 & 64 \\
\bottomrule
\end{tabular}
\end{table}






\begin{figure}[htbp]
    \centering
    \includegraphics[width=\textwidth]{figures/fig_fullft_stabilite.pdf}
    \caption{Lignée \textsc{Réplication-FT} sur 247 époques cumulées, quatre segments séparés par les traits verticaux (chaque segment reprend le meilleur ou le dernier point de contrôle du précédent). (a) Récompense d'entraînement moyennée sur une époque et pass@$1$ test à chaque évaluation, sur la même échelle. (b) Divergence KL vers la référence fixe, médiane sur une époque, échelle logarithmique. (c) Entropie de la politique, moyenne sur une époque. Montée quasi monotone, KL bornée, entropie en décroissance lente : aucun collapse sur la durée.}
    \label{fig:fullft_stabilite}
\end{figure}







\subsubsection{LoRA : motivation, grille et diagnostic}

\paragraph{Pourquoi LoRA.} La lignée précédente a fixé la référence, mais à un coût qui interdisait la suite du programme : à raison d'une soixantaine d'heures pour 72 époques, la campagne de curriculums prévue n'aurait pas tenu dans le stage. L'adaptation de rang faible~\cite{hu2022lora} n'entraîne que de petites matrices additives, ce qui divise la mémoire de l'optimiseur et la taille des points de contrôle par un facteur proche de cent : 120~Mo d'état optimiseur et 180~Mo par meilleur modèle en rang 8, contre 8~Go et 14~Go en \emph{full finetuning}. Cette économie a trois usages dans la suite. Elle rend les reprises après purge presque gratuites. Elle libère de la mémoire GPU pour le moteur de génération, ce qui permettra de doubler la taille des groupes GRPO (section~4.3). Et elle autorise des grilles d'expériences que le \emph{full finetuning} rendait impossibles sur un GPU.

Le blog \emph{LoRA Without Regret}~\cite{schulman2025lora} a apporté une seconde motivation, d'ordre scientifique. Ses auteurs observent qu'en RL, LoRA égale le \emph{full finetuning} même à très petit rang, et proposent une explication : un algorithme de gradient de politique n'apprendrait « qu'environ un bit d'information par épisode », puisqu'un seul scalaire de récompense conclut chaque épisode, si bien qu'un sous-espace de paramètres restreint suffirait à l'absorber. Ils présentent cette lecture comme une hypothèse, et l'appuient d'expériences de raisonnement mathématique mono-tour à récompense binaire. Ils en tirent une recette : un taux d'apprentissage environ dix fois supérieur à celui du \emph{full finetuning}, à peu près indépendant du rang. Nous avons voulu savoir si cette recette et cette explication tenaient dans notre régime, qui en diffère par le multi-tour, ses trajectoires longues et ses observations intercalées.

\paragraph{La grille.} \textsc{Grille-LoRA} reprend la recette du papier et n'en change que l'adaptation : trois rangs ($r \in \{16, 32, 64\}$, $\alpha = 32$), deux taux d'apprentissage ($10^{-5}$, la prédiction du blog, et $3 \times 10^{-6}$) et deux valeurs de $\beta$ ($10^{-3}$ et $10^{-2}$), toujours vers une référence fixe. Six des neuf configurations ont été courues ; les trois autres, prévisibles au vu des premières, ont été abandonnées. Le paysage n'offre aucune zone de fonctionnement satisfaisante.

\begin{itemize}
    \item Le taux d'apprentissage du blog provoque un collapse rapide et complet : pic à 24\,\% vers la troisième époque, puis 0\,\%.
    \item À $3 \times 10^{-6}$ et $\beta = 10^{-3}$, la montée est la plus rapide, 31 à 35\,\% selon le rang, au-dessus de la référence \emph{full finetuning} à budget égal (30 à 32\,\% à l'époque 15), puis chaque run collapse en seconde moitié.
    \item À $3 \times 10^{-6}$ et $\beta = 10^{-2}$, $r=64$ plafonne entre 21 et 29\,\% sans s'effondrer ; $r=16$ atteint 39\,\% puis chute à 10\,\% vers la douzième époque avant de remonter partiellement.
    \item Le rang ne sauve rien, et la relation va à rebours de l'intuition de capacité : à hyperparamètres égaux, $r=16$ fait mieux que $r=32$, qui fait mieux que $r=64$.
\end{itemize}

\paragraph{Diagnostic.} La divergence KL vers la référence rend ce paysage lisible (figure~\ref{fig:kl_lora_vs_fullft}). Tous les runs démarrent au même niveau, $10^{-3}$, la dérive d'un pas d'optimisation. Le \emph{full finetuning} s'en éloigne ensuite régulièrement, d'environ $10^{-3}$ par époque sur la fenêtre de la figure, puis dix fois plus lentement au-delà de la trentième époque (figure~\ref{fig:fullft_stabilite}). Les runs LoRA à référence fixe s'en éloignent géométriquement : à partir de la quatrième époque, leur KL est multipliée par cinq à dix toutes les deux époques, d'autant plus vite que le rang est grand, d'autant plus tard que $\beta$ est fort. Tous franchissent un même niveau, entre 0{,}05 et 0{,}1, et dépassent 1 dans les quatre époques qui suivent ; aucun ne s'est rétabli. Ce niveau n'a rien d'absolu : le \emph{full finetuning} s'y installe après trente époques sans dommage, et les runs à référence mobile de la suite le dépassent parfois transitoirement avant d'y être ramenés. Ce qui le rend fatal ici est la dynamique qui l'atteint, une croissance géométrique que rien ne freine. 





Un ordre de grandeur explique pourquoi $\beta$ ne peut pas l'empêcher. Le terme de politique de la perte vaut environ $3 \times 10^{-3}$ sur un run sain. À KL $= 10^{-3}$ et $\beta = 10^{-2}$, la pénalité vaut $10^{-5}$, trois cents fois moins : elle est inerte en régime sain, ce qui est souhaitable. Elle ne devient comparable au terme de politique que vers KL $= 0{,}3$, quand la croissance géométrique est déjà engagée. Le frein n'engage donc qu'une fois la dérive irréversible, et une valeur de $\beta$ dix fois plus forte plafonnerait des scores que $\beta = 10^{-2}$ plafonne déjà.

Reste à expliquer pourquoi LoRA dérive plus vite que le \emph{full finetuning} à progression égale. Nous n'avons pas de réponse établie, et proposons deux hypothèses.

\begin{itemize}
    \item \textbf{Paramétrisation en produit.} L'incrément de poids s'écrit $BA$, et le gradient de chaque facteur est proportionnel à l'autre. Les auteurs du blog invoquent cette dynamique d'optimisation « moins favorable que celle de la matrice complète » pour expliquer que LoRA tolère mal les grandes collectes en apprentissage supervisé. Notre collecte de 256 trajectoires et nos quatre pas par collecte placent la grille dans ce régime ; l'hypothèse est que le même couplage amplifie chaque pas en RL.
    \item \textbf{Bruit du signal.} Une récompense binaire et rare produit des gradients bruités. Le \emph{full finetuning} les répartit sur des milliards de paramètres, où ils se compensent en partie ; LoRA les concentre sur quelques matrices, où ils s'accumulent. C'est notre hypothèse, non celle du blog, et elle n'a pas été testée isolément.
\end{itemize}

Quelle qu'en soit la cause, la conséquence est la même : à progression égale, la politique LoRA s'éloigne plus vite de sa référence, et le rappel vers une référence fixe est soit trop faible pour empêcher l'emballement, soit assez fort pour l'empêcher mais au prix d'un plafond. Ce diagnostic désigne le correctif de la sous-section suivante. Précisons enfin que le rang $r=8$ retenu ensuite prolonge la tendance de la grille sans avoir été comparé aux autres rangs sous ancre mobile ; nous le classons parmi les variables non contrôlées de l'étude.

\begin{figure}[htbp]
    \centering
    \includegraphics[width=0.9\textwidth]{figures/fig_kl_lora_vs_fullft.pdf}
    \caption{Divergence KL vers la référence fixe pendant les douze premières époques, médiane glissante sur une époque, échelle logarithmique. Le \emph{full finetuning} dérive linéairement ; tous les runs LoRA dérivent géométriquement, d'autant plus vite que le rang est grand (couleur) et d'autant plus tard que $\beta$ est fort (trait plein $10^{-2}$, tirets $10^{-3}$). La bande grise marque le niveau franchi par tous les runs LoRA à référence fixe une à quatre époques avant leur divergence ; le \emph{full finetuning} l'atteint après trente époques sans dommage.}
    \label{fig:kl_lora_vs_fullft}
\end{figure}









\subsubsection{L'ancre KL mobile : correctif et croisement d'expériences}

\paragraph{Principe.} Le diagnostic précédent enferme la référence fixe dans un dilemme : un rappel faible laisse la dérive s'emballer, un rappel fort l'empêche mais plafonne les scores. Dans les deux cas, le défaut est le même : la pénalité KL mesure la distance à un point de plus en plus lointain, le modèle de base, et tire la politique vers l'arrière, contre son progrès. La sortie du dilemme consiste à déplacer la référence. Notre ancre KL mobile procède par fusion et redémarrage périodiques : toutes les quatre époques, l'adaptateur courant est fusionné dans les poids du modèle, un adaptateur vierge le remplace, la pénalité KL se mesure désormais vers ce nouveau point, et l'état de l'optimiseur des paramètres LoRA est remis à zéro. La politique est inchangée au moment du ré-ancrage ; seul le point de référence de la pénalité bouge. La nature du rappel change alors : au lieu de retenir la politique près de son point de départ, il la retient près de son état d'il y a quelques époques, et ne pénalise que la dérive rapide.

Ce mécanisme rejoint une idée ancienne du RL profond. TRPO~\cite{schulman2015trpo} contraignait chaque mise à jour à rester, en divergence KL, dans une région de confiance autour de la politique précédente ; PPO~\cite{schulman2017proximal} a remplacé cette contrainte par le clipping du ratio d'importance (annexe~A). Une pénalité KL vers une référence récente est une forme souple de la même idée, à l'échelle de quelques époques plutôt que d'un pas. Nous l'avons conçue à partir de nos propres courbes, puis retrouvée dans la littérature : ProRL~\cite{liu2025prorl} remet périodiquement sa politique de référence « sur un instantané récent de la politique en ligne » et réinitialise l'optimiseur, parce que « le terme KL finit par dominer la perte et bloquer les mises à jour » ; TR-DPO~\cite{gorbatovski2024trdpo} met à jour la référence en optimisation de préférences ; une ancre en moyenne mobile a été proposée depuis~\cite{zhang2026emapg}. Chez ProRL le déclencheur est la stagnation de la validation ; chez nous la période est fixe.

\paragraph{Résultats.} La configuration \textsc{Ancre-Mobile} ($r=8$, $\beta=10^{-2}$, ré-ancrage toutes les 4 époques) monte de 12 à 35\,\% en 15 époques sans accident, là où les meilleures configurations à référence fixe ont déjà atteint leur pic et commencé à redescendre (figure~\ref{fig:anchor_square}a). Poursuivie, elle atteint 65\,\% vers la centième époque après 25 ré-ancrages ; reprise depuis ce point pour soixante époques supplémentaires, elle atteint 73\,\%. Aucun collapse sur l'ensemble de la lignée. Le panneau (b) de la figure en donne la raison : la KL vers la référence reste à $10^{-3}$ pendant toute la durée, avec une dent de scie à chaque ré-ancrage, quand celle des runs à référence fixe s'emballe entre la cinquième et la huitième époque.

Le croisement \textsc{Carré-Ancre} autorise deux conclusions. D'abord, l'ancre mobile est nécessaire dans notre régime : à référence fixe, la KL diverge quelle que soit la valeur de $\beta$, entre la cinquième et la huitième époque, et aucun run n'a plus progressé ensuite ; le score s'effondre, ou se fige lorsque la politique devenue déterministe continue de résoudre les mêmes tâches (figure~\ref{fig:anchor_square}a et~b). Ensuite, le croisement lève un facteur confondu de la grille : $\beta = 10^{-3}$, dont la KL divergeait à référence fixe, reste à $10^{-3}$ de KL à référence mobile, et le run est stable à 53\,\% après 60 époques. C'est donc le régime de référence, et non la valeur de $\beta$, qui gouverne la stabilité. Les deux valeurs de $\beta$ sont d'ailleurs indiscernables à référence mobile : entre les époques 49 et 58, les évaluations valent en moyenne 48{,}9\,\% pour $\beta=10^{-2}$ et 49{,}7\,\% pour $\beta=10^{-3}$. L'écart final, 65 puis 73 contre 53, tient à la durée des runs, le second ayant été arrêté à l'époque 60. Nous avons conservé $\beta=10^{-2}$ pour la suite par continuité de lignée, non parce qu'il serait meilleur.

\paragraph{Portée.} Cette recette a porté notre modèle au-dessus du score publié, mais nous ne prétendons ni qu'elle soit universelle ni qu'elle soit la meilleure. Elle a été mise au point sur un environnement, avec un rang et une période de ré-ancrage choisis sans grille. Sa robustesse reste à établir sur les autres environnements d'AgentGym-RL, dont la plateforme fournit l'infrastructure. Elle a en revanche un bénéfice immédiat déjà noté : la mémoire libérée par LoRA a permis, dans la phase suivante, de doubler la taille des groupes GRPO à nombre de trajectoires constant, donc d'estimer l'avantage de chaque tâche sur 16 trajectoires au lieu de 8.



\begin{figure}[htbp]
    \centering
    \includegraphics[width=\textwidth]{figures/fig_anchor_square.pdf}
    \caption{Croisement référence fixe ou mobile $\times$ $\beta$ (couleur : $\beta$ ; tirets : référence fixe ; trait plein : référence mobile). (a) Pass@$1$ test sur les seize premières époques : les runs à référence fixe atteignent leur pic entre la sixième et la onzième époque puis s'effondrent, d'autant plus vite que $\beta$ est faible ; les runs à référence mobile montent sans accident. (b) Divergence KL vers la référence, médiane sur une époque, échelle logarithmique : les trois runs à référence fixe franchissent la bande grise entre la cinquième et la huitième époque puis divergent ; les runs à référence mobile restent à $10^{-3}$, avec une dent de scie à chaque ré-ancrage. (c) Pass@$1$ test des deux runs à référence mobile sur toute leur durée : 65\,\% vers l'époque 100 pour $\beta=10^{-2}$, 53\,\% à l'arrêt volontaire de l'époque 60 pour $\beta=10^{-3}$.}
    \label{fig:anchor_square}
\end{figure}






\subsubsection{Signaux précurseurs et bilan}

\paragraph{Deux modes de collapse.} Nous disposons de quinze runs de plus de cinq époques sur la pile actuelle : six stables sur la durée, six collapses à référence fixe, et trois collapses à $G=16$ sans curriculum, détaillés en section~4.3. Leurs hyperparamètres diffèrent, ce qui interdit toute lecture causale ; mais les trajectoires d'entropie et de KL se regroupent en deux signatures nettes (figure~\ref{fig:collapse_signals}). Nous rappelons d'abord le résultat empirique de Cui et al.~\cite{cui2025entropy} qui sert de grille de lecture : en RL à récompense vérifiable, l'entropie de la politique décroît à mesure que la performance monte, la performance « se paie en entropie », et lorsque l'entropie approche de zéro la politique devient déterministe, l'exploration s'éteint et l'apprentissage s'arrête.

\begin{itemize}
    \item \textbf{À référence fixe, la KL précède.} Dès la quatrième époque, la KL médiane croît géométriquement (section~4.2.3) alors que l'entropie reste élevée, et même monte jusqu'à 1{,}4 ou 1{,}5 : la politique devient plus erratique, non plus déterministe. Entre la cinquième et la onzième époque, la KL dépasse 1 et l'entropie s'effondre au même moment, en moins d'une époque. La politique est alors soit détruite, soit figée.
    \item \textbf{À $G=16$ sans curriculum, l'entropie précède.} La KL reste à $10^{-3}$, ramenée à chaque ré-ancrage, tandis que l'entropie décroît régulièrement, de 1{,}1 vers 0{,}6 en six à huit époques. Puis, en moins d'une époque, l'entropie tombe sous 0{,}2 et la KL saute de $10^{-3}$ à plus de $10^{8}$. C'est le collapse d'entropie de Cui et al., que l'ancre mobile ne peut pas voir venir : en se re-basant sur une politique en cours de déterminisation, elle valide la dégénérescence au lieu de la freiner. La section~4.3 en propose le mécanisme.
\end{itemize}

Un troisième cas éclaire les deux premiers. Le run \textsc{Contrôle-G16} relancé avec un ré-ancrage toutes les douze époques suit la première signature, non la seconde : sans ré-ancrage avant l'époque 12, il se comporte comme une référence fixe, et sa KL diverge dès la cinquième époque. La figure~\ref{fig:collapse_signals}b montre d'ailleurs que tous les runs à $G=16$, stables ou non, dérivent vite pendant leurs quatre premières époques, jusqu'à 0{,}1 ou 0{,}3 de KL : c'est le ré-ancrage de la quatrième époque qui les ramène à $10^{-3}$. Espacer les ré-ancrages n'est donc pas une option dans ce régime.

\paragraph{Limites de l'analyse des signaux.} Deux mesures accompagnent l'effondrement sans le précéder. La longueur des générations s'écrase au moment du collapse, sous 250 tokens de sorties stéréotypées, puis explose dans certains runs, jusqu'à plusieurs milliers de tokens de texte dégénéré. La récompense d'entraînement, elle, ne prévient de rien : sur \textsc{Contrôle-G16}, elle se maintient à 0{,}5 pendant les deux époques où la KL vaut déjà $10^{8}$, parce qu'une politique déterministe continue de résoudre les tâches qu'elle connaît. Les deux signaux précurseurs exploitables sont donc la croissance géométrique de la KL à référence fixe, et la décroissance régulière de l'entropie à référence mobile. Nous n'avons pas mené d'analyse de corrélation systématique entre ces signaux, les classes d'erreurs et le moment du collapse ; c'est une analyse directement réalisable sur nos journaux, que nous laissons en travail futur.

\paragraph{Régulation du taux d'apprentissage.} Avant l'ancre mobile, nous avons cherché la stabilité dans la régulation du taux d'apprentissage, sur une lignée de runs LoRA à référence fixe. Des paliers calendaires stabilisent mais interrompent des progressions en cours. Une coupe adaptative, déclenchée sur la stagnation de la récompense d'entraînement, se déclenche souvent sur du bruit ; surtout, couper le taux consolide l'état courant, qui après un pic est souvent déjà une dérive. La variante la plus saine combine des paliers longs et la restauration du meilleur état du palier écoulé. Ces mécanismes et la politique de conservation des meilleurs points de contrôle avec leur optimiseur sont décrits en annexe~C. La recette finale n'en a plus besoin : à taux constant, l'ancre mobile suffit.

\paragraph{Bilan de la section.} Cette section a établi les conditions dans lesquelles un entraînement GRPO multi-tour progresse sans s'effondrer, et les a vérifiées sur nos runs. Le gradient n'est correct qu'à la condition de garder les observations dans la séquence et de les masquer ; sans cela, le modèle n'apprend pas et rien ne le signale. Un terme de récompense ajouté pour densifier le signal a été optimisé pour lui-même : la politique a appris à abandonner les tâches longues, et nous avons conservé la récompense binaire. La recette du papier en \emph{full finetuning} est stable sur 247 époques et atteint 69\,\%, à un coût en temps et en mémoire qui interdisait la suite. LoRA supprime ce coût mais dérive géométriquement en KL vers une référence fixe, quels que soient le rang et le coefficient testés. Déplacer la référence toutes les quatre époques supprime la dérive et porte le même modèle à 73\,\%. Enfin, l'analyse de quinze runs distingue deux modes de collapse, l'un par la KL à référence fixe, l'autre par l'entropie à grande taille de groupe. L'ancre mobile règle le premier. Le second est l'objet de la section suivante.

\begin{figure}[htbp]
    \centering
    \includegraphics[width=\textwidth]{figures/fig_collapse_signals.pdf}
    \caption{Seize premières époques de quinze runs, groupés par issue : stables sur la durée (bleu), collapse à référence fixe (orange, tirets), collapse à $G=16$ sans curriculum (vert). (a) Entropie de la politique, moyenne sur une époque : à référence fixe, elle monte jusqu'à 1{,}4 ou 1{,}5 avant de s'effondrer d'un coup ; à $G=16$, elle décroît régulièrement pendant six à huit époques avant de s'effondrer. (b) Divergence KL vers la référence, médiane sur une époque, échelle logarithmique : croissance géométrique à référence fixe ; à référence mobile, dérive rapide pendant les quatre premières époques puis retour à $10^{-3}$ à chaque ré-ancrage (dents de scie), sauf pour le run ré-ancré toutes les douze époques, qui diverge dès la cinquième.}
    \label{fig:collapse_signals}
\end{figure}


















\vspace{2cm}


\subsection{Expériences RL : les curriculums}
\label{sec:curriculums}

La section~4.2 a laissé une question ouverte. L'ancre KL mobile supprime la dérive de la politique loin de sa référence, mais elle ne voit pas le second mode de collapse : celui où l'entropie s'éteint progressivement, à KL stable, jusqu'à ce que la politique devienne déterministe. Le curriculum est le premier candidat pour y répondre, et c'est à ce titre que nous l'avons d'abord étudié.

Il se prête cependant à une double lecture, et cette dualité est l'idée centrale de la section. D'un côté, le curriculum est un outil d'optimisation : une manière de doser la difficulté présentée à la politique pour qu'elle apprenne plus vite, plus haut, ou plus sûrement. De l'autre, il touche à des questions de recherche plus ouvertes, celles de l'exploration et du méta-apprentissage : la structure de l'expérience peut-elle faire dépasser à une politique ce que son modèle de base savait déjà faire ? Ce qui est acquis sur une difficulté se transfère-t-il aux autres ? L'ordre de présentation change-t-il la manière d'apprendre, la vitesse d'acquisition de chaque profondeur, la généralisation, l'ordre dans lequel les erreurs disparaissent ? Ces questions sont plus risquées que la première, et nous ne prétendons pas les trancher avec trois runs. Mais elles justifient l'intérêt porté au curriculum au-delà de son rôle de stabilisateur : au pire, il stabilise ; au mieux, il apprend autrement.

La section se déroule en trois temps. La sous-section~4.3.1 décrit trois régimes de curriculum comparés à recette d'entraînement identique, face à deux références sans curriculum. La sous-section~4.3.2 présente leurs résultats et la lecture qu'ils autorisent, dont le rôle du curriculum dans la stabilité. La sous-section~4.3.3 propose un mécanisme, au niveau du gradient, qui rend compte à la fois du collapse sans curriculum et de la protection qu'il apporte. La comparaison des régimes à compute contrôlé fait l'objet de la section~4.4. Le transfert entre planification mono-tour et contrôle interactif, initialement prévu ici, n'a pu être mené au-delà de mesures préliminaires ; il est présenté en section~4.5 et repris en ouverture comme une forme de méta-apprentissage.





\subsubsection{Trois régimes de curriculum à recette commune}

\paragraph{Ce qu'est un curriculum.} Dans sa formulation initiale, le \emph{curriculum learning} présente les exemples d'entraînement par difficulté croissante~\cite{bengio2009curriculum}. En apprentissage par renforcement, la notion s'est élargie : la revue de Narvekar et al.~\cite{narvekar2020curriculum} définit le curriculum comme une séquence de tâches, ou d'échantillons, organisée pour apprendre un problème cible, et celle de Portelas et al.~\cite{portelas2020acl} décrit les méthodes de curriculum automatique comme des mécanismes « qui façonnent la trajectoire d'apprentissage d'un agent en lui proposant des tâches adaptées à ses capacités ». Deux leviers sont donc possibles. Le premier ordonne les tâches : quelles tâches présenter, et quand. Le second agit sur les moyens accordés à l'agent pour résoudre une tâche donnée : combien de tours, combien de tokens. ScalingInter~\cite{xi_agentgym-rl_2025} relève de ce second levier, encore peu exploré. Les régimes ci-dessous en couvrent les deux.

\paragraph{Recette commune et références.} Les trois régimes partagent la recette stabilisée en section~4.2 : LoRA de rang 8, taux d'apprentissage $3 \times 10^{-6}$, $\beta = 10^{-2}$, ancre mobile toutes les 4 époques, groupes de $G=16$ trajectoires et 64 trajectoires par mise à jour, horizon final de 30 tours, 80 époques prévues. Ils se comparent à deux références sans curriculum. \textsc{Ancre-Mobile} est la référence à $G=8$, 65 puis 73\,\%. \textsc{Contrôle-G16} reprend la recette commune à l'identique, sans curriculum : c'est le contrôle qui isole l'effet du curriculum de celui de la taille de groupe. Il n'a pas tenu, et ce collapse est l'un des résultats de la section suivante. Conformément au protocole, chaque configuration est un run unique et aucun hyperparamètre n'a été réglé séparément pour un régime : les conclusions sont conditionnelles à cette recette commune.

\begin{itemize}
    \item \textbf{\textsc{Horizon} : le nombre de tours.} Ce régime reprend le principe de ScalingInter : le nombre maximal de tours d'interaction par épisode vaut 10 de l'époque 0 à 15, 20 de l'époque 15 à 30, puis 30. L'hypothèse est la suivante. Un plafond bas ne laisse réussir que les tâches résolubles en peu de tours : les plus faciles, et les tâches moyennes à condition de les résoudre sans détour. Il force donc d'abord l'efficacité sur les tâches courtes, et rend les tâches longues temporairement insolubles ; chaque relèvement du plafond les réintroduit. Les courbes de réussite par profondeur (section~4.3.2) permettront de vérifier cette lecture. Notre calendrier est plus lent que celui du papier, dont les scripts changent de palier toutes les 100 mises à jour, soit environ toutes les deux époques dans leur configuration ; nous avons préféré des paliers longs, pour laisser chaque niveau saturer avant de passer au suivant.
    \item \textbf{\textsc{Budget} : le nombre de tokens par tour.} Ce régime transpose la même idée à la longueur des réponses de l'agent : 256 tokens par tour de l'époque 0 à 15, 512 de l'époque 15 à 35, puis 1\,024, l'horizon restant fixé à 30 tours. Un budget court interdit les raisonnements longs et les sorties dégénérées ; son relèvement autorise ensuite des tours plus élaborés. C'est une extension directe de ScalingInter à un second paramètre du harnais.
    \item \textbf{\textsc{Profondeur} : la structure des tâches.} Ce régime est le plus proche de la définition d'origine : il ordonne les tâches sans toucher aux moyens de l'agent. L'échantillonnage est uniforme sur les recettes de profondeur inférieure ou égale au palier courant, et le palier monte de 1 à 4. Le mécanisme est un échantillonneur pondéré sur l'ensemble des recettes, non un découpage en sous-ensembles : la définition des époques et le reste de la boucle d'entraînement sont identiques aux autres régimes. Une première version fixait les paliers au calendrier, aux époques 0, 12, 25 et 37. Elle a montré qu'un palier maintenu trop longtemps sature : sur le palier 1, la récompense d'entraînement atteint 1 dès la sixième époque, tous les groupes deviennent uniformes et le gradient est nul pendant six époques. La seconde version rend le passage adaptatif : le palier monte dès que la récompense d'entraînement moyenne dépasse 0{,}8 sur une époque complète, ou au plus tard après 10 époques.
\end{itemize}

Ces trois régimes fixent la difficulté à la main, selon une notion de difficulté qui est la nôtre : le nombre de tours, la longueur des réponses, la profondeur des recettes. Rien ne garantit qu'elle coïncide avec la difficulté effective pour le modèle, et dans un environnement réel, de recommandation par exemple, cette difficulté n'est même pas connue à l'avance. Un curriculum automatique, qui sélectionne les tâches d'après le progrès que le modèle y réalise, répond à cette limite ; nous en avons porté et couru une instance, présentée en section~4.5 avec les pistes exploratoires. Enfin, ces régimes agissent chacun sur un seul levier ; leur combinaison n'a pas été testée, et figure parmi les prolongements les plus directs de ce travail (section~5).






\subsubsection{Résultats : le curriculum comme condition de stabilité à $G=16$}

\paragraph{Le contrôle d'abord.} \textsc{Contrôle-G16} applique la recette commune sans curriculum. Il monte vite, de 12 à 54\,\% en treize époques, plus vite que la référence à $G=8$ à nombre égal de trajectoires (figure~\ref{fig:curricula_main}a). Puis il s'effondre selon la seconde signature de la section~4.2.5 : l'entropie de la politique décroît régulièrement de 1{,}0 à 0{,}6 entre la troisième et la dixième époque, la divergence KL restant à $10^{-3}$, puis les deux basculent en moins d'une époque, l'entropie sous 0{,}15 et la KL au-delà de $10^{8}$. Le score de test retombe de 54 à 13\,\% et les générations dégénèrent. Un second essai avec un ré-ancrage toutes les douze époques au lieu de quatre a divergé plus tôt encore, dès la cinquième époque, faute de ré-ancrage avant la douzième. Nous n'avons pas poursuivi : ce collapse est en lui-même le résultat, celui auquel les trois curriculums se comparent.

\paragraph{Les trois curriculums.} À recette identique, les trois régimes traversent la zone où le contrôle s'effondre et poursuivent leur montée sur 80 époques (figure~\ref{fig:curricula_main}a). \textsc{Horizon} atteint 82\,\% à deux reprises, aux époques 65 et 77, et termine sur un plateau dont les dix dernières évaluations valent 78\,\% en moyenne, entre 75 et 82 ; il traverse vingt ré-ancrages. \textsc{Budget} atteint 78\,\% à l'époque 55, puis 80\,\% après une reprise, sur un plateau à 75\,\%. \textsc{Profondeur}, dans sa version adaptative, atteint 72\,\% à l'époque 43, puis 80\,\% après une reprise, sur un plateau à 74\,\%. Les deux reprises ont été imposées par des purges de l'environnement de calcul et sont décrites en annexe~C. Les trois courbes sont presque superposées sur toute la durée ; \textsc{Horizon} se détache légèrement en fin de run. Elles dominent la référence \textsc{Ancre-Mobile} à $G=8$, qui atteint 73\,\% au terme de 170 époques, soit 70\,\% de plateau, à nombre égal de mises à jour. Rappelons l'échelle : la lignée \emph{full finetuning} avait demandé environ 247 époques pour atteindre 69\,\%.

\paragraph{Lecture.} Ces runs livrent la lecture centrale de la section. À recette commune, doubler la taille des groupes accélère l'apprentissage, et avec lui la déterminisation de la politique : sans curriculum, la recette est instable à cette vitesse. Avec un curriculum, quel qu'en soit le levier, elle ne le devient pas. Le curriculum n'est donc pas seulement un accélérateur, il est, dans ce régime, une condition de stabilité. La figure~\ref{fig:curricula_main}b précise le rôle qu'il joue. Il n'empêche pas l'entropie de décroître : celle d'\textsc{Horizon} passe de 0{,}8 à 0{,}2 en 80 époques, celles de \textsc{Budget} et \textsc{Profondeur} suivent la même pente. Il empêche son effondrement brutal, celui qui, sur le contrôle, précède la divergence de la KL. La section~4.3.3 propose le mécanisme de cette protection. Ce constat reste conditionnel : un run par configuration, aucun réglage propre à un régime, et deux essais seulement pour le contrôle.

\paragraph{Ordre d'acquisition et élimination des erreurs.} Deux analyses tirées des évaluations périodiques complètent le tableau. La première stratifie le pass@$1$ par profondeur (figure~\ref{fig:acquisition_depths}). L'ordre d'acquisition est le même dans tous les régimes, curriculum ou non : la profondeur 1 est acquise en moins de cinq époques, la profondeur 2 dépasse 80\,\% vers la trentième et atteint 90\,\% vers la soixantième, la profondeur 3 progresse lentement jusqu'à 30 ou 45\,\%, et la profondeur 4 reste à zéro. Les curriculums n'ont pas modifié cet ordre ; ils ont accéléré chaque étape par rapport à la référence à $G=8$. Le curriculum de profondeur, qui présente explicitement les recettes profondes en dernier, ne fait pas mieux que les autres sur la profondeur 3. La seconde analyse suit les classes d'erreurs renvoyées par l'environnement (figure~\ref{fig:elimination_erreurs}). Les erreurs de format, action mal formée ou plusieurs actions dans un tour, disparaissent en moins de dix époques dans tous les régimes ; ce sont elles qui réapparaissent en masse lors du collapse du contrôle. Les erreurs de planification, recette invalide, ingrédient manquant ou objet introuvable, décroissent lentement sur toute la durée, de douze à cinq par épisode pour \textsc{Horizon} et \textsc{Profondeur}, de douze à sept pour \textsc{Budget} et la référence. Le RL corrige donc d'abord la forme, puis, beaucoup plus lentement, le fond ; le curriculum accélère les deux sans en changer l'ordre.

\begin{figure}[htbp]
    \centering
    \includegraphics[width=\textwidth]{figures/fig_curricula_pass1_entropy.pdf}
    \caption{Trois curriculums (\textsc{Horizon}, \textsc{Profondeur}, \textsc{Budget}), le contrôle sans curriculum à $G=16$ et la référence à $G=8$, en fonction du nombre de mises à jour, chacune de 64 trajectoires. (a) Pass@$1$ test, lissé sur cinq évaluations. (b) Entropie de la politique, moyenne sur une époque. Les lignées \textsc{Profondeur}, \textsc{Budget} et $G=8$ comprennent une reprise après purge (annexe~C). Le contrôle s'effondre vers la millième mise à jour ; les trois curriculums poursuivent leur montée.}
    \label{fig:curricula_main}
\end{figure}

\begin{figure}[htbp]
    \centering
    \includegraphics[width=\textwidth]{figures/fig_curricula_depths.pdf}
    \caption{Pass@$1$ test par profondeur de recette, lissé sur cinq évaluations, pour les mêmes runs. La profondeur 4 (trois tâches de test, une seule à l'entraînement) reste à zéro dans tous les régimes et n'est pas tracée. L'ordre d'acquisition est identique partout : profondeur 1, puis 2, puis 3 ; les curriculums l'accélèrent sans le modifier.}
    \label{fig:acquisition_depths}
\end{figure}

\begin{figure}[htbp]
    \centering
    \includegraphics[width=\textwidth]{figures/fig_curricula_errors.pdf}
    \caption{Nombre moyen d'erreurs par épisode de test, lissé sur cinq évaluations. (a) Erreurs de format : action mal formée, plusieurs actions dans un tour, nom d'objet mal écrit. (b) Erreurs de planification : recette invalide, ingrédients insuffisants, objet introuvable. Les erreurs de format disparaissent en moins de mille mises à jour dans tous les régimes et réapparaissent lors du collapse du contrôle ; les erreurs de planification décroissent lentement sur toute la durée.}
    \label{fig:elimination_erreurs}
\end{figure}



\subsubsection{Mécanisme : qui produit du gradient}

Pourquoi l'entraînement sans curriculum s'effondre-t-il à grande taille de groupe, et pourquoi le curriculum l'en protège-t-il ? Nous proposons un mécanisme au niveau du gradient. Il est cohérent avec toutes nos courbes, mais deux de ses maillons n'ont pas été mesurés directement ; nous le présentons comme une hypothèse.

\paragraph{Le point de départ.} Cui et al.~\cite{cui2025entropy} établissent qu'à chaque pas de gradient de politique, la variation d'entropie est proportionnelle à l'opposé de la covariance entre la log-probabilité d'une action et son avantage. L'entropie baisse quand les actions déjà probables reçoivent un avantage positif et les actions improbables un avantage négatif ; elle monte dans le cas inverse. Tout se joue donc sur quelles actions sont récompensées et lesquelles sont punies, non sur le niveau de la récompense.

\paragraph{Qui produit du gradient.} Dans GRPO, l'avantage est calculé par tâche, sur le groupe des $G$ trajectoires du même prompt. Un groupe entièrement réussi ou entièrement échoué a un avantage nul et ne produit aucun gradient. Seuls les groupes mixtes en produisent. Contrôler l'apprentissage, c'est donc contrôler quels items sont mixtes, et quand.

\paragraph{Sans curriculum, la configuration a le mauvais signe.} Le modèle de base a du support sur les tâches difficiles : à trente tours, une recette de profondeur 3 donne typiquement une à trois réussites sur seize. Le groupe est mixte et produit du gradient. Mais la réussite est courte et faite d'actions routinières, donc probables : elle reçoit un avantage fortement positif, de l'ordre de $+3{,}9$ pour un succès sur seize. Les quinze échecs durent trente tours et se remplissent de dérive, d'actions improbables, et l'avantage négatif s'applique à chacun de leurs tokens, l'attribution se faisant à la trajectoire entière. Avantage positif sur du probable, négatif sur de l'improbable : l'entropie baisse à chaque pas, sur des dizaines de tâches à la fois. Le signal est correct, le modèle apprend, plus vite même qu'avec curriculum au même nombre de mises à jour. Mais l'entropie est une ressource finie : chaque groupe mixte à longs échecs en consomme, et rien ne la renouvelle. Vers 0{,}1, la politique est déterministe, et toute mise à jour fait alors sauter la divergence KL. C'est exactement la seconde signature de la section~4.2.5.

\paragraph{Ce que change le curriculum d'horizon.} Le plafond de dix tours agit de deux façons. D'abord, il transforme les groupes mixtes à longs échecs en groupes entièrement échoués : la réussite d'une profondeur 3 demande plus de dix tours, donc aucune des seize trajectoires ne réussit, l'avantage est nul, et l'entropie de la politique sur ces tâches est préservée parce qu'aucun gradient n'y touche. Ensuite, les seuls groupes mixtes restants sont des tâches résolubles en dix tours, dont les échecs sont courts et contiennent peu d'actions improbables. Au palier suivant, les tâches difficiles redeviennent mixtes, mais dans de meilleures conditions, briques fiables et échecs coupés à vingt tours, et leurs nouvelles réussites viennent d'actions que le modèle produisait peu : crafter un objet intermédiaire au deuxième tour, refaire un \texttt{get} après un \texttt{craft}, enchaîner plus de deux crafts. Avantage positif sur de l'improbable : l'entropie remonte. Le curriculum ne choisit pas les items à la main ; c'est le plafond qui décide quels items sont mixtes et quand. Le même raisonnement vaut pour le budget de tokens, qui coupe les dérives longues, et pour la profondeur, qui retire les tâches difficiles de l'échantillonnage.

\paragraph{Le corollaire pour $G=16$.} Avec seize trajectoires au lieu de huit, la probabilité qu'une tâche difficile compte au moins une réussite est plus grande. Il y a donc, sans curriculum, plus de groupes mixtes à longs échecs dès le départ, plus de gradient du mauvais signe, et un collapse plus tôt : dixième époque pour le contrôle à $G=16$, aucun collapse en 110 époques à $G=8$. Avec curriculum, ces tâches sont entièrement échouées quoi qu'il arrive, et $G=16$ n'a plus cet effet. L'autocurriculum de l'annexe~E, couru à $G=16$ sans plafond, s'est effondré à la septième époque, comme le prédit ce mécanisme.

\paragraph{Ce qui reste à mesurer.} Deux vérifications sont peu coûteuses et n'ont pas été faites : la longueur moyenne des trajectoires échouées du contrôle contre celle du curriculum d'horizon pendant son premier palier ; et la part d'actions rares, craft d'un intermédiaire ou \texttt{get} après un craft, dans les réussites juste après un changement de palier. Tant qu'elles ne sont pas faites, le mécanisme reste une hypothèse illustrée par les courbes d'entropie, non une loi.



\subsection{Comparaison à compute contrôlé}
\label{subsec:compute-controle}

La section~\ref{subsec:protocole} a fixé la convention : chaque mise à jour consomme 64 trajectoires, quelle que soit la taille de groupe, et les runs de tailles de groupe différentes se comparent en mises à jour, jamais en époques. Cette convention égalise le nombre de trajectoires, pas le calcul dépensé pour les produire. Or c'est là que les curriculums de budget prennent leur avantage : un plafond de 10 tours ou de 256 tokens raccourcit les épisodes, donc le temps de génération de chaque mise à jour. Une comparaison en mises à jour sous-estime cet avantage. Nous le mesurons ici avec trois indicateurs. Le premier est le nombre de tokens traités, prompt et trajectoire compris, cumulé sur le run : la génération des trajectoires domine le temps de chaque mise à jour, et ce compteur en est la mesure la plus proche, indépendante du matériel. Le second est le temps GPU, somme des durées de mise à jour : c'est le coût facturé, celui qu'un praticien doit budgéter, mais il dépend du matériel et de l'infrastructure. Le troisième est le nombre de tours d'environnement, c'est-à-dire d'appels au simulateur, qui deviendrait le premier coût dans un environnement réel. Il n'est pas enregistré pendant l'entraînement ; nous l'estimons à partir des évaluations périodiques : le nombre moyen de tours des 100 épisodes de test, plafonné par le palier d'horizon en cours, multiplié par les 64 trajectoires de chaque mise à jour écoulée. L'estimateur suppose que les épisodes d'entraînement ont, à un instant donné, la même longueur que ceux de test ; il donne un ordre de grandeur, non une mesure.

\paragraph{Effet de la taille de groupe sur la vitesse brute.} Un banc dédié, huit mises à jour par configuration à conditions fixées, donne 57{,}1~s par mise à jour à $G=8$ dans le réglage mémoire historique, contre 46{,}9~s à $G=16$ avec le cache de génération élargi, soit 18\,\% de moins. Le mécanisme est le partage du préfixe : les seize trajectoires d'une même tâche réutilisent le cache de son prompt. Élargir davantage le cache est contre-productif (61{,}5~s). Ce banc mélange deux changements, la taille de groupe et le cache, et ne mesure qu'une vitesse à politique fixée. La vitesse moyenne d'un run réel dépend en outre de la longueur des épisodes, donc de la politique et du régime de curriculum, comme le montre le tableau~\ref{tab:compute}.

\begin{table}[htbp]
\centering
\caption{Coût des lignées principales, chiffres arrêtés au gel du 3 septembre. Une mise à jour consomme 64 trajectoires. Heures GPU : somme des durées de mise à jour, évaluations exclues. Tokens : traités par le modèle, prompt et trajectoire compris. Tours : estimés d'après les évaluations périodiques (voir texte). Plateau : moyenne des dix dernières évaluations.}
\label{tab:compute}
\footnotesize
\setlength{\tabcolsep}{4pt}
\begin{tabular}{l c r r r r c c}
\toprule
\textbf{Lignée} & $\boldsymbol{G}$ & \textbf{Mises à jour} & \textbf{h GPU} & \textbf{Tokens ($10^9$)} & \textbf{Tours ($10^6$)} & \textbf{Meilleur} & \textbf{Plateau} \\
\midrule
\textsc{Réplication-FT} & 8 & 11\,564 & 238 & 0{,}87 & 12{,}4 & 69 & 61 \\
\textsc{Ancre-Mobile} & 8 & 8\,018 & 70 & 0{,}62 & 8{,}6 & 73 & 70 \\
\textsc{Horizon} & 16 & 7\,518 & 48 & 0{,}45 & 6{,}3 & \textbf{82} & \textbf{78} \\
\textsc{Profondeur} & 16 & 7\,904 & 66 & 0{,}56 & 7{,}9 & 80 & 74 \\
\textsc{Budget} & 16 & 7\,194 & 51 & 0{,}43 & 7{,}2 & 80 & 75 \\
\textsc{Contrôle-G16} & 16 & 1\,434 & 16 & 0{,}15 & 1{,}9 & 54 & collapse \\
\bottomrule
\end{tabular}
\end{table}

\paragraph{Lecture.} Trois constats se dégagent du tableau et de la figure~\ref{fig:pass1-tokens}, qui place chaque lignée sur l'axe des tokens traités.

\begin{itemize}
    \item \textbf{Les curriculums de budget sont les moins chers par mise à jour.} \textsc{Horizon} coûte 23~s par mise à jour et \textsc{Budget} 26~s, contre 39~s pour le contrôle à horizon et budget pleins dès le départ, et 30~s pour \textsc{Profondeur}, qui joue 30 tours dès le début. Sur 80 époques, \textsc{Horizon} a traité 0{,}45 milliard de tokens là où \textsc{Profondeur} en a traité 0{,}56 pour un plateau inférieur de quatre points. En tours d'environnement, \textsc{Horizon} en consomme moitié moins que le \emph{full finetuning} pour un plateau supérieur de dix-sept points. L'explication est mécanique : les paliers initiaux coupent les épisodes courts.
    \item \textbf{À tokens égaux, les trois curriculums dominent les deux références.} À 0{,}4 milliard de tokens, \textsc{Horizon} et \textsc{Budget} sont entre 75 et 80\,\%, \textsc{Profondeur} vers 68, la référence à $G=8$ vers 55 et le \emph{full finetuning} vers 55 également. Le \emph{full finetuning} atteint son plateau de 61\,\% après 0{,}87 milliard de tokens et 238 heures de GPU ; \textsc{Horizon} dépasse ce niveau après 0{,}1 milliard et une dizaine d'heures.
    \item \textbf{Vitesse ou asymptote.} Sur 80 époques, les trois curriculums ont un plateau supérieur à celui de la référence à $G=8$ après 170 époques, 74 à 78 contre 70. Ils atteignent donc plus haut, et plus vite. Mais la référence montait encore à son arrêt, et aucun run n'a été poursuivi jusqu'à saturation complète : nous ne pouvons pas dire si les asymptotes diffèrent, seulement que les curriculums y arrivent plus tôt.
\end{itemize}

\paragraph{Effet de la taille de groupe seule.} Un seul couple isole cet effet avant collapse : \textsc{Ancre-Mobile} à $G=8$ et \textsc{Contrôle-G16}, à recette identique. À nombre égal de trajectoires, le second monte plus vite, 0{,}45 de récompense d'entraînement contre 0{,}30 à la millième mise à jour. C'est cohérent avec une estimation de l'avantage plus précise à seize trajectoires par tâche, mais c'est une observation sur un couple de runs, non une loi, et elle ouvre plus de questions qu'elle n'en ferme : à partir de quelle taille de groupe le gain de précision est-il mangé par la déterminisation accélérée ? Nous n'avons pas de run à $G=32$ pour le dire.

\begin{figure}[htbp]
    \centering
    \includegraphics[width=0.95\textwidth]{figures/fig_pass1_tokens.pdf}
    \caption{Pass@$1$ test, lissé sur cinq évaluations, en fonction des tokens traités cumulés depuis le début de chaque lignée. À abscisse égale, les runs ont consommé le même volume de génération ; l'écart vertical mesure l'efficacité de chaque régime. Les lignées \textsc{Réplication-FT}, \textsc{Ancre-Mobile}, \textsc{Profondeur} et \textsc{Budget} comprennent des reprises après purge (annexe~C).}
    \label{fig:pass1-tokens}
\end{figure}


\subsection{Expériences complémentaires et pistes exploratoires}
\label{sec:inference_ablations}

Les sections précédentes ont porté sur l'entraînement. Celle-ci rassemble trois types de mesures qui en donnent le cadre. D'abord des mesures sans entraînement : l'oracle pass@$k$, qui met à l'épreuve l'hypothèse du plafond posée en section~2.3, le prompt \emph{few-shot} et la comparaison entre générations de modèles. Ensuite une expérience que nous n'avons pas pu refaire dans la configuration stabilisée, l'interaction entre \emph{few-shot} et RL. Enfin trois pistes exploratoires, menées en marge du programme central et arrêtées avant d'aboutir : la recomposition de trajectoires par SNIS, le transfert entre planification mono-tour et contrôle multi-tour, et l'autocurriculum MAGELLAN. Chacune est résumée ici en un paragraphe et développée en annexe. Le tableau~\ref{tab:recap_section4} récapitule l'ensemble des configurations de la section~4.

\subsubsection{Oracle pass@$k$ et hypothèse du plafond}

Le pass@$k$ est la fraction des tâches résolues au moins une fois en $k$ tirages indépendants. Notre oracle tire, pour chacune des 100 tâches de test, vingt trajectoires à température $1{,}0$ et 30 tours au plus, puis calcule le pass@$k$ pour $k$ de 1 à 20 par l'estimateur non biaisé de Chen et al.~\cite{chen2021codex}, globalement et par profondeur. Le pass@$1$ qui en sort est la moyenne de vingt passes ; c'est la référence retenue pour le modèle non entraîné en section~4.1.

\paragraph{La structure du problème.} Sur Qwen2.5-3B non entraîné, le pass@$1$ vaut 10{,}3\,\% et le pass@20 vaut 47\,\%. Ces deux nombres disent deux choses différentes. Le premier mesure ce que le modèle réussit de façon fiable. Le second mesure ce qu'il sait faire au moins occasionnellement : son support. L'écart entre les deux, 37 points, est la marge qu'un algorithme de renforcement peut espérer convertir en fiabilité sans rien apprendre de nouveau. La décomposition par profondeur sépare deux régimes. En profondeur 1, le pass@20 atteint 94\,\% : presque tout est atteignable. En profondeur 2, il atteint 44\,\%, pour un pass@$1$ de 4\,\% : le modèle sait résoudre près de la moitié de ces tâches, mais rarement. C'est un mur de fiabilité, et une large marge pour le RL. En profondeurs 3 et 4, le pass@20 est nul : zéro succès sur 500 tirages en profondeur 3, zéro sur 60 en profondeur 4. C'est un mur de capacité. Sa conséquence pour GRPO est directe : sans aucun succès, tous les groupes sont uniformément en échec, l'avantage relatif est nul et le gradient aussi. Le modèle de 7 milliards de paramètres, non entraîné, déplace les niveaux sans lever le mur : pass@$1$ de 33\,\%, pass@20 de 68\,\%, et par profondeur 94, 76 et 32\,\%, puis zéro sur les trois tâches de profondeur 4.

\paragraph{L'hypothèse du plafond.} Yue et al.~\cite{yue2025rlvr} observent, sur des tâches de raisonnement mono-tour, que les modèles entraînés par RL dominent leur modèle de base à petit $k$ mais lui deviennent inférieurs à grand $k$, et concluent que le RL concentre la masse de probabilité sur des chemins que le modèle de base savait déjà produire, sans en ajouter. Le pass@$k$ du modèle de base serait ainsi un plafond pour le RL. Nous avons posé cette thèse en section~2.3 comme une hypothèse à tester, et nos mesures la mettent en tension de deux manières.

La première vient d'un modèle entraîné modestement, un point de contrôle LoRA de nos premiers runs à 32\,\% de pass@$1$. Nous l'avons choisi parce qu'il est loin de nos meilleurs modèles : si l'hypothèse cède déjà pour lui, elle ne tient pas à un cas extrême. Son pass@20 atteint 65\,\%, contre 47 pour le modèle de base. Par profondeur, le pass@20 passe de 94 à 100\,\% en profondeur 1, de 44 à 78\,\% en profondeur 2, et de 0 à 8\,\% en profondeur 3. Autrement dit, quatorze tâches de profondeur 2 et deux tâches de profondeur 3 que le modèle de base n'avait résolues dans aucun de ses vingt tirages le sont désormais au moins une fois. Le RL n'a donc pas seulement rendu fiables des succès existants ; il a élargi le support mesuré.

La seconde vient de \textsc{Horizon}. Ce run atteint un pass@$1$ de 82\,\%, alors que le modèle de base plafonne à 47\,\% de pass@20 : la fiabilité après RL dépasse le support avant RL, mesuré à vingt tirages. Nous présentons ce résultat comme une mise en tension, non comme une réfutation, pour trois raisons. Le budget $k=20$ est petit : Yue et al. observent leurs croisements pour des $k$ de l'ordre de la centaine, et le pass@$k$ croît avec $k$. Leurs modèles sont entraînés quelques centaines de mises à jour, les nôtres plusieurs milliers : ils réservent eux-mêmes la question d'un RL plus long, et ProRL~\cite{liu2025prorl}, qui entraîne longtemps, observe comme nous que le RL dépasse la frontière du modèle de base. Enfin l'analyse par tâche, qui dirait sur quelles tâches précises \textsc{Horizon} réussit là où le modèle de base échoue vingt fois, n'a pas été menée : le gel des expériences est intervenu avant. La profondeur 4 reste par ailleurs à zéro dans tous nos runs, et le déséquilibre du jeu d'entraînement borne ce que TextCraft peut prouver sur ce mur.

\subsubsection{Prompt few-shot : gain immédiat et amorçage}

Le prompt \emph{few-shot} fournit au modèle $k$ épisodes résolus avant la tâche courante, sans modifier ses poids~\cite{brown2020gpt3}. Nos exemples proviennent du réservoir de 70 recettes exclu de l'entraînement et du test. Sur Qwen2.5-3B, cinq exemples insérés comme tours de dialogue portent le pass@$1$ de 10{,}3 à 27{,}7\,\%, en moyennes de passes. Le format conditionne le signe du gain : les mêmes exemples concaténés en bloc dans le message système font tomber le score à 5 ou 8\,\%. Le modèle imite alors la forme du bloc et produit lui-même des observations au lieu d'attendre celles de l'environnement.

Le \emph{few-shot} joue aussi un rôle d'amorçage pour le RL. Sans exemples, le pass@20 est nul en profondeur 3 : aucun groupe GRPO mixte n'y est possible. Avec vingt exemples, le pass@10 y atteint 16\,\%. C'est un gain rapide à étudier : il crée du signal là où le gradient était identiquement nul. Nous n'avons qu'une comparaison propre de son effet en entraînement, deux runs anciens à paliers longs, identiques au prompt près, à 41\,\% sans exemples et 49\,\% avec dix exemples ; tous deux ont été interrompus par des pannes d'infrastructure, et le run prévu dans la configuration stabilisée n'a pas été couru. La lecture prudente est que le \emph{few-shot} accélère les premières époques, avec un risque de sur-imitation des recettes des exemples, sans que l'on sache s'il change le plateau.

\subsubsection{Générations de modèles}

Une partie du mur mesuré ne relève pas de l'algorithme d'apprentissage mais du modèle de départ. Dans notre harnais, Qwen3.5-4B non entraîné atteint 78{,}5\,\% de pass@$1$ en moyenne de dix-huit passes, au-dessus de l'objectif du papier, et 97\,\% de pass@18 : presque toutes les tâches lui sont accessibles. Le papier rapporte 45\,\% pour la génération intermédiaire Qwen3-4B. Cette progression s'accorde avec l'exposition croissante des modèles récents aux traces agentiques dès leur pré-entraînement (section~2.3). La leçon pratique est simple : pour un déploiement, partir du modèle le plus récent est le gain le moins coûteux. Nous sommes restés sur Qwen2.5-3B pour une raison de méthode, la comparaison à recette égale avec le papier de référence. Les conditions de stabilité établies en section~4.2 tiennent au RL multi-tour et non au modèle ; il est plausible qu'un modèle récent, déjà entraîné par RL sur des environnements, soit moins instable, mais nous ne l'avons pas mesuré.

\subsubsection{Trois pistes exploratoires}

\paragraph{Recomposition de trajectoires par SNIS.} L'idée, présentée en section~2.6, consiste à réutiliser les sous-trajectoires d'un même groupe GRPO en les recombinant, pondérées par échantillonnage préférentiel auto-normalisé. Le run exploratoire a recollé des sous-trajectoires multi-tours sans alignement causal : les séquences ainsi construites juxtaposent des observations issues de contextes différents, et l'apprentissage s'est effondré. C'est la confirmation empirique du verrou identifié dans la dérivation. La variante qui restreindrait la recomposition à des sous-séquences causalement cohérentes n'a pas été formalisée. L'annexe~B donne la dérivation et le bilan.

\paragraph{Transfert entre planification mono-tour et contrôle multi-tour.} La question est celle de deux régimes d'apprentissage : la planification mono-tour, où le modèle émet un plan complet en une seule génération, et le contrôle multi-tour, où il choisit chaque action après observation. Le premier coûte bien moins cher à entraîner. Apprendre dans l'un aide-t-il dans l'autre, et un régime alterné réduirait-il le coût total ? Nos mesures préliminaires, sans entraînement dédié, montrent une asymétrie : la boucle de rétroaction vaut 23 points pour un modèle non entraîné, le RL interactif améliore modestement la planification, et un RL de planification pure détruit la compétence interactive. Ces résultats étaient attendus et ne testent pas l'alternance ; ils sont détaillés en annexe~D et la question est reprise en ouverture.

\paragraph{Autocurriculum MAGELLAN.} Les trois curriculums de la section~4.3 fixent la difficulté à la main. MAGELLAN~\cite{gaven2025magellan} la fait choisir au modèle : les tâches sont échantillonnées proportionnellement au progrès d'apprentissage qu'un estimateur appris leur prédit. Nous en avons porté le mécanisme et couru un run sur la recette commune. Il a atteint 45\,\% puis s'est effondré à la septième époque, plus tôt que le contrôle sans curriculum. La prédiction écrite avant le run, à savoir que l'estimateur cesserait de présenter la profondeur 4, s'est vérifiée ; mais le même signal, qui mesure du changement, a concentré l'échantillonnage sur les tâches faciles au moment où le collapse faisait osciller les prédictions, et l'a accéléré. L'annexe~E présente la méthode, notre port et cette lecture.

\begin{table}[htbp]
\centering
\caption{Récapitulatif des configurations de la section~4. Pass@$1$ sur les 100 tâches de test ; pour les runs, meilleur score et plateau (moyenne des dix dernières évaluations) ; pour les modèles non entraînés, moyenne des passes de l'oracle. Pass@$k$ à température $1{,}0$. Chiffres arrêtés au gel du 3 septembre.}
\label{tab:recap_section4}
\small
\begin{tabular}{l c c l}
\toprule
\textbf{Configuration} & \textbf{pass@1 (\%)} & \textbf{pass@$k$} & \textbf{Nature} \\
\midrule
Qwen2.5-3B, zéro-shot & 10{,}3 & 47 ($k{=}20$) & référence interne \\
Qwen2.5-3B, few-shot 5 exemples & 27{,}7 & 60 ($k{=}10$) & référence interne \\
Qwen2.5-7B, zéro-shot & 33{,}2 & 68 ($k{=}20$) & référence interne \\
Qwen3.5-4B, zéro-shot & 78{,}5 & 97 ($k{=}18$) & référence inter-générations \\
\textsc{Réplication-FT} & 69 / 61 & --- & RL, \emph{full finetuning} \\
\textsc{Ancre-Mobile} ($G{=}8$) & 73 / 70 & --- & RL, LoRA, sans curriculum \\
\textsc{Contrôle-G16} & 54 / collapse & --- & RL, LoRA, sans curriculum \\
\textsc{Horizon} & \textbf{82 / 78} & --- & RL, curriculum d'horizon \\
\textsc{Profondeur} & 80 / 74 & --- & RL, curriculum de profondeur \\
\textsc{Budget} & 80 / 75 & --- & RL, curriculum de budget \\
\textsc{Magellan} & 45 / collapse & --- & RL, autocurriculum \\
AgentGym-RL-3B & 75 & --- & externe (papier) \\
Gemini 3.5 Flash & 99 & --- & externe, borne haute \\
\bottomrule
\end{tabular}
\end{table}










%---------------------------------------------------------------------------------------



\clearpage
\section{Délimitation du périmètre et ouverture}

\paragraph{Le choix du périmètre.} Le sujet initial de ce stage était l'auto-amélioration des agents fondés sur des modèles de langage, avec pour horizon les agents d'achat que Criteo développe. Ce sujet ouvre sur plusieurs domaines à la fois : l'ingénierie des systèmes agentiques, l'apprentissage par renforcement, l'architecture cognitive, et la recommandation. Nous avons choisi d'en isoler le cœur technique : l'amélioration d'une politique par renforcement dans un environnement interactif, avec ses conditions de stabilité et ses leviers d'efficacité. Ce choix s'appuie sur une analyse du sujet. Un agent d'achat qui s'améliore de lui-même sera, en dernière analyse, un modèle entraîné par renforcement sur un environnement de commerce ; comprendre ce que cet entraînement exige, ce qui le fait échouer et ce qui l'accélère, est le préalable à toute adaptation au domaine. Il s'appuie aussi sur un intérêt personnel pour l'apprentissage par renforcement, dont ce sujet renoue avec les questions classiques. L'objectif était donc de construire une compréhension théorique et pratique solide du mécanisme, sur un environnement simple où il peut être disséqué, pour que le travail suivant puisse se concentrer sur son adaptation à la recommandation et à l'achat. Cette section rappelle ce qui a été laissé de côté, ce que nos résultats appellent directement, et les questions plus larges dans lesquelles ils s'inscrivent.

\paragraph{Ce qui a été laissé de côté.}
\begin{itemize}
    \item \textbf{Mémoire et amélioration non paramétrique.} Un agent peut s'améliorer à poids constants, en accumulant de l'expérience dans une mémoire externe consultée à l'inférence, de la gestion hiérarchique du contexte~\cite{packer_memgpt_2023} aux banques de compétences. Cette voie est complémentaire de la nôtre. Nos mesures \emph{few-shot} (section~4.5) en donnent un indice : le contenu du contexte crée du signal là où le gradient seul n'en avait pas.
    \item \textbf{Architectures multi-agents et optimisation du harnais.} L'orchestration de sous-agents~\cite{kimi_k25_2026} et la découverte automatique d'architectures agentiques~\cite{hu_adas_2025} optimisent le système autour du modèle. Nous nous sommes limités à la capacité décisionnelle d'un agent unique, préalable à toute composition : un orchestrateur ne compense pas des sous-agents dont la politique est instable.
    \item \textbf{Autres algorithmes et densification de la récompense.} Les méthodes à réseau critique, le fine-tuning supervisé itératif sur trajectoires filtrées et les modèles de récompense de processus~\cite{lightman2023verify} n'ont pas été retenus, au profit de la frugalité de GRPO. Notre expérience de récompense façonnée (section~4.2) éclaire ce choix : densifier sans corrompre est le problème que ces modèles cherchent à résoudre, à un coût d'estimation qui reste un verrou en multi-tour.
    \item \textbf{Génération d'environnements.} Traiter l'environnement comme un objet à générer, qu'un agent propose ses propres tâches vérifiables~\cite{zhou2025sca} ou qu'un modèle de monde synthétise les expériences~\cite{dreamgym2025}, est la frontière la plus directement liée au constat de ce rapport : si le passage à l'échelle des environnements est le moteur des capacités agentiques (sections~1.1 et~2.2), leur génération en est la suite logique. L'infrastructure qu'elle exige dépasse ce que nous pouvions construire.
\end{itemize}

\paragraph{Ce que nos résultats appellent directement.}
\begin{itemize}
    \item \textbf{Re-stratifier le jeu d'entraînement.} L'unique recette de profondeur 4 est le principal facteur confondant de l'étude. L'arbre de recettes permet de construire un jeu équilibré par profondeur ; y rejouer les curriculums serait le test le plus direct du mur de capacité.
    \item \textbf{Répéter les runs.} Chaque configuration a été courue une fois. Plusieurs graines par configuration donneraient la variabilité d'un entraînement à l'autre, que nous n'avons pas mesurée.
    \item \textbf{Éprouver la recette ailleurs.} L'ancre mobile et les curriculums ont été mis au point sur un jeu textuel. Leur robustesse reste à établir sur d'autres environnements, ceux d'AgentGym-RL d'abord, puis un environnement d'achat en ligne de type WebShop, plus proche des cas d'usage de Criteo. C'est aussi la condition pour tirer des conclusions générales des analyses de couverture et d'erreurs de la section~4 : l'ordre d'acquisition des difficultés, l'élimination des erreurs de forme avant celles de fond, et l'extension du support au-delà du modèle de base sont des observations sur un environnement ; leur portée théorique demande de les retrouver sur plusieurs.
    \item \textbf{Exploiter davantage chaque trajectoire.} Notre recette est strictement on-policy : chaque trajectoire sert à une mise à jour puis est jetée. Réutiliser les trajectoires collectées pour plusieurs mises à jour, avec la correction hors-politique et le clipping que PPO fournit, est un levier d'efficacité que nous avons délibérément laissé de côté pour stabiliser rapidement. Il mérite d'être repris, en mesurant s'il améliore le rapport entre performance atteinte et trajectoires générées, et s'il modifie la stabilité observée à grande taille de groupe.
    \item \textbf{La récompense pour la recommandation.} Nos conclusions valent pour une récompense binaire exacte. Les cas d'usage de Criteo, où le signal n'est ni binaire ni exact, demanderont un travail propre sur la récompense, éclairé par ce que la section~4.2 a montré des termes façonnés.
\end{itemize}

\paragraph{Ce que nos résultats ouvrent.} Trois questions relient ce travail aux recherches sur le méta-apprentissage, au sens d'un apprentissage dont l'organisation elle-même est apprise ou adaptée, et aux agents d'achat.
\begin{itemize}
    \item \textbf{Combiner et adapter les curriculums.} Nos trois régimes agissent chacun sur un levier, l'horizon, le budget ou la profondeur, et arrivent à des plateaux voisins. Rien ne dit que leurs effets s'additionnent ; rien ne dit le contraire. Leur combinaison, et le passage d'un calendrier fixe à des paliers déclenchés par la performance, comme nous l'avons fait pour la profondeur, sont les extensions les plus simples.
    \item \textbf{Alterner planification et interaction.} L'annexe~D montre une asymétrie : la boucle de rétroaction vaut 23 points, le RL interactif transfère peu vers la planification, et un RL de planification pure détruit le contrôle. La question ouverte est celle d'un régime alterné, où un apprentissage peu coûteux en boucle ouverte préparerait l'apprentissage interactif. Dupoux, LeCun et Malik~\cite{dupoux2026why} proposent une architecture où un système apprenant par observation et un système apprenant par action sont orchestrés par un méta-contrôleur ; planification mono-tour et contrôle multi-tour en sont une instanciation naturelle pour les agents LLM.
    \item \textbf{Faire choisir la difficulté par le modèle.} Nos curriculums fixent la difficulté à la main, selon une notion qui est la nôtre, la profondeur d'une recette. Dans un environnement d'achat, cette difficulté n'est pas connue à l'avance : rien ne dit combien d'appels d'outils, de pages ou de comparaisons une requête demande, et rien ne garantit que la difficulté perçue par le concepteur coïncide avec celle du modèle. Les bancs d'essai les plus exigeants rencontrent le même problème ; Toolathlon~\cite{toolathlon2025} le contourne en mesurant la difficulté d'une tâche par le nombre de tours qu'un modèle propriétaire de référence met à la résoudre, un proxy qui dépend du modèle choisi et fige la difficulté à un instant. Ce détour montre que la difficulté des tâches est au cœur de l'apprentissage sur environnement, puisqu'elle décide de la stabilité du RL. Les agents autotéliques y répondent en laissant l'agent estimer et choisir ses buts : MAGELLAN~\cite{gaven2025magellan}, dont l'annexe~E rapporte notre port, prédit le progrès d'apprentissage à partir des représentations du modèle ; HERAKLES~\cite{carta2025herakles} fait choisir des sous-buts atteignables par une politique de haut niveau et les compile en compétences ; BOSS~\cite{zhang2023boss} fait proposer par un LLM des enchaînements de compétences de difficulté croissante. Notre run MAGELLAN en a montré une condition d'emploi : l'autocurriculum suppose un apprenant stable. Le prolongement le plus cohérent avec ce rapport serait de le brancher sur une recette stabilisée par un curriculum de budget, puis sur un générateur de tâches ou d'environnements, de sorte que la difficulté proposée suive en permanence les capacités du modèle. Pour un agent d'achat, c'est la voie qui dispense de définir la difficulté à la main.
\end{itemize}

Ces frontières dessinent en retour la portée de ce travail. La recette de stabilisation et les régimes de curriculum sont des propriétés de la politique et de son signal d'apprentissage. Elles ne dépendent ni du harnais, ni de la mémoire, ni du mécanisme de génération des tâches, ni de l'environnement particulier sur lequel elles ont été établies, sous réserve des vérifications appelées ci-dessus. C'est ce qui les rend transférables au domaine visé : un agent d'achat s'entraînera par le même algorithme, rencontrera les mêmes instabilités, et aura besoin des mêmes outils pour doser sa difficulté.






\clearpage
\section{Conclusion}

Ce rapport a étudié l'amélioration d'agents fondés sur des modèles de langage par apprentissage par renforcement, dans un environnement interactif à récompense vérifiable. Son point de départ est un constat partagé par les constructeurs de modèles : le facteur limitant du RL agentique n'est plus le volume de données annotées mais l'environnement d'entraînement lui-même, sa difficulté, sa diversité et son vérificateur. Nous avons pris ce constat comme cadre et non comme programme. Le passage à l'échelle des environnements, qui vise à doter les modèles de capacités agentiques générales, relève des constructeurs. Notre question est plus étroite : comment un modèle donné s'améliore-t-il au mieux sur un environnement donné, celui d'un agent d'achat à terme. Nous avons donc travaillé sur un environnement unique et finement instrumenté, TextCraft, pour y chercher ce qui conditionne la stabilité et l'efficience de cet apprentissage, à partir de la réplication d'un système publié et dans les contraintes matérielles de Criteo.

Le premier de nos deux travaux porte sur la stabilisation et l'efficience de l'entraînement. Deux conditions préalables ont été établies : le masquage des observations dans le calcul du gradient, dont le défaut ne se signale par aucune métrique, et l'attention portée à la conception de la récompense, dont nos essais de termes façonnés ont montré qu'un bonus d'apparence raisonnable est optimisé pour lui-même dès que l'on suit sa propagation dans l'avantage, alors que les coefficients de ce type restent, dans la littérature appliquée, trop souvent posés sans justification. La recette du papier de référence en \emph{full finetuning} est stable sur 247 époques et atteint 69\,\%, à un coût en temps et en mémoire qui interdisait toute campagne d'expériences. L'adaptation de rang faible supprime ce coût, mais sa recette publiée ne tient pas dans notre régime : à référence KL fixe, la divergence croît géométriquement et le run s'effondre, quels que soient le rang et le coefficient testés, là où le \emph{full finetuning} dérive lentement. Le correctif que nous proposons déplace la référence toutes les quatre époques sur la politique courante. Il supprime la dérive, porte le même modèle à 73\,\%, et un croisement d'expériences montre que c'est le régime de référence, non la valeur du coefficient, qui gouverne la stabilité. La mémoire libérée a permis de doubler la taille des groupes GRPO. Enfin, l'analyse de quinze runs distingue deux modes de collapse : par la divergence KL à référence fixe, par l'extinction de l'entropie à grande taille de groupe. L'ancre mobile règle le premier, non le second.

Le second travail compare trois régimes de curriculum, sur l'horizon d'interaction, le budget de génération et la profondeur des tâches, à recette identique et à compute contrôlé, face à deux références sans curriculum. À grande taille de groupe, la référence sans curriculum s'effondre par extinction de l'entropie ; les trois curriculums traversent la même zone sans accident et poursuivent leur montée sur 80 époques. Le curriculum d'horizon atteint 82\,\% de réussite, pour un plateau de 78, au-dessus des 75\,\% publiés par le papier de référence sur une infrastructure multi-GPU ; les deux autres atteignent 80\,\%, pour des plateaux de 74 et 75. Le curriculum n'est donc pas seulement un accélérateur ; dans ce régime, il est une condition de stabilité. Il n'empêche pas l'entropie de décroître, il empêche son effondrement brutal. Les curriculums de budget sont aussi les moins coûteux : moitié moins de tours d'environnement que le \emph{full finetuning} pour un plateau supérieur de dix-sept points. Les analyses stratifiées montrent enfin que tous les régimes acquièrent les profondeurs dans le même ordre et éliminent les erreurs de forme avant celles de fond ; les curriculums accélèrent ces étapes sans en changer l'ordre, et aucun ne franchit la profondeur 4.

Un troisième volet confronte ces résultats à l'hypothèse selon laquelle le RL ne ferait que concentrer un support préexistant. Nos mesures vont dans le sens contraire sans le démontrer : un modèle entraîné modestement résout à vingt tirages des tâches que le modèle de base ne résolvait dans aucun tirage, et la fiabilité du meilleur run dépasse le support mesuré du modèle de base. 

Ces résultats ont leurs limites, que le rapport a cherché à rendre explicites : un run par configuration, un seul environnement, une seule recette d'entraînement commune, un jeu d'entraînement déséquilibré en profondeur, et des modèles récents qui résolvent la tâche sans entraînement. Ils ont aussi une portée : la recette de stabilisation et les régimes de curriculum tiennent à la politique et à son signal d'apprentissage, non à l'environnement, et ce sont des outils essentiels à l'amélioration des futurs agents d'achat de Criteo.

Pour Criteo, ce travail laisse trois choses. Une synthèse d'une littérature très riche, qui avance dans de nombreuses directions à la fois et dont le tri, pour des technologies encore en cours de définition, est difficile à faire du point de vue des agents d'achat. Une infrastructure d'expérimentation opérationnelle, entraînement GRPO multi-tour stabilisé, environnement instrumenté et journaux exploitables, qui a porté un modèle de trois milliards de paramètres au-delà du résultat publié sur la même tâche. Et une compréhension des conditions dans lesquelles un tel entraînement réussit ou échoue, qui permettra au travail suivant de se concentrer sur ce qui est propre au domaine : la récompense, l'environnement d'achat et la mesure de sa difficulté.










\newpage



\newpage

\begin{thebibliography}{99}

\bibitem{vaswani2017attention} Vaswani et al. Attention is all you need. NeurIPS 30, 2017.
\bibitem{ouyang2022training} Ouyang et al. Training language models to follow instructions with human feedback. NeurIPS 35, 2022.
\bibitem{zelikman2022star} Zelikman, Wu, Mu, Goodman. STaR: Bootstrapping reasoning with reasoning. NeurIPS 35, 2022.
\bibitem{singh2024beyond} Singh, Co-Reyes, Agarwal et al. Beyond human data (ReST-EM). TMLR 2024. arXiv:2312.06585.
\bibitem{huang2024self} Huang, Block, Foster et al. Self-improvement in language models: The sharpening mechanism. arXiv:2412.01951, 2024.
\bibitem{schulman2017proximal} Schulman et al. Proximal policy optimization algorithms. arXiv:1707.06347, 2017.
\bibitem{shao2024deepseekmath} Shao et al. DeepSeekMath. arXiv:2402.03300, 2024.
\bibitem{rafailov2024direct} Rafailov et al. Direct preference optimization. NeurIPS 36, 2024.
\bibitem{deepseekr1_2025} DeepSeek-AI. DeepSeek-R1. arXiv:2501.12948, 2025.
\bibitem{openai_o1_2024} OpenAI. Learning to reason with LLMs (o1). Blog, 2024.
\bibitem{liu2025drgrpo} Liu et al. Understanding R1-Zero-like training (Dr. GRPO). arXiv:2503.20783, 2025.
\bibitem{yao2023react} Yao et al. ReAct. ICLR 2023.
\bibitem{xi_agentgym_2025} Xi, Ding, Chen et al. AgentGym. ACL 2025. arXiv:2406.04151.
\bibitem{xi_agentgym-rl_2025} Xi, Huang, Liao et al. AgentGym-RL. arXiv:2509.08755, 2025.
\bibitem{xi_agentgymrl_openreview} OpenReview. AgentGym-RL. openreview.net/forum?id=ZgCCDwcGwn, ICLR 2026.
\bibitem{prasad2023adapt} Prasad et al. ADaPT. Findings of NAACL 2024. arXiv:2311.05772.
\bibitem{carta2023glam} Carta, Romac, Wolf, Lamprier, Sigaud, Oudeyer. GLAM. ICML 2023. arXiv:2302.02662.
\bibitem{zhang2023boss} Zhang et al. BOSS. CoRL 2023. arXiv:2310.10021.
\bibitem{gaven2025magellan} Gaven, Carta, Romac, Colas, Lamprier, Sigaud, Oudeyer. MAGELLAN. arXiv:2502.07709, 2025.
\bibitem{carta2025herakles} Carta, Romac, Gaven, Oudeyer, Sigaud, Lamprier. HERAKLES. arXiv:2508.14751, 2025. % à recouper
% foster2026autocurriculum SUPPRIMÉE (placeholder) → remplacée par rajaraman2026curriculum (vérifiée 27/08)
\bibitem{forestier2017imgep} Forestier, Portelas, Mollard, Oudeyer. IMGEP. arXiv:1708.02190, 2017.
\bibitem{colas2022autotelic} Colas, Karch, Sigaud, Oudeyer. Autotelic agents (survey). JAIR 74, 2022.
\bibitem{bengio2009curriculum} Bengio, Louradour, Collobert, Weston. Curriculum learning. ICML 2009.
\bibitem{zhou2025sca} Zhou, Levine, Weston, Li, Sukhbaatar. Self-challenging language model agents. arXiv:2506.01716, 2025.
\bibitem{dreamgym2025} Chen et al. DreamGym. arXiv:2511.03773, 2025.
\bibitem{valmeekam2023planning} Valmeekam et al. On the planning abilities of LLMs. NeurIPS 36, 2023.
\bibitem{kambhampati2024llmmodulo} Kambhampati et al. LLM-modulo frameworks. arXiv:2402.01817, 2024.
\bibitem{brown2020gpt3} Brown et al. Language models are few-shot learners. NeurIPS 33, 2020.
\bibitem{cobbe2021verifiers} Cobbe et al. Training verifiers to solve math word problems. arXiv:2110.14168, 2021.
\bibitem{wang2022selfconsistency} Wang et al. Self-consistency. arXiv:2203.11171, 2022.
\bibitem{yao2023tot} Yao et al. Tree of Thoughts. arXiv:2305.10601, 2023.
\bibitem{madaan2023selfrefine} Madaan et al. Self-Refine. arXiv:2303.17651, 2023.
\bibitem{shinn2023reflexion} Shinn et al. Reflexion. arXiv:2303.11366, 2023.
\bibitem{snell2024scaling} Snell, Lee, Xu, Kumar. Scaling LLM test-time compute. arXiv:2408.03314, 2024.
\bibitem{lightman2023verify} Lightman et al. Let's verify step by step. arXiv:2305.20050, 2023.
\bibitem{skalse2022defining} Skalse, Howe, Krasheninnikov, Krueger. Defining and characterizing reward hacking. NeurIPS 35, 2022.
\bibitem{owen2013montecarlo} Owen. Monte Carlo Theory, Methods and Examples. Stanford, 2013.
\bibitem{schulman2025lora} Schulman et al. LoRA without regret. Blog Thinking Machines, 2025.
\bibitem{qwen3_2025} Qwen Team. Qwen3 technical report. arXiv:2505.09388, 2025.
\bibitem{qwen3coder2025} Qwen Team. Qwen3-Coder. Blog Qwen, 2025.
\bibitem{qwen35_2026} Qwen Team. Qwen3.5 release notes. 2026. % URL à vérifier
\bibitem{kimik2_2025} Kimi Team. Kimi K2: Open agentic intelligence. arXiv:2507.20534, 2025.
\bibitem{kimi_k25_2026} Kimi Team. Kimi K2.5: Visual agentic intelligence. arXiv:2602.02276, 2026.
\bibitem{magistral2025} Mistral AI. Magistral. arXiv:2506.10910, 2025.
\bibitem{taubench2024} Yao, Shinn, Razavi, Narasimhan. Tau-bench. arXiv:2406.12045, 2024.
\bibitem{tau2bench2025} Barres et al. Tau2-bench. arXiv:2506.07982, 2025.
\bibitem{toolathlon2025} Li et al. The Tool Decathlon (Toolathlon). ICLR 2026. arXiv:2510.25726.
\bibitem{mcpuniverse2025} Salesforce AI Research. MCP-Universe. arXiv:2508.14704, 2025.
\bibitem{mcpatlas2026} Scale AI. MCP-Atlas. arXiv:2602.00933, 2026.
\bibitem{zhou2023webarena} Zhou et al. WebArena. arXiv:2307.13854, 2023.
\bibitem{shridhar2020alfworld} Shridhar et al. ALFWorld. arXiv:2010.03768, 2020.
\bibitem{jimenez2023swebench} Jimenez et al. SWE-bench. arXiv:2310.06770, 2023.
\bibitem{openai2024swebenchverified} OpenAI. Introducing SWE-bench Verified. Blog, 2024.
\bibitem{anthropic2024agents} Anthropic. Building effective agents. 2024.
\bibitem{anthropic2024mcp} Anthropic. Introducing the Model Context Protocol. 2024.
\bibitem{hu_adas_2025} Hu, Lu, Clune. ADAS. ICLR 2025. arXiv:2408.08435.
\bibitem{packer_memgpt_2023} Packer et al. MemGPT. arXiv:2310.08560, 2023.

\bibitem{gao2025selfevolving} Gao, Geng, Hua, Hu et al. A survey of self-evolving agents. arXiv:2507.21046, 2025.
\bibitem{rajaraman2026curriculum} Rajaraman, Huang, Dudik, Schapire, Foster, Krishnamurthy. Learning to reason with curriculum I: Provable benefits of autocurriculum. COLT 2026. arXiv:2603.18325.
\bibitem{graves2017automated} Graves, Bellemare, Menick, Munos, Kavukcuoglu. Automated curriculum learning for neural networks. ICML 2017.
\bibitem{wei2022cot} Wei et al. Chain-of-thought prompting. NeurIPS 35, 2022.
\bibitem{xi2023survey} Xi, Chen, Guo et al. The rise and potential of LLM based agents: A survey. arXiv:2309.07864, 2023.

\bibitem{yue2025rlvr} Yue, Chen, Lu, Zhao, Wang, Song, Huang. Does RL really incentivize reasoning capacity in LLMs beyond the base model? NeurIPS 2025. arXiv:2504.13837.
\bibitem{cui2025entropy} Cui et al. The entropy mechanism of RL for reasoning language models. arXiv:2505.22617, 2025.
\bibitem{dupoux2026why} Dupoux, LeCun, Malik. Why AI systems don't learn and what to do about it. arXiv:2603.15381, 2026.
\bibitem{wang2019poet} Wang, Lehman, Clune, Stanley. POET: Endlessly generating increasingly complex and diverse learning environments. arXiv:1901.01753, 2019.


% ---- Ajouts vérifiés le 03/09/2026 (agent de vérification, arXiv/dblp/JMLR) ----
\bibitem{shopr1_2026}
Yimeng Zhang, Tian Wang, Jiri Gesi, Ziyi Wang, Yuxuan Lu, Jiacheng Lin, Sinong Zhan, Vianne Gao, Ruochen Jiao, Junze Liu, Kun Qian, Yuxin Tang, Ran Xue, Houyu Zhang, Qingjun Cui, Yufan Guo, and Dakuo Wang.
\newblock Shop-R1: Rewarding LLMs to simulate human behavior in online shopping via reinforcement learning.
\newblock {\em International Conference on Learning Representations (ICLR)}, 2026. arXiv:2507.17842.

\bibitem{schulman2015trpo}
John Schulman, Sergey Levine, Philipp Moritz, Michael~I. Jordan, and Pieter Abbeel.
\newblock Trust region policy optimization.
\newblock {\em International Conference on Machine Learning (ICML)}, 2015. arXiv:1502.05477.

\bibitem{wei2022flan}
Jason Wei, Maarten Bosma, Vincent~Y. Zhao, Kelvin Guu, Adams~Wei Yu, Brian Lester, Nan Du, Andrew~M. Dai, and Quoc~V. Le.
\newblock Finetuned language models are zero-shot learners.
\newblock {\em International Conference on Learning Representations (ICLR)}, 2022. arXiv:2109.01652.

\bibitem{cs324}
Percy Liang, Tatsunori Hashimoto, and Christopher R\'e.
\newblock CS324: Large Language Models. Cours Stanford, hiver 2022.
\newblock \url{https://stanford-cs324.github.io/winter2022/}.

\bibitem{cs329a}
Aakanksha Chowdhery and Azalia Mirhoseini.
\newblock CS329A: Self-Improving AI Agents. Cours Stanford, automne 2025.
\newblock \url{https://cs329a.stanford.edu/}.

\bibitem{song2026swemaster}
Huatong Song, Lisheng Huang, Shuang Sun, et~al.
\newblock SWE-Master: Unleashing the potential of software engineering agents via post-training.
\newblock {\em arXiv preprint arXiv:2602.03411}, 2026.

\bibitem{liu2025prorl}
Mingjie Liu, Shizhe Diao, Ximing Lu, Jian Hu, Xin Dong, Yejin Choi, Jan Kautz, and Yi Dong.
\newblock ProRL: Prolonged reinforcement learning expands reasoning boundaries in large language models.
\newblock {\em Advances in Neural Information Processing Systems (NeurIPS)}, 2025. arXiv:2505.24864.

\bibitem{gorbatovski2024trdpo}
Alexey Gorbatovski, Boris Shaposhnikov, Alexey Malakhov, Nikita Surnachev, Yaroslav Aksenov, Ian Maksimov, Nikita Balagansky, and Daniil Gavrilov.
\newblock Learn your reference model for real good alignment.
\newblock {\em arXiv preprint arXiv:2404.09656}, 2024.

\bibitem{zhang2026emapg}
Zhang and Ba.
\newblock EMA policy gradient: Taming RL for LLMs with EMA anchor and top-$k$ KL.
\newblock {\em arXiv preprint arXiv:2602.04417}, 2026. % prénoms à compléter avant rendu

\bibitem{vonwerra2020trl}
Leandro von Werra, Younes Belkada, Lewis Tunstall, Edward Beeching, Tristan Thrush, Nathan Lambert, Shengyi Huang, Kashif Rasul, and Quentin Gallou\'edec.
\newblock TRL: Transformer Reinforcement Learning.
\newblock GitHub, 2020. \url{https://github.com/huggingface/trl}.

\bibitem{kwon2023vllm}
Woosuk Kwon, Zhuohan Li, Siyuan Zhuang, Ying Sheng, Lianmin Zheng, Cody~Hao Yu, Joseph~E. Gonzalez, Hao Zhang, and Ion Stoica.
\newblock Efficient memory management for large language model serving with PagedAttention.
\newblock {\em Symposium on Operating Systems Principles (SOSP)}, 2023. arXiv:2309.06180.

\bibitem{hu2022lora}
Edward~J. Hu, Yelong Shen, Phillip Wallis, Zeyuan Allen-Zhu, Yuanzhi Li, Shean Wang, Lu Wang, and Weizhu Chen.
\newblock LoRA: Low-rank adaptation of large language models.
\newblock {\em International Conference on Learning Representations (ICLR)}, 2022. arXiv:2106.09685.

\bibitem{narvekar2020curriculum}
Sanmit Narvekar, Bei Peng, Matteo Leonetti, Jivko Sinapov, Matthew~E. Taylor, and Peter Stone.
\newblock Curriculum learning for reinforcement learning domains: A framework and survey.
\newblock {\em Journal of Machine Learning Research}, 21(181):1--50, 2020. arXiv:2003.04960.

\bibitem{portelas2020acl}
R\'emy Portelas, C\'edric Colas, Lilian Weng, Katja Hofmann, and Pierre-Yves Oudeyer.
\newblock Automatic curriculum learning for deep RL: A short survey.
\newblock {\em International Joint Conference on Artificial Intelligence (IJCAI)}, 2020. arXiv:2003.04664.

\bibitem{chen2021codex}
Mark Chen, Jerry Tworek, Heewoo Jun, Qiming Yuan, Henrique Ponde de Oliveira Pinto, et~al.
\newblock Evaluating large language models trained on code.
\newblock {\em arXiv preprint arXiv:2107.03374}, 2021.




\end{thebibliography}









\appendix



\newpage
% =================================================================
% ANNEXE 1 — Fondements du Deep RL, du cadre POMDP à GRPO
% Condensé structuré du cours CS224R (Stanford, Deep Reinforcement Learning)
% Prérequis LaTeX : amsmath, amssymb
% =================================================================
\section{Annexe 1 : Fondements du Deep Reinforcement Learning, du cadre POMDP à GRPO}
\label{annexe:maths_rl}

Cette annexe présente de manière condensée le socle théorique de l'apprentissage par renforcement profond nécessaire à la lecture de ce document. Elle s'inspire grandement et suit la progression du cours CS224R de Stanford (\emph{Deep Reinforcement Learning}, \href{https://cs224r.stanford.edu/}{cs224r.stanford.edu}) et du cours de Berkley de Sergei Levine: formalisme de décision séquentielle, apprentissage par imitation, gradients de politique, méthodes Acteur-Critique et leurs versions \emph{off-policy} menant à PPO, apprentissage par valeurs, RL hors-ligne, apprentissage de la récompense, et enfin application aux LLMs jusqu'à GRPO et au raisonnement.

% =================================================================
\subsection*{A.1 Cadre Formel : MDP et POMDP}
% =================================================================
L'apprentissage par renforcement modélise un problème de décision séquentielle : un agent observe, agit, observe le résultat de son action, et itère. À chaque pas de temps $t$, l'agent se trouve dans un état $s_t$ (l'état du « monde »), choisit une action $a_t$, reçoit une récompense instantanée $r(s_t, a_t)$ quantifiant la qualité locale du couple état-action, et l'environnement transitionne vers un état $s_{t+1}$ selon une dynamique inconnue $p(s_{t+1} \mid s_t, a_t)$. Cette dynamique satisfait la propriété de Markov : l'état suivant ne dépend que de l'état et de l'action courants, indépendamment du passé. Le tuple $(\mathcal{S}, \mathcal{A}, p, r)$, muni d'une distribution initiale $p(s_1)$, définit un processus de décision markovien (\emph{Markov Decision Process}, MDP).

Le comportement de l'agent est représenté par une politique $\pi_\theta(a \mid s)$, une distribution de probabilité sur les actions conditionnée par l'état, paramétrée par un réseau de neurones. Une séquence $\tau = (s_1, a_1, s_2, a_2, \dots, s_T, a_T)$ est appelée trajectoire (ou \emph{rollout}, ou épisode), et sa probabilité sous la politique se factorise :
\begin{equation}
p_\theta(\tau) = p(s_1) \prod_{t=1}^{T} \pi_\theta(a_t \mid s_t)\, p(s_{t+1} \mid s_t, a_t)
\end{equation}
L'objectif du RL est de maximiser l'espérance de la somme des récompenses :
\begin{equation}
\max_\theta \; J(\theta) = \mathbb{E}_{\tau \sim p_\theta(\tau)} \left[ \sum_{t=1}^{T} r(s_t, a_t) \right]
\label{eq:rl_objective}
\end{equation}
L'espérance est indispensable car deux sources de stochasticité rendent le gain aléatoire : la dynamique de l'environnement et la politique elle-même. La stochasticité de la politique n'est pas un défaut mais une nécessité, tant pour l'exploration (essayer des actions différentes est la condition de l'apprentissage par essai-erreur) que pour la modélisation de comportements variés.

Lorsque l'agent n'a pas accès à l'état complet du monde mais seulement à une observation $o_t$ (une image caméra, l'historique d'une conversation), le cadre devient un processus de décision markovien partiellement observable (\emph{Partially Observable MDP}, POMDP). La conséquence formelle est déterminante : en marginalisant les états, les observations futures dépendent de tout l'historique, $p(o_{t+1} \mid o_{1:t}, a_{1:t}) \neq p(o_{t+1} \mid o_t, a_t)$. La politique doit alors être dotée de mémoire, $\pi_\theta(a_t \mid o_{1:t})$ — ce que réalise naturellement un LLM dont le contexte accumule la trace complète des interactions. C'est ce cadre POMDP qui décrit les agents LLM étudiés dans ce document.

Ce formalisme unifie des applications très diverses : en robotique, $s$ regroupe les images et positions articulaires et $a$ une commande motrice ; pour un agent conversationnel, $o$ est le dernier message de l'utilisateur, $a$ la réponse générée, et $\tau$ la trace de la conversation. Face à l'objectif (\ref{eq:rl_objective}), les familles d'algorithmes se distinguent par leur stratégie : imiter une politique performante (imitation), différencier directement l'objectif (gradients de politique), estimer la valeur de la politique courante pour l'améliorer (Acteur-Critique), estimer la valeur de la politique optimale (méthodes par valeurs), ou apprendre la dynamique elle-même (méthodes par modèle). Chacune fait des compromis différents sur le coût de collecte des données, la stabilité et les hypothèses requises.

% =================================================================
\subsection*{A.2 Apprentissage par Imitation}
% =================================================================
L'approche la plus directe consiste à imiter des démonstrations $\mathcal{D} = \{(s_1, a_1, \dots, s_T, a_T)\}$ collectées par un expert. Le clonage comportemental (\emph{Behavior Cloning}, BC) traite ce problème par apprentissage supervisé génératif, en maximisant la vraisemblance des actions de l'expert sous la politique :
\begin{equation}
\min_\theta \; -\mathbb{E}_{(s,a) \sim \mathcal{D}} \left[ \log \pi_\theta(a \mid s) \right]
\label{eq:bc}
\end{equation}
Cette formulation probabiliste (plutôt qu'une régression $\ell_2$) est essentielle dès que les données proviennent de plusieurs démonstrateurs : une régression vers la moyenne de comportements multimodaux produit une action qui n'appartient à aucun mode, tandis qu'un modèle génératif expressif (distribution catégorielle, mélange de gaussiennes, diffusion) capture la multimodalité.

Le BC est un algorithme \emph{offline} : il n'utilise que des données existantes, sans interaction. Sa limite structurelle est le décalage de covariables (\emph{covariate shift}) : contrairement au cadre supervisé classique où les entrées sont i.i.d., les actions prédites déterminent les états futurs. Une petite erreur emmène l'agent hors de la distribution des démonstrations, où ses erreurs s'amplifient — le phénomène d'erreurs composées (\emph{compounding errors}), $p_{\text{expert}}(s) \neq p_\pi(s)$. Les remèdes reposent sur la collecte de données correctives en ligne (DAgger et ses variantes), où un expert ré-annote ou corrige les états visités par la politique apprise. Par construction, l'imitation ne peut pas dépasser le démonstrateur et n'offre aucun cadre pour s'améliorer par la pratique.

Ce cadre est le pendant exact du SFT des LLMs : la perte (\ref{eq:bc}) est celle du SFT, les démonstrations étant les paires instruction-réponse. Les erreurs composées y ont aussi leur analogue, un modèle dérivant hors distribution au fil d'une génération longue ou d'un dialogue multi-tour.

% =================================================================
\subsection*{A.3 Gradients de Politique (REINFORCE)}
% =================================================================
Les gradients de politique attaquent directement l'objectif (\ref{eq:rl_objective}), avec une difficulté : la récompense n'est en général pas différentiable (elle peut être un vérificateur symbolique, un simulateur ou un jugement humain), et l'espérance porte sur une distribution qui dépend de $\theta$. L'astuce de la dérivée du logarithme, $\nabla_\theta p_\theta(\tau) = p_\theta(\tau) \nabla_\theta \log p_\theta(\tau)$, résout les deux problèmes à la fois. En notant $R(\tau) = \sum_t r(s_t, a_t)$ et en observant que les termes de dynamique de $\log p_\theta(\tau)$ ne dépendent pas de $\theta$ :
\begin{equation}
\nabla_\theta J(\theta) = \mathbb{E}_{\tau \sim p_\theta(\tau)} \left[ \sum_{t=1}^{T} \nabla_\theta \log \pi_\theta(a_t \mid s_t) \, R(\tau) \right]
\label{eq:reinforce}
\end{equation}
Le gradient est désormais une espérance sous $p_\theta$, estimable par échantillonnage Monte-Carlo de $N$ trajectoires. L'interprétation est éclairante : (\ref{eq:reinforce}) est exactement le gradient de l'imitation (\ref{eq:bc}) appliqué aux propres générations de la politique, pondéré par la récompense. Si $R > 0$, la vraisemblance des actions de la trajectoire est augmentée ; si $R < 0$, elle est diminuée. C'est la formalisation de l'apprentissage par essai-erreur : faire davantage ce qui a bien fonctionné.

L'estimateur brut souffre d'une variance élevée, que deux raffinements classiques réduisent sans introduire de biais. D'abord la causalité : l'action au temps $t$ ne peut influencer les récompenses passées, on remplace donc $R(\tau)$ par le retour restant (\emph{reward-to-go}) $\sum_{t' \geq t} r(s_{t'}, a_{t'})$. Ensuite la \emph{baseline} : soustraire une constante $b$ (typiquement la récompense moyenne) ne change pas le gradient en espérance mais recentre le signal, produisant des gradients négatifs pour les comportements sous la moyenne :
\begin{equation}
\nabla_\theta J(\theta) \approx \frac{1}{N} \sum_{i=1}^{N} \sum_{t=1}^{T} \nabla_\theta \log \pi_\theta(a_{i,t} \mid s_{i,t}) \left( \Big( \sum_{t' = t}^{T} r(s_{i,t'}, a_{i,t'}) \Big) - b \right)
\label{eq:pg_baseline}
\end{equation}
Même raffiné, cet estimateur reste bruité — un point critique en récompense \emph{sparse}, où la plupart des trajectoires ne portent aucun signal différencié. Il est de plus strictement \emph{on-policy} : le gradient exige des échantillons de la politique courante, donc chaque pas de gradient impose de recollecter des données, ce qui est prohibitif lorsque la génération est coûteuse.

% =================================================================
\subsection*{A.4 Fonctions de Valeur et Méthodes Acteur-Critique}
% =================================================================
Les méthodes Acteur-Critique améliorent le gradient de politique en apprenant à estimer ce qui est bon plutôt qu'en se contentant de l'observer. Trois objets sont centraux. La fonction de valeur et la Q-fonction d'une politique $\pi$ sont les espérances du retour :
\begin{align}
V^\pi(s) &= \mathbb{E}_\pi \Big[ \textstyle\sum_{t' \geq t} \gamma^{t'-t}\, r(s_{t'}, a_{t'}) \,\Big|\, s_t = s \Big] \\
Q^\pi(s, a) &= \mathbb{E}_\pi \Big[ \textstyle\sum_{t' \geq t} \gamma^{t'-t}\, r(s_{t'}, a_{t'}) \,\Big|\, s_t = s, \, a_t = a \Big]
\end{align}
où $\gamma \in (0, 1]$ est un facteur d'actualisation. $V^\pi(s)$ ne dépend que de l'état : c'est le gain moyen attendu en partant de $s$ et en suivant $\pi$. $Q^\pi(s,a)$ conditionne en plus sur la première action, potentiellement différente de celle que $\pi$ aurait choisie. Elles sont reliées par $V^\pi(s) = \mathbb{E}_{a \sim \pi(\cdot \mid s)}[Q^\pi(s,a)]$, ce qui définit l'avantage :
\begin{equation}
A^\pi(s, a) = Q^\pi(s, a) - V^\pi(s)
\end{equation}
soit le surcroît de gain de l'action $a$ par rapport au comportement moyen de $\pi$ en $s$. D'espérance nulle sous $\pi$, l'avantage est la \emph{baseline} idéale, dépendante de l'état : substitué au retour brut dans (\ref{eq:pg_baseline}), il produit l'estimateur de gradient le moins bruité de cette famille :
\begin{equation}
\nabla_\theta J(\theta) \approx \frac{1}{N} \sum_{i=1}^{N} \sum_{t=1}^{T} \nabla_\theta \log \pi_\theta(a_{i,t} \mid s_{i,t}) \, \hat{A}^\pi(s_{i,t}, a_{i,t})
\label{eq:actor_critic_grad}
\end{equation}
En pratique, seul un réseau Critique $V_\phi \approx V^\pi$ est appris, par régression supervisée. Deux régimes d'estimation s'opposent pour construire ses cibles. L'estimation Monte-Carlo régresse $V_\phi(s_t)$ sur les retours observés $\sum_{t' \geq t} r_{t'}$ : elle est sans biais mais à forte variance. Le \emph{bootstrapping} (apprentissage par différence temporelle, TD) utilise l'estimation courante du réseau lui-même comme cible, $y_t = r(s_t, a_t) + \gamma V_\phi(s_{t+1})$ : la variance est faible mais un biais est introduit puisque la cible est elle-même approximative. Les cibles à $n$ pas, $y_t = \sum_{t'=t}^{t+n-1} \gamma^{t'-t} r_{t'} + \gamma^n V_\phi(s_{t+n})$, interpolent entre les deux et dominent souvent en pratique. L'avantage s'estime alors par $\hat{A}_t = r(s_t, a_t) + \gamma V_\phi(s_{t+1}) - V_\phi(s_t)$ ou sa généralisation à $n$ pas.

% =================================================================
\subsection*{A.5 Acteur-Critique Off-Policy et PPO}
% =================================================================
L'algorithme Acteur-Critique de la section précédente reste \emph{on-policy}. Pour amortir le coût de collecte, on souhaite effectuer plusieurs pas de gradient sur un même lot de trajectoires — mais dès le premier pas, la politique $\pi_\theta$ diffère de la politique $\pi_{\theta_{\text{old}}}$ ayant généré les données. L'échantillonnage préférentiel (\emph{importance sampling}) corrige ce décalage : pour toute distribution de proposition $q$, $\mathbb{E}_{x \sim p}[f(x)] = \mathbb{E}_{x \sim q}[\frac{p(x)}{q(x)} f(x)]$. Appliqué au niveau des trajectoires, le ratio $\prod_t \pi_\theta(a_t \mid s_t)/\pi_{\theta_{\text{old}}}(a_t \mid s_t)$ explose ou s'annule avec l'horizon ; l'approximation standard, en passant à une espérance par pas de temps, ne conserve que le ratio local :
\begin{equation}
r_t(\theta) = \frac{\pi_\theta(a_t \mid s_t)}{\pi_{\theta_{\text{old}}}(a_t \mid s_t)}
\end{equation}
Un danger subsiste : les avantages ayant été estimés sous $\pi_{\theta_{\text{old}}}$, la politique est incitée à s'éloigner arbitrairement d'elle pour maximiser des avantages devenus obsolètes — surapprentissage caractéristique du régime \emph{off-policy}. Deux mécanismes le contiennent. Le premier est une contrainte explicite de région de confiance, $\mathbb{E}_s[D_{\text{KL}}(\pi_\theta(\cdot \mid s) \parallel \pi_{\theta_{\text{old}}}(\cdot \mid s))] \leq \varepsilon$. Le second, retenu par \emph{Proximal Policy Optimization} (PPO) \cite{schulman2017proximal}, borne le ratio par écrêtage (\emph{clipping}) : dès que $r_t(\theta)$ sort de $[1-\epsilon, 1+\epsilon]$ dans la direction favorable, le gradient s'annule et le token ou l'action concernés cessent de contribuer à la mise à jour. Le minimum avec l'objectif non clippé garantit que l'écrêtage n'opère que lorsqu'il rend l'objectif plus conservateur :
\begin{equation}
\mathcal{L}_{\text{PPO}}(\theta) = -\mathbb{E}_{(s_t, a_t) \sim \pi_{\theta_{\text{old}}}} \left[ \min\Big( r_t(\theta)\, \hat{A}_t, \; \text{clip}\big(r_t(\theta),\, 1-\epsilon,\, 1+\epsilon\big)\, \hat{A}_t \Big) \right]
\label{eq:ppo}
\end{equation}
En pratique ($\epsilon \approx 0{,}2$), chaque lot de trajectoires sert à plusieurs époques de mises à jour avant recollecte. PPO n'est que modérément \emph{off-policy} — il réutilise le lot courant, non l'ensemble de l'expérience passée — et c'est précisément ce compromis qui fait sa stabilité, au prix d'une moindre efficacité en données que les méthodes à tampon de rejeu (\emph{replay buffer}) comme SAC. Cette stabilité explique son statut d'algorithme de référence, en robotique comme pour les LLMs.

% =================================================================
\subsection*{A.6 Aparté : Méthodes par Valeurs (Q-Learning)}
% =================================================================
Une famille alternative se passe entièrement de politique explicite. Si l'on dispose d'une estimation exacte de $Q^{\pi^*}$ pour la politique optimale, la politique gloutonne $\pi(s) = \arg\max_a Q(s, a)$ est elle-même optimale. Le Q-learning exploite l'équation d'optimalité de Bellman :
\begin{equation}
Q^{\pi^*}(s, a) = r(s, a) + \gamma \, \mathbb{E}_{s' \sim p(\cdot \mid s, a)} \Big[ \max_{a'} Q^{\pi^*}(s', a') \Big]
\end{equation}
en régressant itérativement $Q_\phi$ vers les cibles bootstrappées $y = r + \gamma \max_{a'} Q_\phi(s', a')$, sur des transitions issues d'un tampon de rejeu — l'algorithme est nativement \emph{off-policy}, donc très efficace en données. Sa stabilisation requiert un ensemble d'artifices (réseaux cibles gelés, double Q-learning contre la surestimation systématique du $\max$, cibles à $n$ pas). Le $\max_{a'}$ le rend surtout adapté aux espaces d'actions discrets de petite taille : il est inapplicable en l'état à la génération de texte, où l'« action » est une séquence de tokens dans un espace combinatoire. Nous ne le développons pas davantage, mais ses concepts (équation de Bellman, bootstrapping, surestimation des valeurs) sont les fondements de la section suivante.

% =================================================================
\subsection*{A.7 RL Hors-Ligne (Offline RL)}
% =================================================================
Le RL hors-ligne pose la question suivante : peut-on apprendre une politique performante à partir d'un jeu de données statique $\mathcal{D}$, collecté par une politique de comportement inconnue $\pi_\beta$, sans aucune interaction nouvelle ? L'enjeu pratique est considérable (réutiliser les données existantes, éviter une collecte coûteuse ou risquée), et l'avantage théorique sur l'imitation est net : en exploitant les récompenses, le RL hors-ligne peut dépasser $\pi_\beta$ et recombiner des fragments de bonnes trajectoires (\emph{trajectory stitching}) qu'aucune trajectoire du jeu de données ne réalise en entier.

Le verrou central est la surestimation des valeurs hors distribution. Si l'on applique naïvement un algorithme \emph{off-policy}, l'amélioration de politique interroge $Q_\phi$ sur des actions $a'$ absentes du jeu de données, où le réseau est arbitrairement peu fiable ; la politique recherche activement ces zones de sur-optimisme, et l'erreur s'auto-amplifie à chaque itération — sans données nouvelles pour la corriger. Toutes les méthodes de RL hors-ligne organisent leur conservatisme autour de ce problème. La ligne la plus simple, l'imitation filtrée, ne conserve que les meilleures trajectoires du jeu de données et les clone. Plus fine, la régression pondérée par l'avantage (\emph{Advantage-Weighted Regression}, AWR) pondère l'imitation de chaque transition par la qualité estimée de l'action :
\begin{equation}
\theta \leftarrow \arg\max_\theta \; \mathbb{E}_{(s,a) \sim \mathcal{D}} \left[ \log \pi_\theta(a \mid s) \, \exp\big( \tfrac{1}{\alpha} \hat{A}(s, a) \big) \right]
\label{eq:awr}
\end{equation}
Cette forme approxime la solution d'un problème d'amélioration sous contrainte KL envers $\pi_\beta$, et surtout ne rétropropage que sur des actions présentes dans les données — le sur-optimisme hors distribution est évité par construction. Des méthodes comme IQL affinent l'estimation de valeur (régression par expectiles asymétriques) pour évaluer une politique meilleure que $\pi_\beta$ sans jamais interroger d'action hors support. Ce cadre éclaire directement certains algorithmes LLM : la perte (\ref{eq:awr}) est la forme générale dont les méthodes d'auto-amélioration par filtrage (imiter ses propres trajectoires réussies, telles STaR ou ReST) sont le cas binaire, et DPO s'analysera comme un algorithme hors-ligne sur données de préférences.

% =================================================================
\subsection*{A.8 Apprentissage de la Récompense}
% =================================================================
Tout ce qui précède suppose une fonction de récompense donnée. Or, hors des jeux et des tâches vérifiables, la récompense ne va pas de soi : elle doit être apprise, et tout signal appris peut être exploité. Un classifieur de succès pré-entraîné, utilisé comme récompense, est systématiquement piraté par l'optimisation RL, qui découvre les états hors de sa distribution d'entraînement où il se trompe avec confiance (\emph{reward hacking}) ; la parade générique est adversariale, en ré-entraînant continûment le signal sur les états que la politique visite.

Pour des juges humains, l'expérience montre que les notes absolues sont bruitées et mal calibrées, tandis que les comparaisons par paires sont fiables. Le modèle de Bradley-Terry postule que la probabilité de préférer $\tau_w$ à $\tau_l$ est $\sigma(r_\theta(\tau_w) - r_\theta(\tau_l))$, où $\sigma$ est la sigmoïde, et le modèle de récompense s'apprend par maximum de vraisemblance sur les préférences annotées :
\begin{equation}
\mathcal{J}_{\text{RM}}(\phi) = -\mathbb{E}_{(x, y_w, y_l) \sim \mathcal{D}} \left[ \log \sigma\big( \text{RM}_\phi(x, y_w) - \text{RM}_\phi(x, y_l) \big) \right]
\label{eq:bradley_terry}
\end{equation}
Le juge peut aussi être un autre LLM (\emph{RL from AI Feedback}, RLAIF), l'observation clé étant que la critique est plus facile que la génération. Dans tous les cas, le modèle de récompense demeure une approximation faillible des préférences réelles — son sur-ajustement par la politique (\emph{reward over-optimization}) est documenté et motive les garde-fous de la section suivante.

% =================================================================
\subsection*{A.9 RL pour les LLMs : RLHF et DPO}
% =================================================================
Le pipeline RLHF \cite{ouyang2022training} assemble les briques précédentes. Un LLM est une politique $\pi_\theta(y \mid x)$ générant une réponse $y = (y_1, \dots, y_T)$ token par token ; le MDP sous-jacent a pour état le contexte $(x, y_{<t})$, pour action le token suivant, et pour dynamique la simple concaténation — déterministe et connue. Après pré-entraînement et SFT (l'étape d'imitation), on collecte des comparaisons par paires sur des réponses échantillonnées, on entraîne un modèle de récompense via (\ref{eq:bradley_terry}), puis on optimise la politique par PPO contre ce modèle. La récompense étant apprise donc imparfaite, l'objectif est régularisé par une pénalité d'ancrage à la politique de référence $\pi_{\text{ref}}$ (le modèle SFT) :
\begin{equation}
\max_\theta \; \mathbb{E}_{y \sim \pi_\theta(\cdot \mid x)} \left[ \text{RM}_\phi(x, y) \right] - \beta \, D_{\text{KL}}\big( \pi_\theta(\cdot \mid x) \parallel \pi_{\text{ref}}(\cdot \mid x) \big)
\label{eq:rlhf_objective}
\end{equation}
Cette KL d'ancrage est distincte de la région de confiance de PPO (section A.5) : la première est un terme de l'objectif, dirigé vers une référence figée, préservant les capacités générales et limitant le piratage de la récompense ; la seconde est un mécanisme d'optimisation, dirigé vers la politique du lot précédent.

\emph{Direct Preference Optimization} (DPO) \cite{rafailov2024direct} court-circuite ce pipeline en observant que le problème (\ref{eq:rlhf_objective}) admet une solution en forme fermée, $\pi^*(y \mid x) \propto \pi_{\text{ref}}(y \mid x) \exp(\frac{1}{\beta}\text{RM}(x,y))$, que l'on peut inverser pour exprimer la récompense en fonction de la politique : $\text{RM}(x, y) = \beta \log \frac{\pi^*(y \mid x)}{\pi_{\text{ref}}(y \mid x)} + \beta \log Z(x)$. En injectant cette reparamétrisation dans la perte de Bradley-Terry (\ref{eq:bradley_terry}) — où la constante de normalisation $Z(x)$ s'annule par différence — on obtient une perte de classification binaire portant directement sur les paramètres du LLM :
\begin{equation}
\mathcal{L}_{\text{DPO}}(\theta; \pi_{\text{ref}}) = -\mathbb{E}_{(x, y_w, y_l)} \left[ \log \sigma \left( \beta \log \frac{\pi_\theta(y_w \mid x)}{\pi_{\text{ref}}(y_w \mid x)} - \beta \log \frac{\pi_\theta(y_l \mid x)}{\pi_{\text{ref}}(y_l \mid x)} \right) \right]
\end{equation}
DPO élimine le modèle de récompense explicite, le Critique et l'échantillonnage en ligne : c'est un algorithme hors-ligne au sens de la section A.7, simple et efficace, mais qui n'exploite pas de données fraîches de la politique courante.

% =================================================================
\subsection*{A.10 RL pour le Raisonnement : RLVR et GRPO}
% =================================================================
La fragilité des récompenses apprises trouve sa résolution la plus radicale dans les domaines à vérification déterministe : en mathématiques ou en programmation, un vérificateur symbolique délivre une récompense binaire impossible à pirater. Ce RL à récompense vérifiable (RLVR) est le moteur des modèles de raisonnement : la politique génère une longue chaîne de pensée avant sa réponse finale, et seule la justesse de cette dernière est récompensée — l'intégralité de la séquence de raisonnement recevant le signal terminal. Ce paradigme a ouvert un second axe de \emph{scaling}, celui du calcul à l'inférence : l'entraînement RLVR fait croître conjointement la précision et la longueur des chaînes de pensée, le modèle apprenant de lui-même à « réfléchir plus longtemps » sur les problèmes difficiles.

L'algorithme emblématique de ce régime est \emph{Group Relative Policy Optimization} (GRPO) \cite{shao2024deepseekmath}. Dans PPO appliqué aux LLMs, le Critique $V_\phi$ est un second réseau de la taille de l'Acteur, coûteux en mémoire, et son estimation de valeur est de qualité douteuse sous récompense terminale. GRPO le supprime en estimant la \emph{baseline} statistiquement : pour chaque prompt $q$, on échantillonne un groupe de $N$ réponses $\{y^{(1)}, \dots, y^{(N)}\}$ sous $\pi_{\theta_{\text{old}}}$, un vérificateur attribue les scores $\{R_1, \dots, R_N\}$, et l'avantage de chaque réponse est son score normalisé au sein du groupe :
\begin{equation}
\hat{A}_i = \frac{R_i - \text{mean}(\mathbf{R})}{\text{std}(\mathbf{R}) + \delta}
\end{equation}
Chaque token de la réponse $i$ hérite du même avantage scalaire $\hat{A}_i$, injecté dans l'objectif clippé de PPO, auquel s'ajoute la pénalité d'ancrage :
\begin{equation}
\mathcal{L}_{\text{GRPO}}(\theta) = -\frac{1}{N} \sum_{i=1}^{N} \frac{1}{|y^{(i)}|} \sum_{t=1}^{|y^{(i)}|} \min\Big( r_{i,t}(\theta)\, \hat{A}_i, \; \text{clip}\big(r_{i,t}(\theta), 1-\epsilon, 1+\epsilon\big)\, \hat{A}_i \Big) + \beta \, D_{\text{KL}}\big( \pi_\theta \parallel \pi_{\text{ref}} \big)
\label{eq:grpo}
\end{equation}
avec $r_{i,t}(\theta) = \pi_\theta(y_t^{(i)} \mid x, y_{<t}^{(i)}) / \pi_{\theta_{\text{old}}}(y_t^{(i)} \mid x, y_{<t}^{(i)})$. Deux propriétés méritent l'attention. D'une part, si toutes les réponses du groupe obtiennent le même score — typiquement $N$ échecs sur un problème trop difficile, ou $N$ succès sur un problème trop facile —, alors $\text{std}(\mathbf{R}) = 0$ et le gradient est nul : aucun apprentissage n'a lieu, ce qui fonde l'importance du calibrage de la difficulté (et motive les approches par curriculum étudiées dans ce document). D'autre part, la diffusion uniforme de $\hat{A}_i$ sur tous les tokens renonce à toute attribution de crédit fine au sein de la trajectoire : c'est le compromis qui rend l'algorithme frugal, et la limite que les modèles de récompense de processus (PRM) cherchent à lever.

En contexte agentique multi-tour, un ajustement supplémentaire s'impose : la séquence contient des tokens émis par l'environnement (observations, résultats d'outils) qui ne sont pas des actions de la politique. Un masque binaire $M_t \in \{0,1\}$ ($1$ sur les tokens générés par l'agent, $0$ sur ceux de l'environnement) restreint la perte aux seules actions, sous peine d'entraîner le modèle à prédire les réponses déterministes de l'environnement et de dégrader la politique — l'exact analogue trajectoriel du masquage de l'instruction en SFT.

Ce socle — gradient de politique, avantage relatif de groupe, ancrage KL, masquage des observations — constitue le cadre algorithmique de l'ensemble des expériences présentées dans ce document.










\newpage


% =================================================================
% ANNEXE B — Dérivation du SNIS (suppose \appendix actif)
% =================================================================
\section{Dérivation Mathématique du SNIS dans l'Optimisation de Politique de Groupe}
\label{sec:annexe_snis}

Cette annexe détaille le développement mathématique de l'échantillonnage préférentiel auto-normalisé (\emph{Self-Normalized Importance Sampling}, SNIS) appliqué à la réestimation du gradient dans les algorithmes de type GRPO. Le formalisme s'appuie sur la théorie des méthodes Monte Carlo \cite[chap.~9]{owen2013montecarlo} et découle d'échanges avec Otmane Sakhi (Criteo AI Lab, équipe Performance Science).

\subsection*{B.1 Cadre et Formulation Initiale}
Soit $x$ l'instruction initiale, $z$ la séquence de pensée et $a$ la réponse finale. On note $\pi_\theta(z, a \mid x) = \pi_\theta(z \mid x)\, \pi_\theta(a \mid x, z)$ la distribution jointe de génération. La fonction de valeur conditionnelle à maximiser s'écrit :
\begin{equation}
V(\pi_\theta \mid x) = \mathbb{E}_{(z,a) \sim \pi_\theta(\cdot \mid x)} \left[ r(x,a) \right]
\end{equation}
où $r(x,a)$ est la récompense vérifiable de la réponse $a$. Par l'astuce de la dérivée du logarithme (\emph{score function estimator}) :
\begin{equation}
\nabla_\theta V(\pi_\theta \mid x) = \mathbb{E}_{(z,a) \sim \pi_\theta(\cdot \mid x)} \left[ \nabla_\theta \log \pi_\theta(z,a \mid x) \, r(x,a) \right]
\label{eq:snis_gradient_base}
\end{equation}

\subsection*{B.2 Changement de Mesure et Recombinaison}
L'objectif est de réévaluer la pertinence de tout couple recomposé $(z_j, a_i)$, formé à partir de la pensée d'un rollout $j$ et de la réponse d'un rollout $i$, au sein d'un même groupe de $G$ générations issues du prompt $x$. En décomposant l'espérance de (\ref{eq:snis_gradient_base}) suivant les lois marginale et conditionnelle, puis en changeant la mesure sur $a$ :
\begin{align}
\nabla_\theta V(\pi_\theta \mid x) &= \mathbb{E}_{z \sim \pi_\theta(\cdot \mid x)} \left[ \mathbb{E}_{a \sim \pi_\theta(\cdot \mid x, z)} \left[ \nabla_\theta \log \pi_\theta(z,a \mid x) \, r(x,a) \right] \right] \\
&= \mathbb{E}_{z \sim \pi_\theta(\cdot \mid x)} \left[ \mathbb{E}_{a \sim \pi_\theta(\cdot \mid x)} \left[ \frac{\pi_\theta(a \mid x, z)}{\pi_\theta(a \mid x)} \nabla_\theta \log \pi_\theta(z,a \mid x) \, r(x,a) \right] \right]
\end{align}
Le ratio de densité $\pi_\theta(a \mid x, z) / \pi_\theta(a \mid x)$ permet d'exprimer $a$ sous la marginale $\pi_\theta(a \mid x)$ plutôt que sous la conditionnelle : les réponses peuvent alors provenir de n'importe quel rollout du groupe.

\subsection*{B.3 Estimateur Monte Carlo et Interversion des Sommes}
À chaque étape GRPO, le modèle génère $G$ complétions indépendantes $(z_i, a_i)_{i \in [1, G]}$. L'approximation Monte Carlo de la double espérance s'écrit :
\begin{equation}
\widehat{\nabla_\theta V}(\pi_\theta \mid x) = \frac{1}{G} \sum_{j=1}^G \frac{1}{G} \sum_{i=1}^G \frac{\pi_\theta(a_i \mid x, z_j)}{\pi_\theta(a_i \mid x)} \nabla_\theta \log \pi_\theta(z_j, a_i \mid x) \, r(x, a_i)
\end{equation}
Les sommes étant finies, on regroupe les termes par réponse $a_i$ :
\begin{equation}
\widehat{\nabla_\theta V}(\pi_\theta \mid x) = \frac{1}{G} \sum_{i=1}^G r(x, a_i) \cdot \frac{1}{G} \sum_{j=1}^G \frac{\pi_\theta(a_i \mid x, z_j)}{\pi_\theta(a_i \mid x)} \nabla_\theta \log \pi_\theta(z_j, a_i \mid x)
\end{equation}

\subsection*{B.4 Auto-Normalisation : des Marginales Inaccessibles aux Poids SNIS}
La marginale $\pi_\theta(a_i \mid x) = \mathbb{E}_{z \sim \pi_\theta(\cdot \mid x)}[\pi_\theta(a_i \mid x, z)]$ n'est pas accessible sous forme fermée, mais les $z_l$ du groupe étant précisément tirés de $\pi_\theta(\cdot \mid x)$, elle admet l'estimateur Monte Carlo $\hat{\pi}_\theta(a_i \mid x) = \frac{1}{G} \sum_{l=1}^G \pi_\theta(a_i \mid x, z_l)$. En substituant cet estimateur au dénominateur, les facteurs $1/G$ se compensent et l'estimateur de gradient prend la forme auto-normalisée :
\begin{equation}
\widehat{\nabla_\theta V}(\pi_\theta \mid x) = \frac{1}{G} \sum_{i=1}^G r(x, a_i) \sum_{j=1}^G \tilde{w}_{ij} \, \nabla_\theta \log \pi_\theta(z_j, a_i \mid x),
\qquad
\tilde{w}_{ij} = \frac{\pi_\theta(a_i \mid x, z_j)}{\sum_{l=1}^G \pi_\theta(a_i \mid x, z_l)}
\end{equation}

\paragraph{Interprétation en termes de cible et de proposition.} Ces poids réalisent un échantillonnage préférentiel auto-normalisé au sens strict : pour chaque réponse $a_i$ fixée, la distribution \emph{cible} est le postérieur $\pi_\theta(z \mid x, a_i) \propto \pi_\theta(z \mid x)\, \pi_\theta(a_i \mid x, z)$ et la distribution de \emph{proposition} est la marginale $\pi_\theta(z \mid x)$ dont sont issus les $z_l$. Les poids d'importance non normalisés sont alors proportionnels à la vraisemblance $\pi_\theta(a_i \mid x, z_j)$, et leur normalisation sur le groupe donne exactement $\tilde{w}_{ij}$.

\paragraph{Propriétés statistiques.} Comme tout estimateur auto-normalisé — ratio de deux estimateurs Monte Carlo corrélés —, $\widehat{\nabla_\theta V}$ est \emph{biaisé à $G$ fini} mais \emph{asymptotiquement consistant} (le biais décroît en $\mathcal{O}(1/G)$), en échange d'une réduction substantielle de variance par rapport à l'échantillonnage préférentiel classique, dont le dénominateur exact est ici de toute façon inaccessible \cite[chap.~9]{owen2013montecarlo}.

\paragraph{Tempérage.} Pour lisser ou accentuer l'attribution de crédit entre pensées du groupe, une température $\alpha > 0$ peut être appliquée \emph{avant} normalisation — l'appliquer après casserait la contrainte $\sum_j \tilde{w}_{ij} = 1$ :
\begin{equation}
\tilde{w}_{ij}^{(\alpha)} = \frac{\pi_\theta(a_i \mid x, z_j)^\alpha}{\sum_{l=1}^G \pi_\theta(a_i \mid x, z_l)^\alpha}
\end{equation}
On retrouve les poids SNIS pour $\alpha = 1$, une pondération uniforme quand $\alpha \to 0$, et une affectation dure à la pensée la plus vraisemblable quand $\alpha \to \infty$. Pour $\alpha \neq 1$, l'estimateur perd l'interprétation de la section précédente et introduit un biais supplémentaire ; le tempérage doit donc être vu comme une heuristique de régularisation, non comme un raffinement statistique.

\subsection*{B.5 Implications Computationnelles et Verrou Multi-tour}
\begin{itemize}
    \item \textbf{Asymétrie computationnelle.} La génération autorégressive de $G$ nouvelles trajectoires coûte $\mathcal{O}(G \cdot L)$ passes de décodage séquentiel ($L$ longueur de séquence), tandis que l'évaluation des $G^2$ vraisemblances $\pi_\theta(a_i \mid x, z_j)$ s'effectue en passes avant parallèles par simple \emph{teacher forcing} sur les paires recomposées. Le SNIS enrichit donc le signal de gradient — $G^2$ paires pondérées au lieu de $G$ trajectoires — à coût marginal faible.
    \item \textbf{Verrou en environnement multi-tour.} En raisonnement mono-tour, la pensée $z_j$ et la réponse $a_i$ ne dépendent que de $x$. En contrôle interactif, l'action au tour $t$ dépend de l'historique strict des observations $o_{1:t-1}$ renvoyées par le simulateur : recomposer $(z_j, a_i)$ à partir de rollouts distincts juxtapose des contextes d'exécution discordants, ce qui viole la cohérence causale de la trajectoire et invalide l'hypothèse de support commun entre cible et proposition. C'est ce verrou, confirmé empiriquement par l'effondrement de l'apprentissage sous recollement naïf (section 4), qui a maintenu cette piste au stade exploratoire ; la voie exécuteur-réviseur esquissée en section 2.6 — restreindre la recomposition à des sous-séquences causalement cohérentes extraites par un modèle juge — reste à formaliser, notamment l'expression des poids de transition entre tours.
\end{itemize}




\newpage
% =================================================================
% ANNEXE C — Analyse critique
% =================================================================
\section{Analyse Critique : Parcours, Reproduction et Dette Technique}
\label{sec:annexe_critique}

L'exploration menée durant ce stage s'est inscrite dans une double complexité : l'appropriation d'une littérature récente à la croisée du Deep RL et des LLMs, et le surcoût d'ingénierie inhérent à l'expérimentation en environnement décisionnel interactif. Cette annexe en retrace les choix, les impasses et les leçons.

\subsection*{C.1 Cadrage de Recherche et Réorientation Stratégique}
La première phase du stage a été consacrée à un travail d'intégration au sein des équipes de Criteo (problématiques ACT, AIAC et Performance Science), axé sur l'état de l'art en systèmes de recommandation (RecSys), les publicités conversationnelles et les identifiants sémantiques. Après un mois et demi d'exploration, un constat s'est imposé : la littérature actuelle en RecSys repose quasi exclusivement sur des formulations mono-tours statiques, sans utilisation d'outils externes. Établir un pont direct entre recommandation et agents LLM interactifs multi-tours s'est révélé prématuré pour un stage de six mois, en l'absence de fondations établies à ce croisement académique.

Cette observation a motivé une réorientation délibérée vers l'étude fondamentale des agents autonomes et du réentraînement par renforcement à récompense vérifiable. Plutôt que de disperser les efforts sur une intersection immature, ce choix a permis de construire un cadre théorique et méthodologique rigoureux autour d'AgentGym-RL et TextCraft. Le travail fournit ainsi une synthèse structurée d'une littérature morcelée et une brique technologique et conceptuelle directement réexploitable par Criteo pour ses futurs développements agentiques.

\subsection*{C.2 Reproduction de la Baseline : Coût et Enseignements}
Une part significative du temps expérimental a été consacrée à tenter de retrouver les performances rapportées par le papier AgentGym-RL sur TextCraft, afin d'asseoir nos comparaisons sur une base équitable. Cet effort s'est révélé beaucoup plus coûteux qu'anticipé, pour des raisons qui rejoignent l'analyse de la section 2.1 : en RL agentique, le résultat dépend autant du harnais — format de prompt, parseur d'actions, budget de tours, gestion des erreurs — que des poids et de l'algorithme, et ces choix sont rarement documentés au niveau de détail qu'exigerait une reproduction exacte. La poursuite du chiffre publié, au-delà de l'établissement d'une baseline saine, a un rendement scientifique décroissant : les semaines investies ont produit peu de connaissances généralisables sur nos questions de recherche.

L'exercice n'a cependant pas été stérile. Il a imposé une exploration systématique des leviers de stabilisation — régimes de taux d'apprentissage, restauration au meilleur état, adaptation LoRA, formats de prompt — dont plusieurs sont devenus des composantes de la configuration de référence (section 4.2), et il a forgé une compréhension fine de la sensibilité du RL multi-tour à des choix d'implémentation en apparence anodins. La leçon méthodologique que nous en tirons, et que nous appliquons dans la feuille de route : traiter la baseline publiée comme un ordre de grandeur, stabiliser notre propre configuration de référence, et concentrer l'effort de rigueur sur les comparaisons \emph{internes} à protocole constant — seules porteuses de conclusions fiables.

\subsection*{C.3 Obstacles Techniques et Leçons Algorithmiques}
Contrairement à l'apprentissage supervisé sur jeux de données statiques, le réentraînement d'un agent en environnement interactif impose de surmonter une dette d'ingénierie importante. Les difficultés rencontrées — matérielles, logicielles ou algorithmiques — se résument en cinq points :

\begin{itemize}
    \item \textbf{Réutilisation d'infrastructure.} Le développement d'environnements de simulation étant un travail d'ingénierie à part entière, le choix d'AgentGym-RL a permis de s'appuyer sur des serveurs HTTP existants et de concentrer l'effort de recherche sur l'optimisation de la politique.
    \item \textbf{Masquage des observations (\texttt{env\_mask}) sous TRL.} L'absence initiale de masquage des tokens renvoyés par l'environnement entraînait le recalcul erroné des log-probabilités, des divergences KL et des ratios d'échantillonnage préférentiel sur des tokens non générés par le modèle, corrompant silencieusement le gradient. Le diagnostic de ce bug — dont le symptôme n'était qu'une convergence anormalement lente puis un effondrement — a représenté l'un des débogages les plus coûteux du stage.
    \item \textbf{Synchronisation vLLM/TRL sous contraintes système.} L'articulation entre l'inférence rapide (vLLM, génération des rollouts) et l'optimisation (TRL, perte GRPO) s'est heurtée à des incompatibilités d'environnement (version ancienne de la \texttt{glibc}), imposant un mécanisme de synchronisation personnalisé entre moteur de génération et pipeline de rétropropagation, avant une migration vers le transfert direct des poids d'adaptateurs en mémoire GPU.
    \item \textbf{Reward shaping et reward hacking.} L'ajout de récompenses intermédiaires denses (bonus de $+0{,}02$ par action syntaxiquement valide) a induit un détournement de récompense caractérisé \cite{skalse2022defining} : la politique a appris à capituler dès le premier tour pour empocher le bonus de conformité, faisant chuter le taux de réussite de 18\,\% à 8\,\%. Dans notre configuration — TextCraft, vérificateur exact, shaping syntaxique —, le signal binaire outcome-only s'est révélé strictement préférable ; la généralisation de ce constat à d'autres schémas de densification reste ouverte (voir PRM, section 5).
    \item \textbf{Stabilisation de LoRA.} La recette LoRA publiée s'est effondrée dans notre régime multi-tour, par divergence de la KL vers une référence fixe. Le passage au rang 8 et surtout le déplacement périodique de la référence (section~4.2.4) ont stabilisé l'entraînement et porté le modèle à 73\,\%, puis à 82\,\% avec un curriculum d'horizon (section~4.3).

\end{itemize}




\section{Annexe D : transfert entre planification mono-tour et contrôle multi-tour}
\label{annexe:transfert}

\subsection*{D.1 Question et motivation}
Deux régimes d'apprentissage sont possibles pour un agent TextCraft. En contrôle multi-tour, le modèle choisit une action, observe le résultat, et recommence jusqu'à trente fois ; c'est le régime de tout le rapport. En planification mono-tour, il émet en une seule génération le plan complet, qui est ensuite exécuté pas à pas dans l'environnement sans retour vers le modèle. Le second régime coûte bien moins cher : une génération par tâche au lieu de trente allers-retours avec le serveur d'environnement, qui dominent le temps de chaque mise à jour. La question posée au début du stage était donc celle d'un méta-apprentissage entre les deux : un entraînement en planification, peu coûteux, accélère-t-il le contrôle interactif ? L'inverse est-il vrai ? Et un régime alterné réduirait-il le coût total sans perte ? La littérature invite à la prudence sur le premier terme : Valmeekam et al.~\cite{valmeekam2023planning} montrent que la planification autonome des LLM se dégrade vite avec la profondeur du problème, et Kambhampati et al.~\cite{kambhampati2024llmmodulo} soutiennent que leur fiabilité vient d'une boucle où un vérificateur externe les corrige. Dans notre cas, ce vérificateur est l'environnement lui-même.

\subsection*{D.2 Protocole des mesures}
La planification mono-tour est évaluée en deux temps. Le modèle reçoit les recettes disponibles et l'objectif, et rédige un plan en texte libre. Un second modèle, à température nulle, convertit ce plan en une liste d'actions, rejouée déterministiquement dans TextCraft ; la tâche est réussie si l'objet cible est dans l'inventaire final. Le score est donc celui du plan, indépendamment de toute réaction à l'environnement. Une variante « aveugle », où le convertisseur ne reçoit ni les recettes ni l'objectif, mesure la part du score due au convertisseur plutôt qu'au plan. Les scores ont été mesurés sur un balayage de onze températures de génération, de 0 à 1, à un tirage par température, pour lisser le bruit d'un tirage unique.

\subsection*{D.3 Résultats}
Trois observations, toutes sans entraînement dédié au transfert.

\begin{itemize}
    \item \textbf{La boucle de rétroaction vaut 23 points.} Qwen3.5-4B, non entraîné, atteint 54\,\% en planification mono-tour contre 77\,\% en contrôle interactif. Observer les résultats de ses actions rend la tâche plus facile que la planifier parfaitement à l'avance. La variante aveugle donne 46\,\% : les plans de ce modèle sont presque auto-suffisants. Pour Qwen2.5-3B, le score de planification passe de 20 à 12\,\% quand le convertisseur est aveugle : une part importante de son score vient de la réparation du plan par le convertisseur.
    \item \textbf{Le RL interactif transfère modestement vers la planification.} La politique Qwen2.5-3B entraînée par RL multi-tour (54\,\% en interactif) planifie mieux que le modèle de base : 19{,}5\,\% contre 12{,}3\,\% en moyenne sur les onze températures, soit sept points, sur onze comparaisons à même température toutes favorables. Le gain est réel mais faible au regard des 36 points gagnés en interactif : ce que le RL a appris est surtout réactif.
    \item \textbf{Le RL de planification détruit la compétence interactive.} Un run GRPO entraîné à produire directement le plan complet atteint 37\,\% en planification, mais retombe à 9\,\% en contrôle interactif, sous le modèle de base. Il ne respecte plus le protocole d'une action par tour. Les profondeurs 3 et 4 restent inchangées dans les deux régimes.
\end{itemize}

\subsection*{D.4 Portée}
Ces résultats étaient attendus : entraîner un modèle à tout planifier d'un coup lui fait perdre l'habitude d'attendre l'observation. Ils ne testent pas la question posée, celle d'un régime alterné ou d'un pré-alignement en planification suivi d'un RL interactif, qui aurait demandé des runs dédiés que le programme n'a pas laissé le temps de mener. Ils fixent en revanche deux repères : la valeur de la boucle de rétroaction, et l'asymétrie du transfert. L'ouverture (section~5) reprend la question sous l'angle du méta-apprentissage, avec la position de Dupoux, LeCun et Malik~\cite{dupoux2026why} pour qui les compétences d'agent s'acquièrent dans l'interaction plutôt que par prédiction hors boucle ; nos mesures vont dans ce sens sans le démontrer.

\section{Annexe E : l'autocurriculum MAGELLAN}
\label{annexe:magellan}

\subsection*{E.1 La méthode}
MAGELLAN~\cite{gaven2025magellan} traite le problème d'un agent autotélique qui doit choisir, dans un grand espace de buts, ceux qu'il va pratiquer. Le critère de choix est le progrès d'apprentissage~\cite{graves2017automated} : la variation de la compétence de l'agent sur un but entre deux instants. Un but déjà maîtrisé ne progresse plus ; un but hors de portée ne progresse pas non plus ; les buts informatifs sont ceux dont la compétence bouge. Estimer cette compétence but par but, en comptant les succès, devient impraticable quand les buts sont nombreux et rarement visités. La contribution de MAGELLAN est d'apprendre un estimateur de compétence : une petite tête de réseau, branchée sur la représentation interne que le modèle de langage forme de l'énoncé du but, prédit la probabilité de succès. Parce que des buts proches ont des représentations proches, l'estimateur généralise d'un but pratiqué aux buts voisins jamais pratiqués. Le progrès est mesuré comme l'écart absolu entre la prédiction courante et celle d'une copie retardée de l'estimateur, et les buts sont échantillonnés proportionnellement à ce progrès, avec une part d'échantillonnage uniforme qui décroît de 1 à 0{,}2 au fil de l'entraînement. Le progrès absolu capture aussi l'oubli : un but dont la compétence chute redevient prioritaire. Les auteurs valident la méthode avec un petit modèle entraîné par un algorithme à critique, sur un environnement textuel dont la majorité des buts est impossible, que le progrès d'apprentissage évacue de lui-même.

\subsection*{E.2 Notre port}
Nous avons transposé la couche de sélection des tâches, seule partie de la méthode indépendante de l'algorithme de RL, sur notre recette GRPO. L'estimateur est une tête à une couche cachée posée sur l'état caché du dernier token de l'énoncé de la tâche ; il est porté par des adaptateurs LoRA dédiés, distincts de ceux de la politique, comme le recommande l'ablation des auteurs, et entraîné en continu par entropie croisée sur les couples tâche-succès des épisodes récents. Le progrès d'une tâche est l'écart absolu entre sa prédiction courante et celle d'un instantané de l'estimateur vieux d'une centaine de mises à jour. Le progrès de toutes les tâches est recalculé toutes les 32 mises à jour, et l'échantillonnage des tâches est proportionnel au progrès, mélangé à une part uniforme décroissante. La recette d'entraînement est celle de la section~4.3 ; seule la distribution des tâches change.

\subsection*{E.3 Résultat et lecture}
Une prédiction a été écrite avant le lancement : la profondeur 4, dont le support de test est nul et qui ne compte qu'une recette d'entraînement, devait voir son progrès prédit tomber à zéro, et l'estimateur cesser de la présenter. Elle s'est vérifiée : en fin de run, la probabilité d'échantillonner la profondeur 4 était de six pour dix mille. Le run lui-même a atteint 45\,\% de pass@$1$ à la sixième époque, puis s'est effondré entre la septième et la huitième, selon la signature d'entropie décrite en section~4.2.5, plus tôt que le contrôle sans curriculum. La distribution d'échantillonnage finale éclaire ce résultat : 79\,\% de la masse sur la profondeur 1, 18\,\% sur la profondeur 2, 3\,\% sur la profondeur 3. Le progrès d'apprentissage mesure du changement, et un collapse est un changement considérable : quand la politique a commencé à se déterminiser, les prédictions de succès sur les tâches faciles se sont mises à osciller, l'estimateur y a vu du progrès, l'échantillonnage s'y est concentré, la diversité des tâches présentées a chuté, et le collapse s'est accéléré. C'est une rétroaction positive entre le signal de progrès et l'instabilité de l'apprenant. Elle ne met pas en cause la méthode, conçue pour un apprenant stable, mais elle en précise une condition d'emploi : l'autocurriculum suppose la stabilité, il ne la fournit pas. Un run à recette stabilisée par ailleurs, par exemple combiné à un curriculum d'horizon, serait le test suivant.





\end{document}